% Problems, Chapter 20: Integral Transforms
% Arfken, Weber & Harris, Mathematical Methods for Physicists, 7th ed.
% Requires: amsmath, amssymb, tikz

\section*{Problems}

\Needspace{9\baselineskip}
\subsection*{20.2 Fourier Transform}

\Needspace{5\baselineskip}
\par\bigskip
\noindent\textbf{20.2.1}
\begin{enumerate}
\item[(a)] Show that $g(-\omega)=g^{*}(\omega)$ is a necessary and
sufficient condition for $f(x)$ to be real.
\item[(b)] Show that $g(-\omega)=-g^{*}(\omega)$ is a necessary and
sufficient condition for $f(x)$ to be pure imaginary.
\end{enumerate}
\par\emph{Note.}~The condition of part (a) is used in the
development of the dispersion relations of Section~12.8.
\par\erratum{Use Fourier inversion in the usual $L^1$ or $L^2$ setting, and interpret reality almost everywhere; continuous representatives give pointwise statements.}
\solution
Use the convention
\[
g(\omega)=\frac1{\sqrt{2\pi}}\int_{-\infty}^{\infty}
f(x)e^{i\omega x}dx,\qquad
f(x)=\frac1{\sqrt{2\pi}}\int_{-\infty}^{\infty}
g(\omega)e^{-i\omega x}d\omega.
\]
Conjugating the first formula, with
$(e^{i\omega x})^*=e^{-i\omega x}$, shows
\[
g^*(\omega)=\frac1{\sqrt{2\pi}}\int f^*(x)e^{-i\omega x}dx
=\mathcal F[f^*](-\omega).
\]
(a) If $f=f^*$, the right side is $\mathcal F[f](-\omega)=g(-\omega)$,
so $g(-\omega)=g^*(\omega)$. Conversely, if
$g(-\omega)=g^*(\omega)$, conjugate the inverse transform and
put $u=-\omega$ (reversing the limits absorbs the sign of
$d\omega=-du$):
\begin{align*}
f^*(x)&=\frac1{\sqrt{2\pi}}\int g^*(\omega)e^{i\omega x}d\omega\\
&=\frac1{\sqrt{2\pi}}\int g(-\omega)e^{i\omega x}d\omega
=\frac1{\sqrt{2\pi}}\int g(u)e^{-iux}du=f(x).
\end{align*}
(b) Purely imaginary means $f^*=-f$. The conjugated forward
formula then gives
$g^*(\omega)=\mathcal F[-f](-\omega)=-g(-\omega)$, i.e.,
$g(-\omega)=-g^*(\omega)$. Conversely, if
$g^*(\omega)=-g(-\omega)$, the steps of the preceding inverse
calculation (conjugation, the hypothesis, then $u=-\omega$) give
$f^*(x)=-(2\pi)^{-1/2}\int g(-\omega)e^{i\omega x}d\omega=-f(x)$.
\erratum{For $L^1$ or $L^2$ transforms the recovered function is
determined almost everywhere; these reality assertions have that
meaning. They hold at every point for continuous representatives.
Changing a function at an isolated point cannot be detected by its
Fourier transform.}

\Needspace{5\baselineskip}
\par\bigskip
\noindent\textbf{20.2.2}\quad The function
\[
f(x)=\begin{cases}
1, & |x| < 1\\
0, & |x| > 1
\end{cases}
\]
is a symmetrical finite step function.
\begin{enumerate}
\item[(a)] Find $g_{c}(\omega)$, Fourier cosine transform of
$f(x)$.
\item[(b)] Taking the inverse cosine transform, show that
\[
f(x)=\frac{2}{\pi}\int_{0}^{\infty}
\frac{\sin\omega\cos\omega x}{\omega}\,d\omega.
\]
\item[(c)] From part (b) show that
\[
\int_{0}^{\infty}\frac{\sin\omega\cos\omega x}{\omega}\,d\omega
=\begin{cases}
0, & |x| > 1,\\[3pt]
\dfrac{\pi}{4}, & |x| = 1,\\[6pt]
\dfrac{\pi}{2}, & |x| < 1.
\end{cases}
\]
\end{enumerate}
\solution
(a) The unitary cosine transform is
\[
g_c(\omega)=\sqrt{\frac2\pi}\int_0^\infty f(x)\cos\omega x\,dx
=\sqrt{\frac2\pi}\left[\frac{\sin\omega x}{\omega}\right]_0^1
=\boxed{\sqrt{\frac2\pi}\frac{\sin\omega}{\omega}}.
\]
At zero its continuous value is $\sqrt{2/\pi}$.
(b) Substitution in the inverse cosine transform gives
\[
f(x)=\sqrt{\frac2\pi}\int_0^\infty g_c(\omega)\cos\omega x\,d\omega
=\frac2\pi\int_0^\infty\frac{\sin\omega\cos\omega x}{\omega}d\omega.
\]
At a jump, inversion gives the mean of the two one-sided limits,
so its value at $x=\pm1$ is $1/2$.

(c) By (b), the integral equals $(\pi/2)f(x)$: it is $\pi/2$ for
$|x|<1$ and $0$ for $|x|>1$, and the midpoint value $1/2$ at
$x=\pm1$ gives $(\pi/2)(1/2)=\pi/4$. To confirm these values,
including the edge value, evaluate the integral directly. For
$\epsilon>0$ first note that
\[
D_\epsilon(b)=\int_0^\infty e^{-\epsilon u}\frac{\sin bu}{u}du,
\qquad D_\epsilon'(b)=\int_0^\infty e^{-\epsilon u}\cos bu\,du
=\frac{\epsilon}{\epsilon^2+b^2}.
\]
Differentiation under the integral sign is allowed because the
$b$-derivative $e^{-\epsilon u}\cos bu$ of the integrand is bounded
by the integrable $e^{-\epsilon u}$; the last integral is $C$ of
Exercise~20.2.3 with $a=\epsilon$, $\omega=b$.
Since $D_\epsilon(0)=0$, integration in $b$ gives
\[
D_\epsilon(b)=\int_0^b\frac{\epsilon\,ds}{\epsilon^2+s^2}
=\arctan\frac b\epsilon.
\]
As $\epsilon\downarrow0$, $b/\epsilon\to\pm\infty$ for $b\gtrless0$.
Thus $D(b)=\lim_{\epsilon\downarrow0}D_\epsilon(b)
=(\pi/2)\operatorname{sgn}b$ for $b\ne0$, while
$D(0)=0$. The improper integral
$\int_0^\infty(\sin bu/u)\,du$ converges by Dirichlet's test, so
Abel's theorem for integrals makes it equal to this limit $D(b)$:
if $\int_0^\infty h\,du$ converges, then
$\int_0^\infty e^{-\epsilon u}h\,du\to\int_0^\infty h\,du$ as
$\epsilon\downarrow0$. To prove it, put $H(u)=\int_u^\infty h\,dv$;
then $H'=-h$, $H(0)=\int_0^\infty h\,du$, and $H$ is continuous with
$H(u)\to0$ as $u\to\infty$, hence bounded. Integration by parts gives
\begin{align*}
\int_0^\infty e^{-\epsilon u}h\,du
&=-\bigl[e^{-\epsilon u}H\bigr]_0^\infty
-\epsilon\int_0^\infty e^{-\epsilon u}H\,du\\
&=H(0)-\epsilon\int_0^\infty e^{-\epsilon u}H\,du.
\end{align*}
Splitting the last integral at any $U>0$ bounds the last term by
$\epsilon U\sup|H|+\sup_{u\ge U}|H(u)|$, which is arbitrarily small
when $U$ is chosen large and then $\epsilon$ small. Now use
\[
\sin\omega\cos\omega x
=\tfrac12\{\sin[(1+x)\omega]+\sin[(1-x)\omega]\}.
\]
Dividing by $\omega$ and integrating term by term,
\[
\int_0^\infty\frac{\sin\omega\cos\omega x}{\omega}d\omega
=\tfrac12[D(1+x)+D(1-x)]
=\frac\pi4[\operatorname{sgn}(1+x)+\operatorname{sgn}(1-x)],
\]
with $\operatorname{sgn}0=0$. For $|x|<1$ both signs are
positive; for $|x|>1$ they cancel; at $|x|=1$ just one is
positive. These are exactly the three stated values.

\Needspace{5\baselineskip}
\par\bigskip
\noindent\textbf{20.2.3}
\begin{enumerate}
\item[(a)] Show that the Fourier sine and cosine transforms of
$e^{-at}$ are
\[
g_{s}(\omega)=\sqrt{\frac{2}{\pi}}\,\frac{\omega}{\omega^{2}+a^{2}},
\qquad
g_{c}(\omega)=\sqrt{\frac{2}{\pi}}\,\frac{a}{\omega^{2}+a^{2}}.
\]
\par\emph{Hint.}~Each of the transforms can be related to the
other by integration by parts.
\item[(b)] Show that
\begin{align*}
\int_{0}^{\infty}\frac{\omega\sin\omega x}{\omega^{2}+a^{2}}\,d\omega
&=\frac{\pi}{2}\,e^{-ax},\qquad x>0,\\
\int_{0}^{\infty}\frac{\cos\omega x}{\omega^{2}+a^{2}}\,d\omega
&=\frac{\pi}{2a}\,e^{-ax},\qquad x>0.
\end{align*}
\end{enumerate}
These results can also be obtained by contour integration
(Exercise~11.8.12).
\par\erratum{Assume $a>0$.}
\solution
\erratum{Take $a>0$ for the ordinary exponentially decaying
transforms and the answers as printed.}

(a) Put
\[
C=\int_0^\infty e^{-at}\cos\omega t\,dt,\qquad
S=\int_0^\infty e^{-at}\sin\omega t\,dt.
\]
Integrate by parts using an antiderivative $-e^{-at}/a$ of the
exponential. The upper-limit terms vanish because $a>0$, so
\begin{align*}
C&=\left[-\frac{e^{-at}}a\cos\omega t\right]_0^\infty
-\frac\omega a\int_0^\infty e^{-at}\sin\omega t\,dt
=\frac1a-\frac\omega aS,\\
S&=\left[-\frac{e^{-at}}a\sin\omega t\right]_0^\infty
+\frac\omega a\int_0^\infty e^{-at}\cos\omega t\,dt
=\frac\omega aC.
\end{align*}
Substitute the second into the first and multiply by $a^2$:
$a^2C=a-\omega^2C$, hence
\[
C=\frac a{a^2+\omega^2},\qquad
S=\frac\omega{a^2+\omega^2}.
\]
Multiplication by $\sqrt{2/\pi}$ gives $g_c$ and $g_s$.
(b) For $x>0$, the inversion formulas
$f(x)=\sqrt{2/\pi}\int_0^\infty g_c(\omega)\cos\omega x\,d\omega$
and $f(x)=\sqrt{2/\pi}\int_0^\infty g_s(\omega)\sin\omega x\,d\omega$,
with $\sqrt{2/\pi}\sqrt{2/\pi}=2/\pi$, therefore read
\[
e^{-ax}=\frac{2a}{\pi}\int_0^\infty
\frac{\cos\omega x}{a^2+\omega^2}d\omega,
\qquad
e^{-ax}=\frac2\pi\int_0^\infty
\frac{\omega\sin\omega x}{a^2+\omega^2}d\omega.
\]
Multiplying the first by $\pi/(2a)$ and the second by $\pi/2$
gives the two integrals of (b).
The cosine integral is absolutely convergent. The sine integral
converges conditionally for $x>0$; it follows from inversion of
the odd extension of $e^{-at}$ on $t>0$. At $x=0$ that odd
extension has a jump and its inversion value is zero, so the
first formula in (b) must not be extended there by substitution.

\Needspace{5\baselineskip}
\par\bigskip
\noindent\textbf{20.2.4}\quad Find the Fourier transform of the
triangular pulse (Fig.~20.6),
\[
f(x)=\begin{cases}
h(1-a|x|), & |x| < 1/a,\\
0, & |x| > 1/a.
\end{cases}
\]
\par\emph{Note.}~This function provides another delta sequence
with $h=a$ and $a\to\infty$.
\par\erratum{Assume $a>0$, so the pulse has positive width.}
\solution
\erratum{Take $a>0$, so that the pulse has the stated positive
width.} Evenness eliminates the sine part:
\[
g(\omega)=\frac{2h}{\sqrt{2\pi}}
\int_0^{1/a}(1-ax)\cos\omega x\,dx.
\]
For $\omega\ne0$, use $\int\cos\omega x\,dx=\sin\omega x/\omega$
and, integrating by parts,
\[
\int x\cos\omega x\,dx
=\frac{x\sin\omega x}{\omega}+\frac{\cos\omega x}{\omega^2}.
\]
Between the limits $0$ and $1/a$ the terms proportional to
$\sin(\omega/a)$ cancel; then $1-\cos\theta=2\sin^2(\theta/2)$ gives
\begin{align*}
\int_0^{1/a}(1-ax)\cos\omega x\,dx
&=\frac{\sin(\omega/a)}\omega-a\left[\frac{\sin(\omega/a)}{a\omega}
+\frac{\cos(\omega/a)-1}{\omega^2}\right]\\
&=\frac{a[1-\cos(\omega/a)]}{\omega^2}
=\frac{2a\sin^2(\omega/2a)}{\omega^2}\\
&=\frac1{2a}\left[\frac{\sin(\omega/2a)}{\omega/2a}\right]^2.
\end{align*}
Multiplying by $2h/\sqrt{2\pi}$,
\[
\boxed{g(\omega)=\frac{h}{a\sqrt{2\pi}}
\left[\frac{\sin(\omega/2a)}{\omega/2a}\right]^2},
\qquad g(0)=\frac{h}{a\sqrt{2\pi}}.
\]
The area of the pulse, a triangle of base $2/a$ and height $h$,
is $h/a$, and $g(0)$ is $(2\pi)^{-1/2}$ times this area. With
$h=a$ the area is one; substitute $u=ax$, $dx=du/a$, in the
action of $f_a(x)=a(1-a|x|)$ on a continuous test function $\eta$:
\[
\int f_a(x)\eta(x)dx=\int_{-1}^1(1-|u|)\eta(u/a)du
\longrightarrow\eta(0)\int_{-1}^1(1-|u|)du=\eta(0),
\]
because $\eta(u/a)\to\eta(0)$ uniformly for $|u|\le1$ by the
continuity of $\eta$ at $0$.
This proves the delta-sequence claim, not merely a pointwise
concentration of the graph. Correspondingly, for $h=a$ the
transform $g(\omega)=(2\pi)^{-1/2}[\sin(\omega/2a)/(\omega/2a)]^2$
tends, for every $\omega$, to $(2\pi)^{-1/2}$, the transform of
$\delta(x)$.

\Needspace{5\baselineskip}
\par\bigskip
\noindent\textbf{20.2.5}\quad Consider the sequence
\[
\delta_{n}(x)=\begin{cases}
n, & |x| < 1/2n,\\
0, & |x| > 1/2n.
\end{cases}
\]
This is Eq.~(1.152). Express $\delta_{n}(x)$ as a Fourier integral
(via the Fourier integral theorem, inverse transform, etc.).
Finally, show that we may write
\[
\delta(x)=\lim_{n\to\infty}\delta_{n}(x)
=\frac{1}{2\pi}\int_{-\infty}^{\infty}e^{-ikx}\,dk.
\]
\par\erratum{The limit and the final Fourier integral are distributional identities, tested against smooth rapidly decreasing functions.}
\solution
The forward transform is
\begin{align*}
g_n(k)&=\frac n{\sqrt{2\pi}}\int_{-1/(2n)}^{1/(2n)}e^{ikx}dx
=\frac n{\sqrt{2\pi}}\,\frac{e^{ik/2n}-e^{-ik/2n}}{ik}\\
&=\frac n{\sqrt{2\pi}}\,\frac{2\sin(k/2n)}{k}
=\frac1{\sqrt{2\pi}}\frac{\sin(k/2n)}{k/2n}.
\end{align*}
Fourier inversion,
$\delta_n(x)=(2\pi)^{-1/2}\int g_n(k)e^{-ikx}dk$, consequently gives
\[
\boxed{\delta_n(x)=\frac1{2\pi}\int_{-\infty}^{\infty}
\frac{\sin(k/2n)}{k/2n}e^{-ikx}dk},
\]
with the midpoint value $n/2$ at either edge of the rectangle.
Those two point values do not affect any integral.
For a smooth rapidly decreasing test function $\eta$ (real or
complex), the interval $|x|<1/(2n)$ on which $\delta_n=n$ has
length $1/n$, so $\eta(0)=n\int_{-1/(2n)}^{1/(2n)}\eta(0)\,dx$ and
\begin{align*}
\Bigl|\int\delta_n(x)\eta(x)dx-\eta(0)\Bigr|
&=\Bigl|n\int_{-1/(2n)}^{1/(2n)}[\eta(x)-\eta(0)]\,dx\Bigr|\\
&\le\max_{|x|\le1/(2n)}|\eta(x)-\eta(0)|\longrightarrow0
\end{align*}
by the continuity of $\eta$ at $0$.
To take the same limit in the Fourier integral, let
$\tilde\eta(k)=(2\pi)^{-1/2}\int\eta(x)e^{ikx}dx$; it is
integrable, and $\eta(x)=(2\pi)^{-1/2}\int\tilde\eta(k)e^{-ikx}dk$.
Insert this in $\int\delta_n\eta\,dx$ and interchange the
absolutely convergent integrations. Since $g_n$ is even,
$\int\delta_n(x)e^{-ikx}dx=\sqrt{2\pi}\,g_n(-k)
=\sin(k/2n)/(k/2n)$, so
\[
\int\delta_n(x)\eta(x)dx
=\frac1{\sqrt{2\pi}}\int_{-\infty}^{\infty}
\frac{\sin(k/2n)}{k/2n}\,\tilde\eta(k)\,dk
\longrightarrow\frac1{\sqrt{2\pi}}\int_{-\infty}^{\infty}
\tilde\eta(k)\,dk=\eta(0).
\]
The limit follows from dominated convergence: the multiplier
tends to one, and $|\sin u/u|\le1$ for real $u$. The last step
is the inversion formula at $x=0$. Written out, the limit is
$(2\pi)^{-1}\int dk\int\eta(x)e^{ikx}dx
=(2\pi)^{-1}\int dk\int\eta(x)e^{-ikx}dx$ ($k\to-k$),
the action on $\eta$ of the kernel $(2\pi)^{-1}\int e^{-ikx}dk$.
Hence the limiting identity is
\[
\boxed{\delta(x)=\frac1{2\pi}\int_{-\infty}^{\infty}e^{-ikx}dk}.
\]
\erratum{This last integral is a distributional Fourier identity,
not a convergent ordinary integral or a pointwise limit of
functions at $x=0$. The test-function argument supplies its
precise meaning.}

\Needspace{5\baselineskip}
\par\bigskip
\noindent\textbf{20.2.6}\quad Using the sequence
\[
\delta_{n}(x)=\frac{n}{\sqrt{\pi}}\exp(-n^{2}x^{2}),
\]
show that
\[
\delta(x)=\frac{1}{2\pi}\int_{-\infty}^{\infty}e^{-ikx}\,dk.
\]
\par\emph{Hint.}~Remember that $\delta(x)$ is defined in terms of
its behavior as part of an integrand.
\par\erratum{Interpret the final Fourier integral as a distribution, rather than an ordinary convergent integral.}
\solution
The exponent completes to a square as
\[
-n^2x^2+ikx=-n^2\left(x-\frac{ik}{2n^2}\right)^2
-\frac{k^2}{4n^2}.
\]
Set $z=x-ik/(2n^2)$; as $x$ runs over $\mathbb R$, $z$ runs
along the horizontal line $\Im z=-k/(2n^2)$. This line can be
shifted back to the real axis: $e^{-n^2z^2}$ is entire, so by
Cauchy's theorem its integral around the rectangle with vertical
sides $z=\pm R+iv$ ($v$ between $0$ and $-k/(2n^2)$) vanishes, and
on those vertical sides the modulus is
$|e^{-n^2(\pm R+iv)^2}|=e^{-n^2R^2+n^2v^2}$, with $v$
in a fixed bounded interval, so those integrals tend to zero.
Consequently the remaining Gaussian integral is
$\int_{\mathbb R}e^{-n^2z^2}dz=\sqrt\pi/n$, and
\begin{align*}
g_n(k)&=\frac1{\sqrt{2\pi}}\frac n{\sqrt\pi}
\int_{-\infty}^{\infty}e^{-n^2x^2+ikx}dx
=\frac1{\sqrt{2\pi}}\frac n{\sqrt\pi}\,e^{-k^2/(4n^2)}
\frac{\sqrt\pi}n\\
&=\frac1{\sqrt{2\pi}}e^{-k^2/(4n^2)}.
\end{align*}
The inverse
$\delta_n(x)=(2\pi)^{-1/2}\int g_n(k)e^{-ikx}dk$ is therefore the
explicit regularized identity
\[
\frac n{\sqrt\pi}e^{-n^2x^2}
=\frac1{2\pi}\int_{-\infty}^{\infty}
e^{-k^2/(4n^2)}e^{-ikx}dk.
\]
For a bounded continuous test function $\eta$, put $u=nx$,
$dx=du/n$:
\[
\int\frac n{\sqrt\pi}e^{-n^2x^2}\eta(x)dx
=\frac1{\sqrt\pi}\int e^{-u^2}\eta(u/n)du
\longrightarrow\frac{\eta(0)}{\sqrt\pi}\int e^{-u^2}du=\eta(0).
\]
Dominated convergence applies, with the integrable majorant
$e^{-u^2}\sup|\eta|$. For a Schwartz test function $\eta$, with
transform $\tilde\eta$, the pairing argument of Exercise~20.2.5
gives likewise
\[
\int\frac n{\sqrt\pi}e^{-n^2x^2}\eta(x)dx
=\frac1{\sqrt{2\pi}}\int_{-\infty}^{\infty}
e^{-k^2/(4n^2)}\tilde\eta(k)\,dk
\longrightarrow\frac1{\sqrt{2\pi}}\int_{-\infty}^{\infty}
\tilde\eta(k)\,dk=\eta(0),
\]
because $0<e^{-k^2/(4n^2)}\le1$, the multiplier tends to one, and
$\tilde\eta$ is integrable. As there, the limit is the action of
the kernel $(2\pi)^{-1}\int e^{-ikx}dk$ on $\eta$. This gives the
required delta identity with the factor $1/(2\pi)$.
\erratum{The unregularized exponential integral is understood
in this weak, distributional sense; removing the Gaussian does
not make it an ordinary improper integral.}

\Needspace{5\baselineskip}
\par\bigskip
\noindent\textbf{20.2.7}\quad The formula
\[
\delta(t-x)=\frac{1}{2\pi}\int_{-\infty}^{\infty}
e^{i\omega(t-x)}\,d\omega
=\frac{1}{2\pi}\int_{-\infty}^{\infty}
e^{i\omega t}e^{-i\omega x}\,d\omega
\]
can be identified as the continuum limit of an eigenfunction
expansion. Derive sine and cosine representations of $\delta(t-x)$
that are comparable to the exponential representation just given.
\par\nobreak\hfill\emph{ANS.}\enspace
$\displaystyle\frac{2}{\pi}\int_{0}^{\infty}
\sin\omega t\sin\omega x\,d\omega$,\quad
$\displaystyle\frac{2}{\pi}\int_{0}^{\infty}
\cos\omega t\cos\omega x\,d\omega$.
\par\erratum{The two half-line answers above mean distributions for $x,t>0$. On the full line they require reflected delta terms, as derived in the solution.}
\solution
On a finite interval $(0,L)$ the orthonormal Dirichlet
eigenfunctions are $\sqrt{2/L}\sin(n\pi x/L)$, $n\ge1$, and the
orthonormal Neumann eigenfunctions are $1/\sqrt L$ and
$\sqrt{2/L}\cos(n\pi x/L)$, $n\ge1$; indeed
$\int_0^L\sin^2(n\pi x/L)dx=\int_0^L\cos^2(n\pi x/L)dx=L/2$.
With $\varphi_n$ denoting either orthonormal set, the completeness
relations $\sum_n\varphi_n(t)\varphi_n(x)=\delta(t-x)$ give, for
interior $x,t$,
\[
\delta(t-x)=\frac2L\sum_{n=1}^\infty
\sin\frac{n\pi t}L\sin\frac{n\pi x}L,
\]
\[
\delta(t-x)=\frac1L+\frac2L\sum_{n=1}^\infty
\cos\frac{n\pi t}L\cos\frac{n\pi x}L.
\]
Set $\omega_n=n\pi/L$, so $\Delta\omega=\omega_{n+1}-\omega_n=\pi/L$.
The factor $2/L=(2/\pi)\Delta\omega$ turns the sums into the
Riemann sums
$(2/\pi)\sum_n\sin\omega_nt\sin\omega_nx\,\Delta\omega$ and
$(2/\pi)\sum_n\cos\omega_nt\cos\omega_nx\,\Delta\omega$, which become
integrals over $0<\omega<\infty$ as $L\to\infty$; the constant
Neumann term $1/L$ vanishes. In the weak sense this gives the two
quoted kernels.

Their domains can be checked directly from the product-to-sum
formulas
$\sin\omega t\sin\omega x=\frac12[\cos\omega(t-x)-\cos\omega(t+x)]$
and
$\cos\omega t\cos\omega x=\frac12[\cos\omega(t-x)+\cos\omega(t+x)]$:
\begin{align*}
\frac2\pi\int_0^\infty\sin\omega t\sin\omega x\,d\omega
&=\delta(t-x)-\delta(t+x),\\
\frac2\pi\int_0^\infty\cos\omega t\cos\omega x\,d\omega
&=\delta(t-x)+\delta(t+x).
\end{align*}
Here $\delta(u)=\pi^{-1}\int_0^\infty\cos\omega u\,d\omega$, which is
the given exponential representation with $t-x$ replaced by $u$:
in $(2\pi)^{-1}\int_{-\infty}^\infty(\cos\omega u+i\sin\omega u)\,d\omega$
the odd sine term integrates to zero, and the even cosine term
gives twice its integral over $(0,\infty)$.
\erratum{The stated sine and cosine answers represent
$\delta(t-x)$ on the positive half-line, with $x,t>0$.
On the whole real line they contain the reflected source
$\mp\delta(t+x)$ shown above. Endpoint values at zero require
the corresponding half-line boundary convention.}
For fixed $x>0$ the reflected source $\delta(t+x)$ is concentrated
at $t=-x<0$, outside the support of any test function supported
inside $(0,\infty)$, so it contributes nothing there; this proves
exactly the needed positive-half-line identities.

\Needspace{5\baselineskip}
\par\bigskip
\noindent\textbf{20.2.8}\quad In a resonant cavity, an
electromagnetic oscillation of frequency $\omega_{0}$ dies out as
\[
A(t)=\begin{cases}
A_{0}e^{-\omega_{0}t/2Q}e^{-i\omega_{0}t}, & t>0,\\[3pt]
0, & t<0.
\end{cases}
\]
The parameter $Q$ is a measure of the ratio of stored energy to
energy loss per cycle. Calculate the frequency distribution of the
oscillation, $a^{*}(\omega)a(\omega)$, where $a(\omega)$ is the
Fourier transform of $A(t)$.
\par\emph{Note.}~The larger $Q$ is, the sharper your resonance
line will be.
\par\nobreak\hfill\emph{ANS.}\enspace
$a^{*}(\omega)a(\omega)=\dfrac{\erratum{|A_{0}|^{2}}}{2\pi}\,
\dfrac{1}{(\omega-\omega_{0})^{2}+(\omega_{0}/2Q)^{2}}$.
\par\erratum{Take $\omega_0>0$, $Q>0$. The amplitude $A_0$ may be complex, so its squared modulus is required in the answer.}
\solution
Put $\gamma=\omega_0/(2Q)>0$, the amplitude decay
rate. In the direct transformation of the causal signal the
kernel $e^{i\omega t}$ combines with $e^{-\gamma t-i\omega_0t}$
into a single exponential, and its value at the upper limit
vanishes because
$|e^{-[\gamma-i(\omega-\omega_0)]t}|=e^{-\gamma t}\to0$:
\begin{align*}
a(\omega)&=\frac{A_0}{\sqrt{2\pi}}\int_0^\infty
 e^{-[\gamma-i(\omega-\omega_0)]t}dt
=\frac{A_0}{\sqrt{2\pi}}\left[
-\frac{e^{-[\gamma-i(\omega-\omega_0)]t}}{\gamma-i(\omega-\omega_0)}
\right]_0^\infty\\
&=\frac{A_0}{\sqrt{2\pi}}
\frac1{\gamma-i(\omega-\omega_0)}.
\end{align*}
The conjugate $a^*(\omega)$ has the denominator
$\gamma+i(\omega-\omega_0)$, and the product of the two
denominators is $\gamma^2+(\omega-\omega_0)^2$. Hence
\[
\boxed{|a(\omega)|^2=\frac{|A_0|^2}{2\pi}
\frac1{(\omega-\omega_0)^2+(\omega_0/2Q)^2}.}
\]
\erratum{The printed numerator $A_0^2$ presumes a real initial
amplitude; for a complex amplitude it is $|A_0|^2$. Decay requires
$\omega_0/Q>0$, conventionally $\omega_0>0$, $Q>0$.}
The peak is at $\omega_0$, where the last factor equals
$1/\gamma^2$. Half its height, $1/(2\gamma^2)$, occurs where
$(\omega-\omega_0)^2=\gamma^2$, i.e., at
$|\omega-\omega_0|=\gamma$, so the full width at half maximum
is $2\gamma=\omega_0/Q$. Increasing $Q$ narrows the line.
As a normalization check (Parseval's theorem), use
$\int_{-\infty}^{\infty}d\omega/[(\omega-\omega_0)^2+\gamma^2]=\pi/\gamma$
and $\int_0^\infty e^{-2\gamma t}dt=1/(2\gamma)$:
\[
\int_{-\infty}^{\infty}|a(\omega)|^2d\omega
=\frac{|A_0|^2}{2\pi}\,\frac\pi\gamma=\frac{|A_0|^2}{2\gamma}
=|A_0|^2\int_0^\infty e^{-2\gamma t}dt
=\int_0^\infty|A(t)|^2dt.
\]

\Needspace{5\baselineskip}
\par\bigskip
\noindent\textbf{20.2.9}\quad Prove that
\[
\frac{\hbar}{2\pi i}\int_{-\infty}^{\infty}
\frac{e^{-i\omega t}\,d\omega}{E_{0}-i\Gamma/2-\hbar\omega}
=\begin{cases}
\exp\left(-\dfrac{\Gamma t}{2\hbar}\right)
\exp\left(-i\dfrac{E_{0}t}{\hbar}\right), & t>0,\\[12pt]
0, & t<0.
\end{cases}
\]
This Fourier integral appears in a variety of problems in quantum
mechanics: barrier penetration, scattering, time-dependent
perturbation theory, and so on.
\par\erratum{(The formula as given assumes that $\Gamma>0$.)}
\par\emph{Hint.}~Try contour integration.
\solution
Let
\[
\omega_p=\frac{E_0-i\Gamma/2}{\hbar},\qquad
\hbar>0,\quad E_0\in\mathbb R,\quad\Gamma>0.
\]
Since $E_0-i\Gamma/2-\hbar\omega=-\hbar(\omega-\omega_p)$, the
factor $\hbar$ cancels and the integral is
\[
I(t)=-\frac1{2\pi i}\int_{-\infty}^{\infty}
\frac{e^{-i\omega t}}{\omega-\omega_p}d\omega.
\]
The pole lies below the real axis: $\Im\omega_p=-\Gamma/(2\hbar)<0$.
If $t>0$, write $\omega=u+iv$; the modulus of the exponential is
$|e^{-iut}e^{vt}|=e^{vt}$, which decays for $v<0$,
so close the contour in the lower half-plane. The closure is
clockwise, and, because $1/(\omega-\omega_p)\to0$ uniformly on
large arcs, Jordan's lemma makes its arc integral vanish.
The residue of $e^{-i\omega t}/(\omega-\omega_p)$ at the simple
pole is $e^{-i\omega_pt}$; consequently
\[
I(t)=-\frac1{2\pi i}(-2\pi i)e^{-i\omega_pt}
=e^{-i(E_0-i\Gamma/2)t/\hbar}
=e^{-\Gamma t/(2\hbar)}e^{-iE_0t/\hbar}.
\]
If $t<0$, the factor $e^{vt}$ decays for $v>0$, so close above
instead. There is no pole inside, so the
integral is zero. This proves both cases, including the sign
from the contour orientation.
\erratum{The Instructor's Manual, p.~334, writes $e^{+i\omega_pt}$ in its
residue and final step. The exponent must be negative, as the
original integrand shows; the positive sign would produce
growth for $t>0$. The published condition $\Gamma>0$ is
necessary for the causal contour choice.}
At $t=0$ the symmetric principal-value inversion gives $1/2$,
the mean across the jump. Indeed, $\omega-\omega_p$ stays in the
upper half-plane, so along the real axis the principal
$\ln(\omega-\omega_p)$ is continuous and $\arg(\omega-\omega_p)$
decreases from $\pi$ to $0$, while
$\ln|R-\omega_p|-\ln|{-R}-\omega_p|\to0$ as $R\to\infty$; hence
\[
I(0)=-\frac1{2\pi i}\lim_{R\to\infty}
\bigl[\ln(\omega-\omega_p)\bigr]_{-R}^{R}
=-\frac{-i\pi}{2\pi i}=\frac12.
\]
The problem itself states only $t\ne0$.
For $\Gamma=0$, a pole prescription would
have to be supplied before the real-axis integral is defined.

\Needspace{5\baselineskip}
\par\bigskip
\noindent\textbf{20.2.10}\quad Verify that the following are
Fourier integral transforms of one another:
\begin{enumerate}
\item[(a)] $\left\{\begin{aligned}
&\sqrt{\frac{2}{\pi}}\cdot\frac{1}{\sqrt{a^{2}-x^{2}}},
&&|x| < a,\\[6pt]
&0, &&|x| > a,
\end{aligned}\right\}$ and $J_{0}(ay)$,
\item[(b)] $\left\{\begin{aligned}
&0, &&|x| < a,\\[6pt]
&-\sqrt{\frac{2}{\pi}}\,\frac{1}{\sqrt{\erratum{x^{2}-a^{2}}}},
&&|x| > a,
\end{aligned}\right\}$ and $Y_{0}(a|y|)$,
\item[(c)] $\sqrt{\dfrac{\pi}{2}}\,\dfrac{1}{\sqrt{x^{2}+a^{2}}}$
\quad and\quad $K_{0}(a|y|)$.
\item[(d)] Can you suggest why $I_{0}(ay)$ is not included in
this list?
\end{enumerate}
\par\erratum{(It is helpful to use the integral representation
$J_{0}(x)=\dfrac{2}{\pi}\displaystyle\int_{0}^{1}
\dfrac{\cos xt}{\sqrt{1-t^{2}}}\,dt$.)}
\par\emph{Hint.}~$J_{0}$, $Y_{0}$, and $K_{0}$ may be transformed
most easily by using an exponential representation, reversing the
order of integration, and employing the Dirac delta function
exponential representation, Eq.~(20.20). These cases can be
treated equally well as Fourier cosine transforms.
\par\erratum{Take $a>0$. In (b),(c), ordinary transform integrals are asserted for $y\ne0$; the identities also have a tempered-distribution interpretation.}
\solution
\erratum{Take $a>0$ for these transform pairs.} Since all the displayed functions are even,
the full unitary transform is
$\sqrt{2/\pi}\int_0^\infty f(x)\cos(xy)dx$: in
$(2\pi)^{-1/2}\int f(x)(\cos xy+i\sin xy)\,dx$ the odd sine part
integrates to zero, and the even cosine part gives twice its
integral over $(0,\infty)$.

(a) For the compactly supported member the prefactors give
$\sqrt{2/\pi}\sqrt{2/\pi}=2/\pi$, so
\begin{align*}
\mathcal F f(y)&=\frac2\pi\int_0^a
\frac{\cos(xy)}{\sqrt{a^2-x^2}}dx\\
&=\frac2\pi\int_0^1\frac{\cos(ayu)}{\sqrt{1-u^2}}du
=J_0(ay).
\end{align*}
The substitution $x=au$ gives $dx=a\,du$ and
$\sqrt{a^2-x^2}=a\sqrt{1-u^2}$, so both factors of $a$ cancel.
The last equality is the Poisson representation stated in the
published erratum.

(b) With the corrected square root, the prefactors give $-2/\pi$.
Since $\cos(xy)=\cos(x|y|)$, the substitution $x=au$
($dx=a\,du$, $\sqrt{x^2-a^2}=a\sqrt{u^2-1}$) gives
\[
\mathcal F f(y)=-\frac2\pi\int_a^\infty
\frac{\cos(xy)}{\sqrt{x^2-a^2}}dx
=-\frac2\pi\int_1^\infty
\frac{\cos(a|y|u)}{\sqrt{u^2-1}}du
=Y_0(a|y|),\quad y\ne0.
\]
The last equality is the real oscillatory representation
$Y_0(z)=-(2/\pi)\int_1^\infty\cos(zt)(t^2-1)^{-1/2}dt$, $z>0$,
derived in Exercise~14.4.8(b), with $z=a|y|$. Its endpoint
singularity is integrable and its infinite tail is conditionally
convergent for $y\ne0$.

(c) The prefactors now multiply to one,
$\sqrt{2/\pi}\sqrt{\pi/2}=1$, giving
\[
\mathcal F f(y)=\int_0^\infty
\frac{\cos(xy)}{\sqrt{x^2+a^2}}dx=K_0(a|y|),\quad y\ne0.
\]
To prove the last equality, write Exercise~14.5.13, with its $y$
replaced by $a$, in the unitary normalization:
\[
\frac1{\sqrt{x^2+a^2}}=\frac2\pi\int_0^\infty
K_0(aq)\cos(xq)dq
=\sqrt{\frac2\pi}\int_0^\infty
\Bigl[\sqrt{\frac2\pi}K_0(aq)\Bigr]\cos(xq)dq.
\]
Thus $1/\sqrt{x^2+a^2}$ is the cosine transform of
$\sqrt{2/\pi}K_0(aq)$, which is integrable on $(0,\infty)$ (its
singularity at $q=0$ is only logarithmic) and smooth for $q>0$.
The cosine transform is its own inverse, so for $q>0$
\[
\sqrt{\frac2\pi}\int_0^\infty\frac{\cos(xq)}{\sqrt{x^2+a^2}}dx
=\sqrt{\frac2\pi}K_0(aq).
\]
At $q=|y|$, with $\cos(xy)=\cos(x|y|)$, this is precisely the
asserted full-transform relation.
Applying the full transform twice reflects the original function,
$\mathcal F^2f(x)=f(-x)$: by the inversion formula,
$(2\pi)^{-1/2}\int g(y)e^{ixy}dy$ is the inverse transform of $g$
evaluated at $-x$, and for $g=\mathcal Ff$ it is $f(-x)$. All
three members here are even, so this also verifies the reverse
directions.
\erratum{The logarithmic singularities in (b) and (c) at
$y=0$ are not finite point values. Their transform identities
hold as tempered distributions and as the indicated ordinary
improper integrals away from zero; values at the support edges
in (a),(b) are likewise singular.}

(d) For $a>0$, the asymptotic form $I_0(z)\sim e^z/\sqrt{2\pi z}$ as
$z\to+\infty$ (Section~14.6), together with the evenness of $I_0$
(its series contains only even powers of $z$), gives
$I_0(ay)\sim e^{a|y|}/\sqrt{2\pi a|y|}$ as $|y|\to\infty$.
It therefore has neither an ordinary Fourier integral nor a
Fourier transform as a tempered distribution represented by this
function. This exponential growth is the obstruction to an
analogous pair in the present Fourier setting. At $a=0$ the
constant $I_0(0)=1$ instead has the distributional transform
$\sqrt{2\pi}\delta$, a separate degenerate case.

\Needspace{5\baselineskip}
\par\bigskip
\noindent\textbf{20.2.11}\quad \erratum{Show that the left-hand expression is the forward Fourier transform of the right-hand expression, and find the forward transform of the left-hand expression:}
\[
i^{n}J_{n}(t)\qquad\text{and}\qquad
\begin{cases}
\sqrt{\dfrac{2}{\pi}}\,T_{n}(x)(1-x^{2})^{-1/2},
& |x| < 1,\\[10pt]
0, & |x| > 1.
\end{cases}
\]
$T_{n}(x)$ is the $n$th-order Chebyshev polynomial.
\par\erratum{(The left-hand-side quantities are the transforms of
their right-hand-side counterparts, but the right-hand-side
quantities are $(-1)^{n}$ times the transforms of the
left-hand-side expressions.)}
\par\emph{Hint.}~With $T_{n}(\cos\theta)=\cos n\theta$, the
transform of $T_{n}(x)(1-x^{2})^{-1/2}$ leads to an integral
representation of $J_{n}(t)$.
\solution
Denote the compactly supported right-hand expression
by $f_n(x)$. Its transform, with prefactor
$(2\pi)^{-1/2}\sqrt{2/\pi}=1/\pi$, is
\[
\mathcal Ff_n(t)=\frac1\pi\int_{-1}^1
\frac{T_n(x)e^{itx}}{\sqrt{1-x^2}}dx.
\]
Set $x=\cos\theta$, where $0\le\theta\le\pi$; then
$\sqrt{1-x^2}=\sin\theta$ and $T_n(x)=\cos n\theta$. Reversing
the limits removes the sign of $dx=-\sin\theta\,d\theta$, and
$\sin\theta$ cancels:
\[
\mathcal Ff_n(t)=\frac1\pi\int_0^\pi
 e^{it\cos\theta}\cos(n\theta)d\theta.
\]
For integer $n\ge0$ the integral is half the full-circle
integral, because reflection $\theta\mapsto2\pi-\theta$ leaves
both factors unchanged. The sine part of $e^{in\theta}$ integrates
to zero under the same reflection, because $\sin n\theta$ changes
sign. To evaluate the full-circle integral, use the
Jacobi--Anger expansion of Exercise~14.1.5,
\[
e^{it\cos\theta}=\sum_{m=-\infty}^{\infty}i^mJ_m(t)e^{im\theta}.
\]
Multiply by $e^{in\theta}$ and integrate term by term; as
$(2\pi)^{-1}\int_0^{2\pi}e^{i(m+n)\theta}d\theta=\delta_{m,-n}$,
only $m=-n$ survives. With $J_{-n}=(-1)^nJ_n$ and
$i^{-n}(-1)^n=i^n$,
\[
\frac1\pi\int_0^\pi e^{it\cos\theta}\cos n\theta\,d\theta
=\frac1{2\pi}\int_0^{2\pi}e^{it\cos\theta+in\theta}d\theta
=i^{-n}J_{-n}(t)=\boxed{i^nJ_n(t)}.
\]
To check the direction of the transform, apply it again and use
$\mathcal F^2f(x)=f(-x)$ (Exercise~20.2.10):
\[
\mathcal F[i^nJ_n](x)=f_n(-x)=(-1)^nf_n(x),
\]
since $T_n(-x)=(-1)^nT_n(x)$. Thus the inverse transform of
$i^nJ_n$ is $f_n$, while its forward transform is $(-1)^nf_n$.
This proves, and explains, the published sign correction. The
square-root singularities are integrable; reverse inversion can
be read in the tempered-distribution sense at the singular edges.

\Needspace{5\baselineskip}
\par\bigskip
\noindent\textbf{20.2.12}\quad Show that the Fourier exponential
transform of
\[
f(\mu)=\begin{cases}
P_{n}(\mu), & |\mu| \le 1\\
0, & |\mu| > 1
\end{cases}
\]
is \erratum{$i^{n}\sqrt{2/\pi}\,j_{n}(\omega)$, where $\omega$ is
the transform variable (the book prints
$\left(2i^{n}/2\pi\right)j_{n}(kr)$)}. Here $P_{n}(\mu)$ is a
Legendre polynomial and $j_{n}$ is a spherical Bessel function.
\solution
The forward transform, with variable $\omega$, is
\[
g(\omega)=\frac1{\sqrt{2\pi}}\int_{-1}^1
P_n(\mu)e^{i\omega\mu}d\mu.
\]
Use Rodrigues' formula
$P_n(\mu)=(2^nn!)^{-1}d^n(\mu^2-1)^n/d\mu^n$.
Each integration by parts moves one derivative from
$(\mu^2-1)^n$ onto $e^{i\omega\mu}$, contributing a factor $-1$
and, from $de^{i\omega\mu}/d\mu=i\omega e^{i\omega\mu}$, a factor
$i\omega$. There are no endpoint terms: the polynomial
$(\mu^2-1)^n=(\mu-1)^n(\mu+1)^n$ has zeros of order $n$ at
$\mu=\pm1$, so it and its first $n-1$ derivatives vanish
at both endpoints. After $n$ steps the factor is
$(-1)^n(i\omega)^n=(-i\omega)^n$; in the second line below,
$(\mu^2-1)^n=(-1)^n(1-\mu^2)^n$ and $(-1)^n(-i\omega)^n=i^n\omega^n$.
Hence
\begin{align*}
g(\omega)
&=\frac{(-i\omega)^n}{2^nn!\sqrt{2\pi}}
\int_{-1}^1(\mu^2-1)^ne^{i\omega\mu}d\mu\\
&=\frac{i^n\omega^n}{2^nn!\sqrt{2\pi}}
\int_{-1}^1(1-\mu^2)^ne^{i\omega\mu}d\mu.
\end{align*}
The spherical Poisson representation of Exercise~14.7.4 is
\[
j_n(\omega)=\frac{\omega^n}{2^{n+1}n!}
\int_{-1}^1(1-\mu^2)^ne^{i\omega\mu}d\mu.
\]
Multiplied by $2i^n/\sqrt{2\pi}$, its right side becomes the
second line for $g(\omega)$, since $2/2^{n+1}=1/2^n$. Hence
\[
\boxed{g(\omega)=\frac{2i^n}{\sqrt{2\pi}}j_n(\omega)
=i^n\sqrt{\frac2\pi}\,j_n(\omega)}.
\]
No division by $\omega$ occurs, so the result holds at $\omega=0$
as well: there both sides equal $\sqrt{2/\pi}\,\delta_{n0}$, since
$\int_{-1}^1P_n\,d\mu=2\delta_{n0}$ and $j_n(0)=\delta_{n0}$.
This derivation verifies
both the scaling and the transform variable in the marked
correction; $kr$ is only a possible later choice of $\omega$.

\Needspace{5\baselineskip}
\par\bigskip
\noindent\textbf{20.2.13}
\begin{enumerate}
\item[(a)] Show that $f(x)=x^{-1/2}$ is a \emph{self-reciprocal}
under both Fourier cosine and sine transforms; that is,
\begin{align*}
\sqrt{\frac{2}{\pi}}\int_{0}^{\infty}
x^{-1/2}\cos xt\,dx&=t^{-1/2},\\
\sqrt{\frac{2}{\pi}}\int_{0}^{\infty}
x^{-1/2}\sin xt\,d\erratum{x}&=t^{-1/2}.
\end{align*}
\item[(b)] Use the preceding results to evaluate the Fresnel
integrals
\[
\int_{0}^{\infty}\cos\big(y^{2}\big)\,dy
\quad\text{and}\quad
\int_{0}^{\infty}\sin\big(y^{2}\big)\,dy.
\]
\end{enumerate}
\par\erratum{In part (a), take $t>0$.}
\solution
\erratum{The differential in the second integral must be
$dx$: the textbook, p.~979, prints $ds$.
Take $t>0$; both displayed integrals then converge conditionally.}
(a) Introduce positive damping before taking the square-root limit:
\[
I_\epsilon(t)=\int_0^\infty x^{-1/2}e^{-(\epsilon-it)x}dx
=\Gamma(1/2)(\epsilon-it)^{-1/2}.
\]
The gamma integral applies because $\Re(\epsilon-it)>0$: for real
$s>0$ the substitution $u=sx$ gives
$\int_0^\infty x^{-1/2}e^{-sx}dx=s^{-1/2}\int_0^\infty u^{-1/2}e^{-u}du
=\Gamma(1/2)s^{-1/2}$, with $\Gamma(1/2)=\sqrt\pi$, and both sides
are analytic in $s$ for $\Re s>0$, so the identity persists there
with the principal branch of $s^{-1/2}$.
As $\epsilon\downarrow0$, $|\epsilon-it|\to t$ and, on the principal
branch, $\arg(\epsilon-it)\to-\pi/2$, so
$(\epsilon-it)^{-1/2}\to t^{-1/2}e^{i\pi/4}$ and, with
$e^{i\pi/4}=(1+i)/\sqrt2$,
\[
I(t)=\lim_{\epsilon\downarrow0}I_\epsilon(t)
=\sqrt\pi\,t^{-1/2}e^{i\pi/4}
=\sqrt{\frac\pi2}\,t^{-1/2}(1+i).
\]
The ordinary integrals converge at zero by the exponent $-1/2$
and at infinity by Dirichlet's test, so Abel's theorem for integrals
(Exercise~20.2.2(c), here with $h(x)=x^{-1/2}e^{ixt}$) identifies
this limit with their values:
$I(t)=\int_0^\infty x^{-1/2}\cos xt\,dx+i\int_0^\infty x^{-1/2}\sin xt\,dx$.
Equating real and imaginary parts, each of these integrals is
$\sqrt{\pi/2}\,t^{-1/2}$; multiplying by $\sqrt{2/\pi}$ proves both
self-reciprocal identities.
\erratum{The Instructor's Manual, p.~336, has two further slips: the change
$u=-izt$ (its integration variable $z$ is our $x$) maps the
positive real $z$ axis to the negative imaginary $u$ axis, and
its final factor $t^{1/2}$ must be $t^{-1/2}$.
The damped gamma calculation above fixes both the phase and
the scaling without an ambiguous contour rotation.}

(b) For the Fresnel integrals take $t=1$ in
$\int_0^\infty x^{-1/2}\cos xt\,dx=\int_0^\infty x^{-1/2}\sin xt\,dx
=\sqrt{\pi/2}\,t^{-1/2}$ and set $x=y^2$. Then $dx=2y\,dy$ and
$x^{-1/2}dx=2dy$, with the limits $0,\infty$ unchanged, giving
\[
2\int_0^\infty\cos(y^2)dy
=2\int_0^\infty\sin(y^2)dy=\sqrt{\frac\pi2}.
\]
Thus each is $\boxed{\sqrt{\pi/8}}$. The substitution is first
made at a finite upper limit, followed by the convergent improper
limit; it does not require absolute convergence.

\Needspace{5\baselineskip}
\par\bigskip
\noindent\textbf{20.2.14}\quad Show that
$\left[\dfrac{1}{r^{2}}\right]^{T}\!\!(\mathbf{k})
=\left(\dfrac{\pi}{2}\right)^{1/2}\dfrac{1}{k}$.
\par\erratum{Here $k=|\mathbf k|>0$. The equality extends distributionally, with a singular zero-frequency value.}
\solution
Let $k=|\mathbf k|>0$ and use a polar axis parallel
to $\mathbf k$. To make the large-radius integration explicit,
first multiply $r^{-2}$ by $e^{-\epsilon r}$, $\epsilon>0$.
Since $d^3r=r^2dr\,d\Omega$, the radial singularity cancels
against the volume factor. With
$d\Omega=\sin\theta\,d\theta\,d\phi$ and $u=\cos\theta$, the
angular integration is
\[
\int d\Omega\,e^{ikr\cos\theta}
=2\pi\int_{-1}^1e^{ikru}du
=2\pi\frac{e^{ikr}-e^{-ikr}}{ikr}=4\pi\frac{\sin kr}{kr}.
\]
Thus
\begin{align*}
\mathcal F\left[\frac{e^{-\epsilon r}}{r^2}\right](\mathbf k)
&=\frac{4\pi}{(2\pi)^{3/2}k}
\int_0^\infty e^{-\epsilon r}\frac{\sin kr}{r}dr\\
&=\sqrt{\frac2\pi}\frac1k\arctan\frac{k}{\epsilon}.
\end{align*}
Here $4\pi/(2\pi)^{3/2}=\sqrt{2/\pi}$, and the last integral is
$D_\epsilon(k)=\arctan(k/\epsilon)$ of Exercise~20.2.2 (with $b=k$,
$u=r$). As $\epsilon\downarrow0$, $\arctan(k/\epsilon)\to\pi/2$ and
$\sqrt{2/\pi}\,\pi/2=\sqrt{\pi/2}$. The limit also holds after
pairing with a Schwartz test function $\eta(\mathbf k)$. Put
$\tilde\eta(\mathbf r)=(2\pi)^{-3/2}\int\eta(\mathbf k)
e^{i\mathbf k\cdot\mathbf r}d^3k$, again a Schwartz function. Since
$\eta$ and $e^{-\epsilon r}r^{-2}$ are integrable
($\int e^{-\epsilon r}r^{-2}d^3r=4\pi\int_0^\infty e^{-\epsilon r}dr
=4\pi/\epsilon$), the order of integration may be interchanged:
\[
\int\mathcal F\Bigl[\frac{e^{-\epsilon r}}{r^2}\Bigr](\mathbf k)\,
\eta(\mathbf k)\,d^3k
=\int\frac{e^{-\epsilon r}}{r^2}\,\tilde\eta(\mathbf r)\,d^3r.
\]
As $\epsilon\downarrow0$, dominated convergence takes the left side
to $\int\sqrt{\pi/2}\,k^{-1}\eta\,d^3k$ (majorant
$\sqrt{\pi/2}\,k^{-1}|\eta|$, since $0<\arctan(k/\epsilon)<\pi/2$)
and the right side to $\int r^{-2}\tilde\eta\,d^3r$ (majorant
$r^{-2}|\tilde\eta|$). Both majorants are integrable, because
$k^{-1}$ and $r^{-2}$ are locally integrable in three dimensions
($d^3k=k^2dk\,d\Omega$, $d^3r=r^2dr\,d\Omega$) and $\eta$,
$\tilde\eta$ decrease rapidly. The transform of the distribution
$r^{-2}$ is defined by this pairing rule,
$\int\mathcal F[r^{-2}]\,\eta\,d^3k=\int r^{-2}\tilde\eta\,d^3r$, so
the equality of the two limits for every $\eta$ states that
\[
\boxed{\mathcal F[r^{-2}](\mathbf k)
=\sqrt{\frac\pi2}\frac1k}.
\]
\erratum{Here $k$ is the nonnegative magnitude of the wave
vector. The formula is for $k>0$, or is a tempered-distribution
identity on all of three-dimensional momentum space. Its value
at $k=0$ is singular, and the original spatial integral is not
absolutely convergent at infinity.}

\Needspace{5\baselineskip}
\par\bigskip
\noindent\textbf{20.2.15}\quad The Fourier transform formulas for a
function of two variables are
\begin{align*}
F(u,v)&=\frac{1}{2\pi}\iint
f(x,y)e^{i(ux+vy)}\,dx\,dy,\\
f(x,y)&=\frac{1}{2\pi}\iint
F(u,v)e^{-i(ux+vy)}\,du\,dv,
\end{align*}
where the integrations are over the entire $xy$ or $uv$ plane. For
$f(x,y)=f\big([x^{2}+y^{2}]^{1/2}\big)=f(r)$, show that the
zero-order Hankel transforms
\begin{align*}
F(\rho)&=\int_{0}^{\infty}r f(r)J_{0}(\rho r)\,dr,\\
f(r)&=\int_{0}^{\infty}\rho F(\rho)J_{0}(\rho r)\,d\rho,
\end{align*}
are a special case of the Fourier transforms.
\par\emph{Note.}~This technique may be generalized to derive the
Hankel transforms of order $\nu=0,\tfrac{1}{2},1,\tfrac{3}{2},
\ldots$. See the two texts by Sneddon (Additional Readings). It
might also be noted that the Hankel transforms of half-integral
orders $\nu=\pm\tfrac{1}{2}$ reduce to Fourier sine and cosine
transforms.
\solution
Put $x=r\cos\theta$, $y=r\sin\theta$, and write
$u=\rho\cos\phi$, $v=\rho\sin\phi$. Then
$ux+vy=\rho r(\cos\theta\cos\phi+\sin\theta\sin\phi)
=\rho r\cos(\theta-\phi)$ and $dx\,dy=r\,dr\,d\theta$.
The forward transform is consequently
\begin{align*}
F(u,v)&=\frac1{2\pi}\int_0^\infty rf(r)dr
\int_0^{2\pi}e^{i\rho r\cos(\theta-\phi)}d\theta\\
&=\int_0^\infty rf(r)J_0(\rho r)dr.
\end{align*}
The angular integral is $2\pi J_0(\rho r)$, independent of
$\phi$: the integrand has period $2\pi$ in $\theta$, so the shift
$\theta-\phi\to\theta$ does not change the integral, and
$(2\pi)^{-1}\int_0^{2\pi}e^{iz\cos\theta}d\theta=J_0(z)$ is the
case $n=0$ of the integral evaluated in Exercise~20.2.11. This
proves that the transform is radial: $F=F(\rho)$.
For the inverse, use $du\,dv=\rho\,d\rho\,d\phi$:
\begin{align*}
f(x,y)&=\frac1{2\pi}\int_0^\infty\rho F(\rho)d\rho
\int_0^{2\pi}e^{-i\rho r\cos(\theta-\phi)}d\phi\\
&=\int_0^\infty\rho F(\rho)J_0(\rho r)d\rho,
\end{align*}
because the same steps give $2\pi J_0(-\rho r)$ for the angular
integral and $J_0$ is even. This is the claimed inverse Hankel
transform. The factors $r$ and $\rho$ are the two polar Jacobians;
the Fourier factors $1/(2\pi)$ cancel the angular integrals.
For radial Schwartz functions these operations (polar coordinates
and the order of integration) are justified directly. By Parseval's
theorem, $\iint|f|^2dx\,dy=\iint|F|^2du\,dv$; in polar coordinates
both sides carry the factor $2\pi$ of the angular integration, so
$\int_0^\infty|f(r)|^2r\,dr=\int_0^\infty|F(\rho)|^2\rho\,d\rho$.
The Hankel pair therefore extends, as the Fourier transform does,
to all radial $f$ with $\int_0^\infty|f(r)|^2r\,dr<\infty$.

\Needspace{5\baselineskip}
\par\bigskip
\noindent\textbf{20.2.16}\quad Show that the 3-D Fourier
exponential transform of a radially symmetric function may be
rewritten as a Fourier sine transform:
\[
\frac{1}{(2\pi)^{3/2}}
\erratum{\int_{\mathbb R^3} f(r)e^{i\mathbf{k}\cdot\mathbf{r}}\,d^{3}r}
=\frac{1}{k}\sqrt{\frac{2}{\pi}}\int_{0}^{\infty}
r f(r)\sin kr\,dr.
\]
\par\erratum{Take $k=|\mathbf k|>0$, and use a continuous limit at $k=0$ when it exists.}
\solution
Choose the $z$ axis along $\mathbf k$, with
$k=|\mathbf k|>0$, so that $\mathbf k\cdot\mathbf r=kr\cos\theta$.
In spherical coordinates,
\begin{align*}
\mathcal Ff(\mathbf k)
&=\frac1{(2\pi)^{3/2}}\int_0^\infty r^2f(r)dr
\int_0^{2\pi}d\phi\int_0^\pi e^{ikr\cos\theta}\sin\theta\,d\theta\\
&=\frac{2\pi}{(2\pi)^{3/2}}\int_0^\infty r^2f(r)dr
\int_{-1}^1 e^{ikru}du\\
&=\frac{4\pi}{(2\pi)^{3/2}}\int_0^\infty
r^2f(r)\frac{\sin kr}{kr}dr\\
&=\boxed{\frac1k\sqrt{\frac2\pi}\int_0^\infty rf(r)\sin kr\,dr}.
\end{align*}
In the second line $u=\cos\theta$, $du=-\sin\theta\,d\theta$; in
the third,
$\int_{-1}^1e^{ikru}du=(e^{ikr}-e^{-ikr})/(ikr)=2\sin kr/(kr)$;
in the last, $4\pi/(2\pi)^{3/2}=\sqrt{2/\pi}$ and $r^2/(kr)=r/k$.
Thus $k\mathcal Ff(k)$ is the unitary sine transform of $rf(r)$.
The marked volume element must be $d^3r$ over all space, as in
the published correction. If $\int_0^\infty r^2|f(r)|dr<\infty$,
the zero-frequency limit follows from the third line, since
$\sin(kr)/(kr)\to1$ with $|\sin(kr)/(kr)|\le1$ (dominated
convergence):
\[
\mathcal Ff(0)=\sqrt{\frac2\pi}\int_0^\infty r^2f(r)dr.
\]
Consequently the displayed $1/k$ does not create a singularity
when this finite limit exists.

\Needspace{9\baselineskip}
\subsection*{20.3 Properties of Fourier Transforms}

\Needspace{5\baselineskip}
\par\bigskip
\noindent\textbf{20.3.1}\quad Write the 1-D equivalents of the
equations for translation, scale change, sign change, and complex
conjugation that were given for 3-D transforms in Eqs.~(20.48) to
(20.51).
\solution
Write $g(\omega)=\mathcal Ff(\omega)$ with the positive
forward exponential of Eq.~(20.10). The four one-dimensional
identities are
\begin{align*}
\mathcal F[f(x-X)](\omega)&=e^{i\omega X}g(\omega),\\
\mathcal F[f(\alpha x)](\omega)&=\frac1{|\alpha|}g(\omega/\alpha),
\quad \alpha\in\mathbb R\setminus\{0\},\\
\mathcal F[f(-x)](\omega)&=g(-\omega),\\
\mathcal F[f^*(-x)](\omega)&=g^*(\omega).
\end{align*}
For translation set $u=x-X$, so $dx=du$ and
$e^{i\omega x}=e^{i\omega X}e^{i\omega u}$; the factor
$e^{i\omega X}$ comes out of the integral.
For scaling set $u=\alpha x$, so $dx=du/\alpha$ and
$e^{i\omega x}=e^{i(\omega/\alpha)u}$. When $\alpha<0$ the
integration limits reverse, and restoring their order changes
$1/\alpha$ into $-1/\alpha=1/|\alpha|$; in either case
\[
\mathcal F[f(\alpha x)](\omega)
=\frac1{|\alpha|}\frac1{\sqrt{2\pi}}\int f(u)e^{i(\omega/\alpha)u}du
=\frac1{|\alpha|}g(\omega/\alpha).
\]
The third identity is the second at $\alpha=-1$.
For the fourth, set $u=-x$ (the reversal of the limits cancels
the sign of $dx=-du$) and obtain
\[
\frac1{\sqrt{2\pi}}\int f^*(-x)e^{i\omega x}dx
=\frac1{\sqrt{2\pi}}\int f^*(u)e^{-i\omega u}du
=\left[\frac1{\sqrt{2\pi}}\int f(u)e^{i\omega u}du\right]^*
=g^*(\omega).
\]
Equivalently, applying the third identity to $f^*(-x)$,
$\mathcal F[f^*](\omega)=g^*(-\omega)$.
\erratum{A scale-change factor without an absolute value
presumes a positive scale. The formula above covers either sign;
$\alpha=0$ is not an invertible change of variables.}
\erratum{The one-dimensional translation line of the
Instructor's Manual, p.~337, prints $e^{i\omega t}$ for a shift
$f(t-a)$.
The phase must contain the fixed displacement, $e^{i\omega a}$,
not the integration variable.}

\Needspace{5\baselineskip}
\par\bigskip
\noindent\textbf{20.3.2}
\begin{enumerate}
\item[(a)] Show that by replacement of $\mathbf{r}$ by
$\mathbf{r}-\mathbf{R}$ in the formula for the Fourier transform
of $f(\mathbf{r})$, one can derive the translation formula,
Eq.~(20.48).
\item[(b)] Using methods similar to those for part (a), establish
the formulas for scale change, sign change, and complex
conjugation, Eqs.~(20.49) to (20.51).
\end{enumerate}
\solution
Let
$g(\mathbf k)=(2\pi)^{-3/2}\int f(\mathbf r)e^{i\mathbf k\cdot\mathbf r}d^3r$.
(a) Put $\mathbf u=\mathbf r-\mathbf R$. Translation has unit
Jacobian, and the exponential separates:
\begin{align*}
\mathcal F[f(\mathbf r-\mathbf R)](\mathbf k)
&=\frac1{(2\pi)^{3/2}}\int f(\mathbf u)
e^{i\mathbf k\cdot(\mathbf u+\mathbf R)}d^3u\\
&=\boxed{e^{i\mathbf k\cdot\mathbf R}g(\mathbf k)}.
\end{align*}
(b) For real nonzero $\alpha$, set $\mathbf u=\alpha\mathbf r$.
The volume element is $d^3r=|\alpha|^{-3}d^3u$ (the Jacobian of
$\mathbf r=\mathbf u/\alpha$ is $|\det(\alpha^{-1}I)|=|\alpha|^{-3}$),
and $\mathbf k\cdot\mathbf r=(\mathbf k/\alpha)\cdot\mathbf u$, so
\begin{align*}
\mathcal F[f(\alpha\mathbf r)](\mathbf k)
&=\frac{|\alpha|^{-3}}{(2\pi)^{3/2}}\int f(\mathbf u)
e^{i(\mathbf k/\alpha)\cdot\mathbf u}d^3u\\
&=\boxed{|\alpha|^{-3}g(\mathbf k/\alpha)}.
\end{align*}
In particular, at $\alpha=-1$ this gives
$\mathcal F[f(-\mathbf r)](\mathbf k)=g(-\mathbf k)$.
Finally, put $\mathbf u=-\mathbf r$, again with unit Jacobian,
so that $\mathbf k\cdot\mathbf r=-\mathbf k\cdot\mathbf u$:
\begin{align*}
\mathcal F[f^*(-\mathbf r)](\mathbf k)
&=\frac1{(2\pi)^{3/2}}\int f^*(\mathbf u)e^{-i\mathbf k\cdot\mathbf u}d^3u\\
&=\left[\frac1{(2\pi)^{3/2}}\int f(\mathbf u)e^{i\mathbf k\cdot\mathbf u}d^3u\right]^*
=\boxed{g^*(\mathbf k)}.
\end{align*}
\erratum{The printed $\alpha^{-3}$ in Eq.~(20.49) needs
$\alpha>0$. For negative real scaling the correct Jacobian is
$|\alpha|^{-3}$, as the sign-change case itself confirms.}

\Needspace{5\baselineskip}
\par\bigskip
\noindent\textbf{20.3.3}\quad Derive Eq.~(20.53), the formula for
the Fourier transform of $\nabla^{2}f(\mathbf{r})$.
\solution
For a sufficiently decaying differentiable function,
integrate by parts in $r_j$, holding the other coordinates fixed
($d^2r_\perp$ denotes their area element) and using
$\partial_je^{i\mathbf k\cdot\mathbf r}=ik_je^{i\mathbf k\cdot\mathbf r}$:
\begin{align*}
\mathcal F[\partial_j f](\mathbf k)
&=\frac1{(2\pi)^{3/2}}\int d^2r_\perp
\left([f e^{i\mathbf k\cdot\mathbf r}]_{r_j=-\infty}^{\infty}
-\int f\,ik_j e^{i\mathbf k\cdot\mathbf r}dr_j\right)\\
&=-ik_j g(\mathbf k).
\end{align*}
The boundary bracket vanishes because $f\to0$ as
$|r_j|\to\infty$. Applying the same calculation to
$\partial_j f$ gives
$\mathcal F[\partial_j^2 f]=-ik_j\mathcal F[\partial_jf]
=(-ik_j)^2g=-k_j^2g$.
Summing the three coordinate equations yields
\[
\boxed{\mathcal F[\nabla^2f](\mathbf k)
=-(k_x^2+k_y^2+k_z^2)g(\mathbf k)=-k^2g(\mathbf k)}.
\]
Schwartz functions satisfy the required endpoint and interchange
conditions. The same identity extends to tempered distributions
by defining their derivatives through integration by parts, so
singular sources can also be transformed consistently.

\Needspace{5\baselineskip}
\par\bigskip
\noindent\textbf{20.3.4}\quad Verify Eqs.~(20.55) and (20.56), the
formulas for the derivatives of 1-D Fourier transforms.
\solution
Starting from the one-dimensional definition and integrating by
parts with $de^{i\omega t}/dt=i\omega e^{i\omega t}$,
\begin{align*}
\mathcal F[f'](\omega)
&=\frac1{\sqrt{2\pi}}\int_{-\infty}^{\infty}f'(t)e^{i\omega t}dt\\
&=\frac1{\sqrt{2\pi}}[f(t)e^{i\omega t}]_{-\infty}^{\infty}
-i\omega\frac1{\sqrt{2\pi}}\int_{-\infty}^{\infty}f(t)e^{i\omega t}dt\\
&=-i\omega g(\omega),
\end{align*}
since the bracket vanishes when $f(t)\to0$ as $t\to\pm\infty$.
For $n$ derivatives repeat the operation, applying it to
$f^{(n-1)},f^{(n-2)},\ldots$. Each repetition moves
one derivative to the exponential and contributes $-i\omega$:
$\mathcal F[f^{(n)}]=-i\omega\mathcal F[f^{(n-1)}]=\cdots$, so
\[
\boxed{\mathcal F[f^{(n)}](\omega)=(-i\omega)^n g(\omega)}.
\]
More explicitly, $n$ integrations by parts on $(-R,R)$ give
\begin{align*}
\int_{-R}^{R}f^{(n)}(t)e^{i\omega t}dt
&=\sum_{j=0}^{n-1}(-i\omega)^j
[f^{(n-1-j)}(t)e^{i\omega t}]_{-R}^{R}\\
&\quad+(-i\omega)^n\int_{-R}^{R}f(t)e^{i\omega t}dt;
\end{align*}
under the usual decaying-function assumptions ($f,\ldots,f^{(n-1)}$
tend to zero at $\pm\infty$, and $f$ and $f^{(n)}$ are integrable)
each endpoint term tends to zero
as $R\to\infty$. This explains why no initial-value terms occur
for this two-sided transform, unlike the one-sided Laplace
transform. If $f$ has a jump at $t_0$, integrating by parts
separately on each side of $t_0$ leaves the extra endpoint term
$-[f(t_0^+)-f(t_0^-)]e^{i\omega t_0}/\sqrt{2\pi}$ in the transform
of the pointwise derivative, so that derivative alone violates the
identity. The identity still holds when $f'$ is the distributional
derivative, which contains the term
$[f(t_0^+)-f(t_0^-)]\,\delta(t-t_0)$; its transform
$[f(t_0^+)-f(t_0^-)]e^{i\omega t_0}/\sqrt{2\pi}$ cancels the endpoint
term. Jumps of $f',\ldots,f^{(n-1)}$ are treated in the same way.

\Needspace{5\baselineskip}
\par\bigskip
\noindent\textbf{20.3.5}\quad Derive the inverse of Eq.~(20.56),
namely that
\[
\left[t^{n}f(t)\right]^{T}\!\!(\omega)
=i^{-n}\,\frac{d^{n}}{d\omega^{n}}\,g(\omega).
\]
\solution
Differentiate the forward transform with respect to
its frequency argument under the integral sign; since
$\partial e^{i\omega t}/\partial\omega=it\,e^{i\omega t}$,
\[
g'(\omega)=\frac1{\sqrt{2\pi}}\int f(t)
\frac\partial{\partial\omega}e^{i\omega t}dt
=\frac i{\sqrt{2\pi}}\int tf(t)e^{i\omega t}dt
=i\,\mathcal F[tf](\omega).
\]
After $n$ derivatives, each bringing down another factor $it$,
\[
g^{(n)}(\omega)=\frac1{\sqrt{2\pi}}\int (it)^nf(t)e^{i\omega t}dt
=i^n\mathcal F[t^nf](\omega).
\]
Dividing by $i^n$ proves
\[
\boxed{\mathcal F[t^nf](\omega)=i^{-n}g^{(n)}(\omega)}.
\]
For example, $t^nf\in L^1$ justifies these differentiations
by dominated convergence, with lower moments controlled by
$f\in L^1$ and the highest moment. Schwartz functions satisfy
all such conditions. This is the duality between multiplication
in direct space and differentiation in frequency space.

\Needspace{5\baselineskip}
\par\bigskip
\noindent\textbf{20.3.6}\quad The 1-D neutron diffusion equation
with a (plane) source is
\[
-D\,\frac{d^{2}\varphi(x)}{dx^{2}}+K^{2}D\varphi(x)=Q\delta(x),
\]
where $\varphi(x)$ is the neutron flux, $Q\delta(x)$ is the
(plane) source at $x=0$, and $D$ and $K^{2}$ are constants. Apply
a Fourier transform. Solve the equation in transform space.
Transform your solution back into $x$-space.
\par\nobreak\hfill\emph{ANS.}\enspace
$\varphi(x)=\dfrac{Q}{2KD}\,e^{-|Kx|}$.
\par\erratum{For the displayed answer assume $D>0$, $K>0$, and seek the whole-line solution bounded at both infinities. These conditions select the decaying solution.}
\solution
\erratum{For the decaying whole-line solution take $D>0$
and $K^2>0$, and put $\kappa=\sqrt{K^2}>0$. The printed
coefficient $1/K$ presumes $K>0$. Decay, or boundedness at both
infinities, must be part of the boundary specification.}
Write $\Phi=\mathcal F\varphi$. The delta transform is
$\mathcal F[\delta](\omega)=(2\pi)^{-1/2}\int\delta(x)e^{i\omega x}dx
=1/\sqrt{2\pi}$, and
$\mathcal F[\varphi'']=(-i\omega)^2\Phi=-\omega^2\Phi$ by
Exercise~20.3.4 with $n=2$. Transforming the equation term by
term gives $-D(-\omega^2\Phi)+\kappa^2D\Phi=Q/\sqrt{2\pi}$, i.e.,
\[
D(\omega^2+\kappa^2)\Phi(\omega)=\frac Q{\sqrt{2\pi}},
\qquad \Phi(\omega)=\frac{Q}{D\sqrt{2\pi}(\omega^2+\kappa^2)}.
\]
Invert with $\varphi(x)=(2\pi)^{-1/2}\int\Phi(\omega)e^{-i\omega x}d\omega$.
Because $1/(\omega^2+\kappa^2)$ is even, the odd part
$-i\sin\omega x$ of $e^{-i\omega x}$ integrates to zero and the
cosine part gives twice its integral over $(0,\infty)$. The second
formula of Exercise~20.2.3(b), with $a=\kappa$ and $x$ replaced by
$|x|$ (the cosine integral is even in $x$, and at $x=0$ it equals
$\int_0^\infty d\omega/(\omega^2+\kappa^2)=\pi/(2\kappa)$), then gives
$\int_0^\infty\cos\omega x\,(\omega^2+\kappa^2)^{-1}d\omega
=(\pi/2\kappa)e^{-\kappa|x|}$:
\begin{align*}
\varphi(x)&=\frac{Q}{2\pi D}\int_{-\infty}^{\infty}
\frac{e^{-i\omega x}}{\omega^2+\kappa^2}d\omega\\
&=\frac{Q}{\pi D}\int_0^\infty
\frac{\cos\omega x}{\omega^2+\kappa^2}d\omega
=\boxed{\frac{Q}{2\kappa D}e^{-\kappa|x|}}.
\end{align*}
As a check, away from zero $\varphi''=\kappa^2\varphi$. For
$x\gtrless0$, $\varphi'=\mp\kappa\varphi$, so
$\varphi'(0^\pm)=\mp Q/(2D)$ and the derivative jump is
\[
\varphi'(0^+)-\varphi'(0^-)
=-\frac{Q}{2D}-\frac{Q}{2D}=-\frac QD.
\]
As a distribution, therefore,
$\varphi''=\kappa^2\varphi+[\varphi'(0^+)-\varphi'(0^-)]\delta(x)
=\kappa^2\varphi-(Q/D)\delta(x)$, so
$-D\varphi''+\kappa^2D\varphi=Q\delta(x)$ exactly.
A homogeneous addition is $A e^{\kappa x}+B e^{-\kappa x}$;
boundedness at both infinities forces $A=B=0$.
For $K=0$ and nonzero $Q$, the formal fundamental solution is
$-Q|x|/(2D)$ plus an affine function (since $|x|''=2\delta(x)$),
and no solution decays at
both ends. Substituting $K=0$ in the printed formula is invalid.

\Needspace{9\baselineskip}
\subsection*{20.4 Fourier Convolution Theorem}

\Needspace{5\baselineskip}
\par\bigskip
\noindent\textbf{20.4.1}\quad Work out the convolution equation
corresponding to Eq.~(20.71) for
\begin{enumerate}
\item[(a)] Fourier sine transforms
\[
\frac{1}{2}\int_{0}^{\infty}g(y)\big[f(y+x)+f(y-x)\big]\,dy
=\int_{0}^{\infty}F_{s}(s)G_{s}(s)\cos sx\,ds,
\]
where $f$ and $g$ are odd functions.
\item[(b)] Fourier cosine transforms
\[
\frac{1}{2}\int_{0}^{\infty}g(y)\big[f(y+x)+f(x-y)\big]\,dy
=\int_{0}^{\infty}F_{c}(s)G_{c}(s)\cos sx\,ds,
\]
where $f$ and $g$ are even functions.
\end{enumerate}
\solution
Use the unitary half-line transforms, for example
$F_s(s)=\sqrt{2/\pi}\int_0^\infty f(u)\sin(su)du$,
with the identical normalization in the inverse.

(a) Extend $f$ oddly. Its inverse sine representation
\[
f(u)=\sqrt{\frac2\pi}\int_0^\infty F_s(s)\sin(su)\,ds
\]
then holds for all real $u$, since both sides are odd in $u$.
Apply it at $u=y+x$ and $u=y-x$, and use the identity
$\sin(A+B)+\sin(A-B)=2\sin A\cos B$ with $A=sy$, $B=sx$:
\[
f(y+x)+f(y-x)=2\sqrt{\frac2\pi}\int_0^\infty
F_s(s)\sin(sy)\cos(sx)ds.
\]
Multiply by $g(y)/2$, integrate on $y>0$, and interchange the
order of integration:
\begin{align*}
\frac12\int_0^\infty g(y)[f(y+x)+f(y-x)]dy
&=\sqrt{\frac2\pi}\int_0^\infty F_s(s)\cos(sx)
\left[\int_0^\infty g(y)\sin(sy)dy\right]ds\\
&=\int_0^\infty F_s(s)G_s(s)\cos(sx)ds.
\end{align*}
By the definition of $G_s$ the bracket is
$\sqrt{\pi/2}\,G_s(s)$, and the factors $\sqrt{2/\pi}$ and
$\sqrt{\pi/2}$ cancel.

(b) Extend $f$ evenly; then
$f(u)=\sqrt{2/\pi}\int_0^\infty F_c(s)\cos(su)\,ds$ for all
real $u$. Since $f(x-y)=f(y-x)$, the identity
$\cos(A+B)+\cos(A-B)=2\cos A\cos B$, again with $A=sy$,
$B=sx$, gives
\[
f(y+x)+f(x-y)=2\sqrt{\frac2\pi}\int_0^\infty
F_c(s)\cos(sy)\cos(sx)ds.
\]
The same multiplication and integration replace the inner
integral $\int_0^\infty g(y)\cos(sy)dy$ by $\sqrt{\pi/2}\,G_c(s)$,
giving exactly the second
formula. The calculation is justified first for sufficiently
regular decaying transforms, and then by the $L^2$ pairing
limit where appropriate; odd/even extensions are essential for
arguments that become negative.
\erratum{The Instructor's Manual's intermediate two-delta
kernels on p.~339, which arise when both transforms on the
right-hand side are written as integrals and the $s$ integral is
done first, omit reflected sources. For example, the sine kernel
on the full range of its arguments is}
\[
\erratum{\begin{aligned}
\int_0^\infty\sin(sy)\sin(st)\cos(sx)ds
=\frac\pi4\big[&\delta(t-y-x)+\delta(t-y+x)\\
&-\delta(t+y+x)-\delta(t+y-x)\big].
\end{aligned}}
\]
\erratum{This follows by applying the product identity and
the reflected-source formula in Exercise~20.2.7.}
Indeed, $\sin(sy)\sin(st)=\frac12[\cos s(t-y)-\cos s(t+y)]$,
and multiplying each cosine by $\cos(sx)$ with
$\cos\alpha\cos\beta=\frac12[\cos(\alpha+\beta)+\cos(\alpha-\beta)]$
turns the integrand into
\[
\tfrac14\big[\cos s(t-y+x)+\cos s(t-y-x)
-\cos s(t+y+x)-\cos s(t+y-x)\big];
\]
each term is then integrated with
$\int_0^\infty\cos(su)\,ds=\pi\delta(u)$ (Exercise~20.2.7).
\erratum{The cosine
kernel has the same four delta arguments with all plus signs.
The reflected terms cannot all be discarded on $t,y>0$ when
$x$ is arbitrary; the odd/even extensions above include them.}
Since $F_s(s)G_s(s)$ is $2/\pi$ times the integral of
$f(t)g(y)\sin(st)\sin(sy)$ over $t,y>0$, inserting the sine
kernel and doing the $t$ integral, in which each delta
contributes only where its root $t$ is positive ($\theta$ is the
unit step function), gives
\[
\begin{aligned}
\int_0^\infty F_sG_s\cos(sx)\,ds
=\frac12\int_0^\infty g(y)\big[&f(y+x)\theta(y+x)+f(y-x)\theta(y-x)\\
&-f(-y-x)\theta(-y-x)-f(x-y)\theta(x-y)\big]dy.
\end{aligned}
\]
For odd $f$, $-f(-y-x)=f(y+x)$ and $-f(x-y)=f(y-x)$, so, since
$\theta(u)+\theta(-u)=1$, the bracket is $f(y+x)+f(y-x)$, which
gives (a); without the two
reflected terms the contributions from $y+x<0$ and from $y<x$
would be lost. The cosine kernel, with all plus signs and even
$f$, gives $f(y+x)+f(x-y)$ in the same way.

\Needspace{5\baselineskip}
\par\bigskip
\noindent\textbf{20.4.2}\quad Show that for both Fourier sine and
Fourier cosine transforms Parseval's relation has the form
\[
\int_{0}^{\infty}F(t)G(t)\,dt=\int_{0}^{\infty}f(y)g(y)\,dy.
\]
\solution
Set $x=0$ in each of the two formulas just proved.
For the sine version the bracket is $f(y)+f(y)=2f(y)$.
For the cosine version it is $f(y)+f(-y)=2f(y)$ because the
extension is even. On either right-hand side $\cos(sx)=1$.
Thus in both cases
\[
\boxed{\int_0^\infty f(y)g(y)dy
=\int_0^\infty F(s)G(s)ds}.
\]
The same result follows by inserting the inverse transform of
$f$ into the left-hand side and identifying the transform of
$g$; for the sine case,
\begin{align*}
\int_0^\infty f(y)g(y)dy
&=\int_0^\infty g(y)\sqrt{\frac2\pi}\int_0^\infty F(s)\sin(sy)\,ds\,dy\\
&=\int_0^\infty F(s)\left[\sqrt{\frac2\pi}\int_0^\infty
g(y)\sin(sy)\,dy\right]ds
=\int_0^\infty F(s)G(s)\,ds,
\end{align*}
and the cosine case is identical with $\cos(sy)$; all
normalization factors cancel. For $L^2$ functions,
Cauchy--Schwarz controls both pairings and permits approximation
by regular functions. This displayed version is bilinear and is
valid also for complex functions. To obtain the norm form set
$g=\overline f$: the sine and cosine kernels are real, so
$G=\overline F$, and
$\int|f|^2=\int|F|^2$.

\Needspace{5\baselineskip}
\par\bigskip
\noindent\textbf{20.4.3}
\begin{enumerate}
\item[(a)] A rectangular pulse is described by
\[
f(x)=\begin{cases}
1, & |x| < a,\\
0, & |x| > a.
\end{cases}
\]
Show that the Fourier exponential transform is
\[
F(t)=\sqrt{\frac{2}{\pi}}\,\frac{\sin at}{t}.
\]
This is the single-slit diffraction problem of physical optics.
The slit is described by $f(x)$. The diffraction pattern amplitude
is given by the Fourier transform $F(t)$.
\item[(b)] Use the Parseval relation to evaluate
\[
\int_{-\infty}^{\infty}\frac{\sin^{2}t}{t^{2}}\,dt.
\]
This integral may also be evaluated by using the calculus of
residues (Exercise~11.8.9).
\end{enumerate}
\par\nobreak\hfill\emph{ANS.}\enspace(b) $\pi$.
\par\erratum{Take $a>0$.}
\solution
(a) For $a>0$, direct integration gives
\begin{align*}
F(t)&=\frac1{\sqrt{2\pi}}\int_{-a}^{a}e^{itx}dx
=\frac1{\sqrt{2\pi}}\frac{e^{iat}-e^{-iat}}{it}\\
&=\frac1{\sqrt{2\pi}}\frac{2\sin at}{t}
=\boxed{\sqrt{\frac2\pi}\frac{\sin at}{t}},
\end{align*}
using $e^{i\theta}-e^{-i\theta}=2i\sin\theta$ and
$2/\sqrt{2\pi}=\sqrt{2/\pi}$.
At $t=0$ take the limit $F(0)=a\sqrt{2/\pi}$, since
$\sin at/t\to a$.
(b) Parseval's relation in the unitary form
$\int|f(x)|^2dx=\int|F(t)|^2dt$ gives
\[
2a=\int_{-a}^a1\,dx
=\int_{-\infty}^{\infty}|F(t)|^2dt
=\frac2\pi\int_{-\infty}^{\infty}\frac{\sin^2at}{t^2}dt.
\]
Set $u=at$, so $dt=du/a$ and $t^{-2}=a^2/u^2$:
\[
2a=\frac{2a}{\pi}\int_{-\infty}^{\infty}\frac{\sin^2u}{u^2}du.
\]
Canceling $2a$ gives the required value $\boxed{\pi}$.
The apparent singularity at zero is removable and the tails are
bounded by $u^{-2}$, so the final integral is absolutely convergent.

\Needspace{5\baselineskip}
\par\bigskip
\noindent\textbf{20.4.4}\quad Solve Poisson's equation,
$\nabla^{2}\psi(\mathbf{r})=-\rho(\mathbf{r})/\varepsilon_{0}$, by
the following sequence of operations:
\begin{enumerate}
\item[(a)] Take the Fourier transform of both sides of this
equation. Solve for the Fourier transform of $\psi(\mathbf{r})$.
\item[(b)] Carry out the Fourier inverse transform.
\end{enumerate}
\par\erratum{Seek the free-space solution that vanishes at infinity, for sufficiently localized charge. Without this boundary condition, harmonic additions are possible.}
\solution
(a) Let $\widehat\psi$ and $\widehat\rho$ denote
unitary three-dimensional transforms. By Exercise~20.3.3,
$\mathcal F[\nabla^2\psi]=-k^2\widehat\psi$, so the transformed
equation is
\[
-k^2\widehat\psi(\mathbf k)=-\frac{\widehat\rho(\mathbf k)}{\varepsilon_0},
\qquad
\widehat\psi(\mathbf k)=\frac{\widehat\rho(\mathbf k)}{\varepsilon_0k^2}.
\]
(b) To invert the multiplier, first evaluate the absolutely
convergent regulated radial transform, $\eta>0$, with the radial
formula of Exercise~20.2.16 (here $rf(r)=e^{-\eta r}$) and the
integral $S$ of Exercise~20.2.3 (with $a=\eta$, $\omega=k$):
\begin{align*}
\mathcal F\left[\frac{e^{-\eta r}}r\right](\mathbf k)
&=\sqrt{\frac2\pi}\frac1k\int_0^\infty e^{-\eta r}\sin kr\,dr
=\sqrt{\frac2\pi}\frac1k\,\frac{k}{k^2+\eta^2}\\
&=\sqrt{\frac2\pi}\frac1{k^2+\eta^2}.
\end{align*}
Letting $\eta\downarrow0$ in distributions gives
$\mathcal F[1/r]=\sqrt{2/\pi}/k^2$, so by Fourier inversion
$\mathcal F^{-1}[\sqrt{2/\pi}/k^2]=1/r$: the
inverse transform of $k^{-2}$ is $\sqrt{\pi/2}/r$.
For our normalization, if $A=\mathcal Fa$ and $B=\mathcal Fb$, the
convolution theorem reads
\[
\mathcal F^{-1}[AB](\mathbf r)
=(2\pi)^{-3/2}\int a(\mathbf u)b(\mathbf r-\mathbf u)d^3u.
\]
Indeed, transforming the right-hand side and substituting
$\mathbf r=\mathbf u+\mathbf v$ at fixed $\mathbf u$ factorizes it:
\[
\frac1{(2\pi)^3}\int a(\mathbf u)e^{i\mathbf k\cdot\mathbf u}d^3u
\int b(\mathbf v)e^{i\mathbf k\cdot\mathbf v}d^3v
=\frac{(2\pi)^{3/2}A\,(2\pi)^{3/2}B}{(2\pi)^3}=AB.
\]
Use $a=\rho$ and $b=\sqrt{\pi/2}/r$, so that $B=k^{-2}$:
\[
\psi(\mathbf r)=\frac1{\varepsilon_0}\,
\mathcal F^{-1}\!\left[\frac{\widehat\rho}{k^2}\right]
=\frac1{\varepsilon_0}\,\frac{\sqrt{\pi/2}}{(2\pi)^{3/2}}
\int\frac{\rho(\mathbf u)}{|\mathbf r-\mathbf u|}d^3u.
\]
Since
$\sqrt{\pi/2}/(2\pi)^{3/2}=\sqrt\pi/(\sqrt2\cdot2\sqrt2\,\pi^{3/2})
=1/(4\pi)$,
\[
\boxed{\psi(\mathbf r)=\frac1{4\pi\varepsilon_0}
\int\frac{\rho(\mathbf u)}{|\mathbf r-\mathbf u|}d^3u}.
\]
\erratum{This is the free-space solution with vanishing
potential at infinity for suitably localized charge. Without a
boundary condition, one may add any harmonic function; division
by $k^2$ does not establish uniqueness at the zero mode.}
For smooth localized charge the Coulomb singularity is locally
integrable, and $\nabla^2(1/r)=-4\pi\delta^{(3)}$ verifies the
source normalization of the answer.

\Needspace{5\baselineskip}
\par\bigskip
\noindent\textbf{20.4.5}
\begin{enumerate}
\item[(a)] Given $f(x)=1-|x/2|$ for $-2\le x\le2$, with $f(x)=0$
elsewhere, show that the Fourier transform of $f(x)$ is
\[
F(t)=\sqrt{\frac{2}{\pi}}\left(\frac{\sin t}{t}\right)^{2}.
\]
\item[(b)] Using the Parseval relation, evaluate
\[
\int_{-\infty}^{\infty}\left(\frac{\sin t}{t}\right)^{4}dt.
\]
\end{enumerate}
\par\nobreak\hfill\emph{ANS.}\enspace(b) $\dfrac{2\pi}{3}$.
\solution
(a) Let $q(x)$ equal one on $[-1,1]$ and zero
elsewhere. In this solution $*$ denotes the unnormalized
convolution integral, so the Fourier convolution theorem includes
its factor $\sqrt{2\pi}$. In
$(q*q)(x)=\int q(y)q(x-y)dy$ the integrand is one exactly when
$y\in[-1,1]$ and $y\in[x-1,x+1]$. For $0\le x\le2$ the overlap
of the two intervals is $[x-1,1]$, of length $2-x$; for
$-2\le x\le0$ it is $[-1,x+1]$, of length $2+x$; for $|x|>2$
the intervals do not overlap. Hence $(q*q)(x)=2-|x|$ for
$|x|\le2$ and zero outside, and $f=(q*q)/2$.
From Exercise~20.4.3 with $a=1$, $Q(t)=\sqrt{2/\pi}\sin t/t$.
The unitary convolution theorem reads
$\mathcal F[q*q]=\sqrt{2\pi}\,Q^2$: substituting $x=y+z$ in the
transform of $q*q$ gives
$(2\pi)^{-1/2}\int q(y)e^{ity}dy\int q(z)e^{itz}dz
=(2\pi)^{-1/2}(\sqrt{2\pi}\,Q)^2$. Therefore
\[
F(t)=\frac12\sqrt{2\pi}\,Q(t)^2
=\frac12\sqrt{2\pi}\,\frac2\pi\left(\frac{\sin t}{t}\right)^2
=\boxed{\sqrt{\frac2\pi}\left(\frac{\sin t}{t}\right)^2}.
\]
(b) Since $f$ is even, its direct-space squared norm is twice the
integral over $0\le x\le2$; there put $u=1-x/2$, $dx=-2\,du$, so
$u$ runs from $1$ to $0$, and reversing the limits absorbs the minus
sign:
\[
\int|f(x)|^2dx=2\int_0^2(1-x/2)^2dx
=4\int_0^1u^2du=\frac43.
\]
In transform space it is
\[
\int|F(t)|^2dt=\frac2\pi\int_{-\infty}^{\infty}
\left(\frac{\sin t}{t}\right)^4dt.
\]
Equating the norms yields the value
$\boxed{(\pi/2)(4/3)=2\pi/3}$.
\erratum{The Instructor's Manual's Parseval line on p.~341 prints
$(\sin t/t)^2$ where it must print $(\sin t/t)^4$:
the already squared amplitude $F(t)$ is squared again to
form $|F(t)|^2$.}

\Needspace{5\baselineskip}
\par\bigskip
\noindent\textbf{20.4.6}\quad With $F(t)$ and $G(t)$ the Fourier
transforms of $f(x)$ and $g(x)$, respectively, show that
\[
\int_{-\infty}^{\infty}\big|f(x)-g(x)\big|^{2}dx
=\int_{-\infty}^{\infty}\big|F(t)-G(t)\big|^{2}dt.
\]
If $g(x)$ is an approximation to $f(x)$, the preceding relation
indicates that the mean square deviation in $t$-space is equal to
the mean square deviation in $x$-space.
\par\erratum{Assume $f,g\in L^2$, or more generally that the difference and its transform are defined in $L^2$.}
\solution
Let $h=f-g$. Linearity gives
$\mathcal Fh=F-G$. Parseval's norm identity applied to $h$ is
\[
\int|h(x)|^2dx=\int|\mathcal Fh(t)|^2dt,
\]
which is exactly the claimed equation. To display the complex
cross terms explicitly,
\[
|f-g|^2=|f|^2+|g|^2-f\overline g-\overline f g.
\]
The unitary inner-product identity
$\int u\overline v\,dx=\int U\overline V\,dt$ sends each of these
four integrals respectively to the integrals of $|F|^2$, $|G|^2$,
$F\overline G$, and $\overline F G$, and with the same signs
$|F|^2+|G|^2-F\overline G-\overline FG=|F-G|^2$.
\erratum{The finite mean-square statement requires $f-g\in L^2$
and its $L^2$ transform; for example $f,g\in L^2$ suffice.
It is an equality of integrated squared errors, not pointwise
errors at corresponding coordinates.}

\Needspace{5\baselineskip}
\par\bigskip
\noindent\textbf{20.4.7}\quad Use the Parseval relation to
evaluate
\[
\text{(a)}\ \int_{-\infty}^{\infty}
\frac{d\omega}{(\omega^{2}+a^{2})^{2}},
\qquad
\text{(b)}\ \int_{-\infty}^{\infty}
\frac{\omega^{2}\,d\omega}{(\omega^{2}+a^{2})^{2}}.
\]
\par\emph{Hint.}~Compare Exercise~20.2.3.
\par\nobreak\hfill\emph{ANS.}\enspace
(a) $\dfrac{\pi}{2a^{3}}$,\quad(b) $\dfrac{\pi}{2a}$.
\par\erratum{The displayed answers use $a>0$; for nonzero real $a$ replace their $a$ by $|a|$.}
\solution
Take $a>0$ and apply the half-line Parseval identities
to $f(x)=e^{-ax}$. Its direct squared norm is
\[
\int_0^\infty e^{-2ax}dx=\frac1{2a}.
\]
(a) The cosine transform of Exercise~20.2.3 is
$g_c(\omega)=\sqrt{2/\pi}\,a/(\omega^2+a^2)$, so Parseval's
relation of Exercise~20.4.2, applied with both functions equal to
$f$ (so that both transforms are $g_c$), gives
\[
\frac1{2a}=\int_0^\infty g_c(\omega)^2d\omega
=\frac2\pi\int_0^\infty
\frac{a^2}{(\omega^2+a^2)^2}d\omega.
\]
Dividing by $2a^2/\pi$ gives
$\int_0^\infty d\omega/(\omega^2+a^2)^2=\pi/(4a^3)$. The
integrand is even, so doubling the half-line integral gives
\[
\boxed{\int_{-\infty}^{\infty}\frac{d\omega}{(\omega^2+a^2)^2}
=\frac\pi{2a^3}}.
\]
(b) Using the sine transform
$g_s(\omega)=\sqrt{2/\pi}\,\omega/(\omega^2+a^2)$ instead,
\[
\frac1{2a}=\int_0^\infty g_s(\omega)^2d\omega
=\frac2\pi\int_0^\infty
\frac{\omega^2}{(\omega^2+a^2)^2}d\omega,
\]
so the half-line integral is $\pi/(4a)$ and, the integrand being
even, the whole-line answer is $\boxed{\pi/(2a)}$.
\erratum{If $a$ is allowed to be negative, the integrands are
unchanged and the answers contain $|a|^3$ and $|a|$, respectively.
At $a=0$ both integrals diverge. The printed forms use $a>0$.}

\Needspace{5\baselineskip}
\par\bigskip
\noindent\textbf{20.4.8}\quad The nuclear form factor
$F(\mathbf{k})$ and the charge distribution $\rho(\mathbf{r})$
are 3-D Fourier transforms of each other:
\[
F(\mathbf{k})=\frac{1}{(2\pi)^{3/2}}\int
\rho(\mathbf{r})e^{i\mathbf{k}\cdot\mathbf{r}}\,d^{3}r.
\]
If the measured form factor is
\[
F(\mathbf{k})=(2\pi)^{-3/2}\left(1+\frac{k^{2}}{a^{2}}\right)^{-1},
\]
find the corresponding charge distribution.
\par\nobreak\hfill\emph{ANS.}\enspace
$\rho(r)=\dfrac{a^{2}}{4\pi}\,\dfrac{e^{-ar}}{r}$.
\par\erratum{The parameter $a$ in the displayed answer is the positive scale, $a>0$.}
\solution
Assume $a>0$ first. Rewrite the measured transform as
\[
F(\mathbf k)=\frac{a^2}{(2\pi)^{3/2}}\frac1{k^2+a^2}.
\]
The radial formula of Exercise~20.2.16, with $rf(r)=e^{-ar}$ and
$\int_0^\infty e^{-ar}\sin kr\,dr=k/(k^2+a^2)$ (the integral $S$
of Exercise~20.2.3 with $\omega=k$), gives
\[
\mathcal F\left[\frac{e^{-ar}}r\right](\mathbf k)
=\sqrt{\frac2\pi}\frac1k\int_0^\infty e^{-ar}\sin kr\,dr
=\sqrt{\frac2\pi}\frac1{k^2+a^2}.
\]
Therefore multiplication by $a^2/(4\pi)$ produces the given
$F$, because $(4\pi)^{-1}\sqrt{2/\pi}=(2\pi)^{-3/2}$.
Since a function is determined by its Fourier transform (the
inversion theorem), the charge distribution is
\[
\boxed{\rho(r)=\frac{a^2}{4\pi}\frac{e^{-ar}}r}.
\]
Its total charge is
\[
4\pi\int_0^\infty r^2\rho(r)dr
=a^2\int_0^\infty re^{-ar}dr=a^2\cdot\frac1{a^2}=1,
\]
in agreement with $\int\rho\,d^3r=(2\pi)^{3/2}F(0)=1$ (put
$\mathbf k=0$ in the definition of $F$). The $1/r$ origin
singularity is integrable with the volume element.
\erratum{The scale in the printed decaying exponential is the
positive root. For nonzero real $a$ of either sign use
$e^{-|a|r}$; the form factor only determines $a^2$.}

\Needspace{5\baselineskip}
\par\bigskip
\noindent\textbf{20.4.9}\quad Using convolution methods, find an
integral whose value is the electrostatic interaction energy
between a charge distribution $\rho(\mathbf{r}-\mathbf{A})$ and a
unit point charge at $\mathbf{C}$.
\solution
Let $H(\mathbf r)=1/|\mathbf r|$, write
$(\rho*H)(\mathbf v)=\int\rho(\mathbf u)H(\mathbf v-\mathbf u)d^3u$
for the unnormalized convolution, and use SI normalization, so the potential of unit charge is
$H/(4\pi\varepsilon_0)$. The potential at $\mathbf C$ of
the translated distribution is its interaction energy with the
unit point charge:
\[
U=\frac1{4\pi\varepsilon_0}\int
\frac{\rho(\mathbf r-\mathbf A)}{|\mathbf C-\mathbf r|}d^3r.
\]
Putting $\mathbf u=\mathbf r-\mathbf A$ identifies a convolution:
\[
\boxed{U=\frac1{4\pi\varepsilon_0}(\rho*H)(\mathbf C-\mathbf A)
=\frac1{4\pi\varepsilon_0}\int
\frac{\rho(\mathbf u)}{|\mathbf C-\mathbf A-\mathbf u|}d^3u}.
\]
If $R(\mathbf k)=\mathcal F\rho(\mathbf k)$, then
$\mathcal FH=\sqrt{2/\pi}/k^2$ (Exercise~20.4.4), and the
convolution theorem of Exercise~20.4.4, in the form
$\mathcal F(\rho*H)=(2\pi)^{3/2}R\,\mathcal FH$, gives
$\mathcal F(\rho*H)=4\pi R/k^2$, because
$(2\pi)^{3/2}\sqrt{2/\pi}=2^{3/2}\pi^{3/2}2^{1/2}\pi^{-1/2}=4\pi$.
Inverting,
\[
(\rho*H)(\mathbf v)=\frac1{(2\pi)^{3/2}}
\int\frac{4\pi R(\mathbf k)}{k^2}e^{-i\mathbf k\cdot\mathbf v}d^3k,
\]
and multiplying by $1/(4\pi\varepsilon_0)$ at
$\mathbf v=\mathbf C-\mathbf A$ yields the equivalent
transform-space integral
\[
\boxed{U=\frac1{\varepsilon_0(2\pi)^{3/2}}
\int\frac{R(\mathbf k)}{k^2}
e^{-i\mathbf k\cdot(\mathbf C-\mathbf A)}d^3k}.
\]
Both formulas show that only the relative displacement enters.
For real charge the frequency-domain integrand has conjugate
symmetry ($R(-\mathbf k)=R(\mathbf k)^*$, so its value at
$-\mathbf k$ is the complex conjugate of its value at
$\mathbf k$), so $U$ is real. These formulas assume the potential
integral exists; point-charge singularities are understood through
the corresponding distributions. There is no factor $1/2$ in
the mutual energy of these two specified charge distributions.

\Needspace{5\baselineskip}
\par\bigskip
\noindent\textbf{20.4.10}\quad With $\psi(\mathbf{r})$ a wave
function in ordinary space and $\varphi(\mathbf{p})$ the
corresponding momentum function, show that
\begin{enumerate}
\item[(a)] $\displaystyle\frac{1}{(2\pi\hbar)^{3/2}}
\int\mathbf{r}\psi(\mathbf{r})
e^{-i\mathbf{r}\cdot\mathbf{p}/\hbar}\,d^{3}r
=i\hbar\nabla_{p}\,\varphi(\mathbf{p})$,
\item[(b)] $\displaystyle\frac{1}{(2\pi\hbar)^{3/2}}
\int\mathbf{r}^{2}\psi(\mathbf{r})
e^{\erratum{-i}\mathbf{r}\cdot\mathbf{p}/\hbar}\,d^{3}r
=\big(i\hbar\nabla_{p}\big)^{2}\varphi(\mathbf{p})$.
\end{enumerate}
\par\erratum{(This exercise assumes that the units and scaling of
the momentum wave function correspond to the formula
$\varphi(\mathbf{p})=\dfrac{1}{(2\pi\hbar)^{3/2}}
\displaystyle\int\psi(\mathbf{r})
e^{-i\mathbf{r}\cdot\mathbf{p}/\hbar}\,d^{3}r$.)}
\par\emph{Note.}~$\nabla_{p}$ is the gradient in momentum space:
\[
\hat{\mathbf{e}}_{x}\frac{\partial}{\partial p_{x}}
+\hat{\mathbf{e}}_{y}\frac{\partial}{\partial p_{y}}
+\hat{\mathbf{e}}_{z}\frac{\partial}{\partial p_{z}}.
\]
These results may be extended to any positive integer power of
$\mathbf{r}$ \erratum{on the appropriate moment domain, and to an
analytic function only when its Maclaurin series can be
interchanged with the transform for the wave function in use}.
\solution
Use precisely the momentum-space convention stated
in the erratum, whose forward exponential has a negative sign:
\[
\varphi(\mathbf p)=\frac1{(2\pi\hbar)^{3/2}}\int
\psi(\mathbf r)e^{-i\mathbf r\cdot\mathbf p/\hbar}d^3r.
\]
(a) Differentiation of the exponential gives
\[
\nabla_p e^{-i\mathbf r\cdot\mathbf p/\hbar}
=-\frac{i\mathbf r}{\hbar}e^{-i\mathbf r\cdot\mathbf p/\hbar}.
\]
Move the derivative under the integral and multiply by $i\hbar$;
since $i\hbar(-i\mathbf r/\hbar)=\mathbf r$,
\[
\boxed{i\hbar\nabla_p\varphi
=\frac1{(2\pi\hbar)^{3/2}}\int
\mathbf r\psi(\mathbf r)e^{-i\mathbf r\cdot\mathbf p/\hbar}d^3r}.
\]
(b) Apply another component derivative: differentiating twice
with respect to $p_j$ brings down
$(-ir_j/\hbar)^2=-r_j^2/\hbar^2$. Summing over $j=x,y,z$,
\[
\nabla_p^2 e^{-i\mathbf r\cdot\mathbf p/\hbar}
=-\frac{r_x^2+r_y^2+r_z^2}{\hbar^2}e^{-i\mathbf r\cdot\mathbf p/\hbar}.
\]
Since $(i\hbar\nabla_p)^2=-\hbar^2\nabla_p^2$ and
$-\hbar^2(-r^2/\hbar^2)=r^2$, differentiating under the integral
proves
\[
\boxed{(i\hbar\nabla_p)^2\varphi
=\frac1{(2\pi\hbar)^{3/2}}\int r^2\psi(\mathbf r)
e^{-i\mathbf r\cdot\mathbf p/\hbar}d^3r}.
\]
\erratum{The exponent printed in part (b) omits the factor $i$:
it reads $-\mathbf r\cdot\mathbf p/\hbar$. The Fourier kernel
must be $e^{-i\mathbf r\cdot\mathbf p/\hbar}$ in both parts.}
The same argument works for every polynomial in the coordinate
components. Appropriate integrable moments of $\psi$ justify
ordinary differentiation; otherwise operator domains or a
well-defined distributional interpretation must be specified.

\Needspace{5\baselineskip}
\par\bigskip
\noindent\textbf{20.4.11}\quad The ordinary space wave function
$\psi(\mathbf{r},t)$ satisfies the time-dependent Schr\"odinger
equation,
\[
i\hbar\,\frac{\partial\psi(\mathbf{r},t)}{\partial t}
=-\frac{\hbar^{2}}{2m}\nabla^{2}\psi+V(\mathbf{r})\psi.
\]
Show that the corresponding time-dependent momentum wave function
satisfies the analogous equation
\[
i\hbar\,\frac{\partial\varphi(\mathbf{p},t)}{\partial t}
=\frac{p^{2}}{2m}\,\varphi+V(i\hbar\nabla_{p})\varphi.
\]
\par\emph{Note.}~Assume that $V(\mathbf{r})$ may be expressed by
a Maclaurin series \erratum{whose termwise transformation is
justified on the chosen wave function (a polynomial potential
suffices)}, and use Exercise~20.4.10.
$V(i\hbar\nabla_{p})$ is the same function of the variable
$i\hbar\nabla_{p}$ that $V(\mathbf{r})$ is of the variable
$\mathbf{r}$.
\solution
Transform with the negative-exponent momentum
convention of Exercise~20.4.10; for any function $\chi$ write
\[
\mathcal F_p[\chi](\mathbf p)=\frac1{(2\pi\hbar)^{3/2}}\int\chi(\mathbf r)
e^{-i\mathbf r\cdot\mathbf p/\hbar}d^3r,
\]
so that $\varphi=\mathcal F_p[\psi]$. The time derivative commutes
with the spatial transform, producing $i\hbar\partial_t\varphi$.
For the Laplacian, integrate by parts in $x$; the boundary terms
vanish because $\psi\to0$ at infinity, and
$\partial_xe^{-i\mathbf r\cdot\mathbf p/\hbar}
=-(ip_x/\hbar)e^{-i\mathbf r\cdot\mathbf p/\hbar}$, so
\[
\mathcal F_p\Big[\frac{\partial\psi}{\partial x}\Big]
=-\frac1{(2\pi\hbar)^{3/2}}\int\psi\,
\frac{\partial}{\partial x}e^{-i\mathbf r\cdot\mathbf p/\hbar}d^3r
=\frac{ip_x}\hbar\varphi.
\]
Applying this twice in each coordinate and summing
$(ip_j/\hbar)^2=-p_j^2/\hbar^2$ over $j=x,y,z$ gives
\[
\mathcal F_p[\nabla_r^2\psi]
=-\frac{p^2}{\hbar^2}\varphi,
\]
so the kinetic term becomes
$-(\hbar^2/2m)(-p^2/\hbar^2)\varphi=p^2\varphi/(2m)$.
Write a convergent coordinate expansion with multi-index
$\alpha=(\alpha_x,\alpha_y,\alpha_z)$,
$\mathbf r^\alpha=x^{\alpha_x}y^{\alpha_y}z^{\alpha_z}$ and
$|\alpha|=\alpha_x+\alpha_y+\alpha_z$:
\[
V(\mathbf r)=\sum_\alpha v_\alpha\mathbf r^\alpha.
\]
Exercise~20.4.10(a), applied component by component, gives
$\mathcal F_p[x\chi]=i\hbar\,\partial_{p_x}\mathcal F_p[\chi]$
for any admissible $\chi$; applying it to
$\chi=\psi,x\psi,x^2\psi,\ldots$ and likewise for $y$ and $z$,
\[
\mathcal F_p[\mathbf r^\alpha\psi]
=(i\hbar)^{|\alpha|}\partial_{p_x}^{\alpha_x}
\partial_{p_y}^{\alpha_y}\partial_{p_z}^{\alpha_z}\varphi
\equiv(i\hbar\nabla_p)^\alpha\varphi.
\]
Consequently
\[
\mathcal F_p[V\psi]=\sum_\alpha v_\alpha(i\hbar\nabla_p)^\alpha\varphi
=V(i\hbar\nabla_p)\varphi.
\]
Combining the terms gives
\[
\boxed{i\hbar\frac{\partial\varphi}{\partial t}
=\frac{p^2}{2m}\varphi+V(i\hbar\nabla_p)\varphi}.
\]
\erratum{For a nonpolynomial potential, a local Maclaurin series
alone does not justify interchange with an integral over all
space. Its action on the chosen wave function must converge
and permit that interchange. Polynomial potentials satisfy the
algebraic step without an infinite-series issue.}
When a general potential is treated by its transform instead,
the rigorous replacement is the momentum convolution
\[
\mathcal F_p[V\psi](\mathbf p)
=\frac1{(2\pi\hbar)^{3/2}}\int
\widetilde V(\mathbf p-\mathbf q)\varphi(\mathbf q)d^3q,
\]
where $\widetilde V=\mathcal F_pV$ with the same convention; it
follows by inserting
$\psi(\mathbf r)=(2\pi\hbar)^{-3/2}\int\varphi(\mathbf q)
e^{i\mathbf r\cdot\mathbf q/\hbar}d^3q$ into $\mathcal F_p[V\psi]$
and doing the $\mathbf r$ integral first, which produces
$\widetilde V(\mathbf p-\mathbf q)$. This holds
whenever the convolution, possibly distributional, is defined.
The differential expression is its polynomial or convergent-series
form.

\Needspace{9\baselineskip}
\subsection*{20.5 Signal-Processing Applications}

\Needspace{5\baselineskip}
\par\bigskip
\noindent\textbf{20.5.1}\quad Find the transfer function
$\varphi(\omega)$ for the circuit shown in the left panel of
Fig.~20.10. Is this a high-pass, a low-pass, or a more complicated
filter?
\solution
The left circuit has a series resistor $R$ and a capacitor $C$
across the output terminals. With the $e^{i\omega t}$ phasor
convention used in Section~20.5, $Z_C=1/(i\omega C)$. The
current is $I=V_{\rm in}/(R+Z_C)$, and $V_{\rm out}=IZ_C$.
Therefore, dividing numerator and denominator by $Z_C$ and using
$R/Z_C=i\omega RC$,
\[
\boxed{\varphi(\omega)=\frac{V_{\rm out}}{V_{\rm in}}
=\frac{Z_C}{R+Z_C}=\frac1{1+i\omega RC}}.
\]
For positive $R,C$, since $|1+i\omega RC|^2=1+\omega^2R^2C^2$,
\[
|\varphi|^2=\frac1{1+\omega^2R^2C^2},\qquad
\arg\varphi=-\tan^{-1}(\omega RC)\quad(\omega\ge0).
\]
The gain decreases from $1$ at zero frequency to $0$ at high
frequency: this is a low-pass filter. The power gain is $1/2$
at $\omega_c=1/(RC)$. These expressions assume an unloaded
output, as in the drawing. The sign of $i\omega$ follows the
section's phasor convention; the opposite time convention
conjugates the transfer function.

\Needspace{5\baselineskip}
\par\bigskip
\noindent\textbf{20.5.2}\quad Find the transfer function
$\varphi(\omega)$ for the circuit shown in the right panel of
Fig.~\erratum{20.10}.
\par\emph{Hint.}~The potential difference across an inductor is
given by $L\,dI/dt$.
\solution
\erratum{The intended reference is the right panel of
Fig.~20.10, which contains a series inductor $L$ and an output
resistor $R$. Figure~20.11 has a single, different circuit.}
For a current $I e^{i\omega t}$, the voltage across the
inductor is $L\,dI/dt=i\omega LI$. Kirchhoff's voltage law and
the output voltage are consequently
\[
V_{\rm in}=(R+i\omega L)I,\qquad V_{\rm out}=RI.
\]
Dividing gives
\[
\boxed{\varphi(\omega)=\frac{R}{R+i\omega L}
=\frac1{1+i\omega L/R}}.
\]
Its power gain is $|\varphi|^2=[1+(\omega L/R)^2]^{-1}$, which
decreases from $1$ at $\omega=0$ to $0$ as $\omega\to\infty$ and
equals $1/2$ at $\omega L/R=1$; so it is a
low-pass filter with half-power frequency $R/L$, for positive
$R,L$.

\Needspace{5\baselineskip}
\par\bigskip
\noindent\textbf{20.5.3}\quad Find the transfer function
$\varphi(\omega)$ for the circuit shown in Fig.~20.11. This is a
band-pass filter.
\par\emph{Hint.}~Assume the currents in the various parts of the
circuit to have the values shown in the figure.
\solution
Use the currents specified in Fig.~20.11: $I_1$ flows
through $R_1$, $I_2$ through the series combination $R_2,C_2$,
and $I_1+I_2$ through $C_1$. Put $p=i\omega$. The two node
voltages satisfy
\[
V_1=R_1I_1=\left(R_2+\frac1{pC_2}\right)I_2,
\quad V_{\rm out}=\frac{I_2}{pC_2},
\quad V_{\rm in}=V_1+\frac{I_1+I_2}{pC_1}.
\]
The middle relation gives $I_2=pC_2V_{\rm out}$, and the
first then gives
\[
V_1=(1+pR_2C_2)V_{\rm out},\qquad
I_1=\frac{1+pR_2C_2}{R_1}V_{\rm out}.
\]
Hence
$I_1+I_2=[(1+pR_2C_2)/R_1+pC_2]V_{\rm out}$. Substitute $V_1$
and $I_1+I_2$ into the input relation and multiply
by $pR_1C_1$:
\[
pR_1C_1V_{\rm in}
=\big[pR_1C_1(1+pR_2C_2)+1+pR_2C_2+pR_1C_2\big]V_{\rm out}.
\]
The bracket expands to
$1+p(R_1C_1+R_1C_2+R_2C_2)+p^2R_1R_2C_1C_2$, and
$p^2=-\omega^2$. Thus
\[
\boxed{\varphi(\omega)=
\frac{i\omega R_1C_1}
{1+i\omega(R_1C_1+R_1C_2+R_2C_2)-\omega^2R_1R_2C_1C_2}}.
\]
To check the filter type, write $A=R_1C_1$,
$B=R_1R_2C_1C_2$, and $D=R_1C_1+R_1C_2+R_2C_2$, all
positive. The denominator of $\varphi$ is $(1-B\omega^2)+iD\omega$,
so
\[
|\varphi|^2=\frac{A^2\omega^2}{(1-B\omega^2)^2+D^2\omega^2}
=\frac{A^2}{\omega^{-2}+B^2\omega^2+D^2-2B},
\]
where the second form divides numerator and denominator by
$\omega^2$ and expands $(1-B\omega^2)^2=1-2B\omega^2+B^2\omega^4$.
It vanishes as $\omega\to0$ and $\omega\to\infty$. Its
maximum occurs where the denominator is smallest. By the
inequality of arithmetic and geometric means,
$\omega^{-2}+B^2\omega^2\ge2\sqrt{\omega^{-2}B^2\omega^2}=2B$,
with equality exactly when $\omega^{-2}=B^2\omega^2$ (equivalently
$-2\omega^{-3}+2B^2\omega=0$), namely
$\omega=1/\sqrt B$. There the denominator is
$2B+D^2-2B=D^2$, so the maximum gain is $|\varphi|=A/D$.
It therefore passes an intermediate band of frequencies.

\Needspace{5\baselineskip}
\par\bigskip
\noindent\textbf{20.5.4}\quad The circuit elements for
Exercise~20.5.3 correspond to the successive transfer functions
shown in Fig.~20.12. Explain why the transfer function for this
exercise \erratum{is uniformly approximated by the product of the
individual unloaded transfer functions when $R_2\gg R_1$}.
\erratum{Explain the loading correction. The condition
$R_2\gg R_1$ is sufficient for a uniformly accurate product
approximation but, contrary to the textbook's word ``only'', it is
not necessary.}
\solution
Taken separately with unloaded outputs, the high-pass
and low-pass stages have transfer functions
\[
\varphi_1=\frac{pR_1C_1}{1+pR_1C_1},\qquad
\varphi_2=\frac1{1+pR_2C_2},\qquad p=i\omega.
\]
Their product is
\[
\varphi_1\varphi_2=\frac{pR_1C_1}
{1+p(R_1C_1+R_2C_2)+p^2R_1R_2C_1C_2}.
\]
The actual denominator found in Exercise~20.5.3 contains the
additional term $pR_1C_2$. This appears because the second
stage draws current from the first: the current through $C_1$
is $I_1+I_2$, not $I_1$. Quantitatively,
\[
\frac{I_2}{I_1}=\frac{R_1}{R_2+1/(i\omega C_2)},\qquad
\left|\frac{I_2}{I_1}\right|
=\frac{R_1}{\sqrt{R_2^2+1/(\omega^2C_2^2)}}\le\frac{R_1}{R_2}.
\]
Hence $R_2\gg R_1$ makes the load current small at every real
frequency. More explicitly, let
$D_0=(1+pR_1C_1)(1+pR_2C_2)=1+p(R_1C_1+R_2C_2)+p^2R_1R_2C_1C_2$,
so that $\varphi_1\varphi_2=pR_1C_1/D_0$. The denominator of
$\varphi$ is $D_0+pR_1C_2=D_0(1+\eta)$, hence
\[
\varphi=\frac{\varphi_1\varphi_2}{1+\eta},\qquad
\eta=\frac{pR_1C_2}{D_0},\qquad |\eta|\le R_1/R_2.
\]
Indeed,
\[
|\eta|=\frac{\omega R_1C_2}
{\sqrt{1+\omega^2R_1^2C_1^2}\sqrt{1+\omega^2R_2^2C_2^2}}
\le\frac{R_1}{R_2}\quad(\omega\ge0),
\]
since the first square root is at least $1$ and the second
at least $\omega R_2C_2$ (the case $\omega=0$ has $\eta=0$).
For $\delta=R_1/R_2<1$, the relative discrepancy, wherever
the product is nonzero, therefore satisfies
\[
\left|\frac{\varphi}{\varphi_1\varphi_2}-1\right|
=\left|\frac1{1+\eta}-1\right|
=\frac{|\eta|}{|1+\eta|}\le\frac{|\eta|}{1-|\eta|}
\le\frac{\delta}{1-\delta},
\]
by the reverse triangle inequality $|1+\eta|\ge1-|\eta|$ and
because $x/(1-x)$ increases on $0\le x<1$.
\erratum{The quoted limit is a sufficient condition for a
uniformly accurate product approximation, not a necessary one.
At sufficiently low and high frequencies the loading is small in
any case; an ideal buffer between stages makes the product
exact. For finite positive elements in the unbuffered circuit,
the product is not an identity in frequency.}
The exact size of the correction shows that $R_2\gg R_1$ is not
necessary. With $a=R_1C_1$ and
$b=R_2C_2$, dividing the numerator and denominator of $|\eta|^2$
by $\omega^2$ and expanding $(1+\omega^2a^2)(1+\omega^2b^2)$ gives
\[
|\eta|^2=\frac{R_1^2C_2^2}{\omega^{-2}+a^2+b^2+a^2b^2\omega^2}
\le\frac{R_1^2C_2^2}{(a+b)^2},
\]
since $\omega^{-2}+a^2b^2\omega^2\ge2ab$ (the inequality of
Exercise~20.5.3, where $B=ab$), with equality at the passband
center $\omega=1/\sqrt{ab}$. Thus
$|\eta|\to0$ as $\omega\to0$ or $\omega\to\infty$, and
\[
\max_\omega|\eta|=\frac{R_1C_2}{R_1C_1+R_2C_2}
<\frac{R_1C_2}{R_2C_2}=\frac{R_1}{R_2}.
\]
The product is therefore uniformly accurate whenever
$R_1C_2\ll R_1C_1+R_2C_2$, for example also when $C_2\ll C_1$
with comparable resistances.

\Needspace{20\baselineskip}
\subsection*{20.6 Discrete Fourier Transform}

\Needspace{5\baselineskip}
\par\bigskip
\noindent\textbf{20.6.1}\quad Derive the trigonometric forms of
discrete orthogonality corresponding to Eq.~(20.120):
\begin{align*}
\sum_{k=0}^{N-1}\cos(2\pi pk/N)\sin(2\pi qk/N)&=0\\
\sum_{k=0}^{N-1}\cos(2\pi pk/N)\cos(2\pi qk/N)&=
\erratum{\begin{cases}
N, & p=q=(0\text{ or }N/2);\\[3pt]
N/2, & (p+q=0\text{ or }N)\text{ or }p=q\text{ but not both};\\[3pt]
0, & \text{otherwise}.
\end{cases}}\\
\sum_{k=0}^{N-1}\sin(2\pi pk/N)\sin(2\pi qk/N)&=
\erratum{\begin{cases}
N/2, & p=q\text{ and }p+q\ne(0\text{ or }N);\\[3pt]
-N/2, & p\ne q\text{ and }p+q=N;\\[3pt]
0, & \text{otherwise}.
\end{cases}}
\end{align*}
\par\emph{Note.}~If $N$ is odd, $p$ and $q$ will never have the
value $N/2$.
\par\emph{Hint.}~Consider the use of trigonometric identities
such as
\[
\sin A\cos B=\frac{1}{2}\big[\sin(A+B)+\sin(A-B)\big].
\]
\par\erratum{Take integer indices $0\le p,q<N$ and a positive integer $N$ in these piecewise formulas.}
\solution
For integers $m$ define the discrete selector
\[
\Delta_N(m)=\begin{cases}1,&m\equiv0\pmod N,\\0,&\text{otherwise}.
\end{cases}
\]
With $r=e^{2\pi im/N}$ the finite geometric series gives
\[
\sum_{k=0}^{N-1}e^{2\pi imk/N}=\sum_{k=0}^{N-1}r^k=N\Delta_N(m).
\]
Indeed, if $m\equiv0\pmod N$, every term is $r^k=1$; otherwise $r\ne1$
and the sum is $(1-r^N)/(1-r)=0$, because
$r^N=e^{2\pi im}=1$. Its imaginary part
$\sum_k\sin(2\pi mk/N)$ is therefore zero and its real part
$\sum_k\cos(2\pi mk/N)$ is $N\Delta_N(m)$. With
$\theta_k=2\pi k/N$, the product identities
$\cos A\sin B=\frac12[\sin(B+A)+\sin(B-A)]$ and
$\cos A\cos B,\ \sin A\sin B=\frac12[\cos(A-B)\pm\cos(A+B)]$
then give
\begin{align*}
\sum_k\cos(p\theta_k)\sin(q\theta_k)
&=\frac12\sum_k\{\sin[(q+p)\theta_k]+\sin[(q-p)\theta_k]\}=0,\\
\sum_k\cos(p\theta_k)\cos(q\theta_k)
&=\frac12\sum_k\{\cos[(p-q)\theta_k]+\cos[(p+q)\theta_k]\}\\
&=\frac N2\{\Delta_N(p-q)+\Delta_N(p+q)\},\\
\sum_k\sin(p\theta_k)\sin(q\theta_k)
&=\frac12\sum_k\{\cos[(p-q)\theta_k]-\cos[(p+q)\theta_k]\}\\
&=\frac N2\{\Delta_N(p-q)-\Delta_N(p+q)\}.
\end{align*}
These formulas hold for all integer $p,q$ and positive integer
$N$. \erratum{The displayed piecewise formulas in the question
use the representative range $0\le p,q<N$. Without that range,
congruences modulo $N$ must replace the equalities.}
In this range $|p-q|<N$, so $\Delta_N(p-q)=1$ exactly when $p=q$.
Also $0\le p+q\le2N-2<2N$, so $\Delta_N(p+q)=1$ exactly when
$p+q=0$ or $N$.
Both selectors equal $1$ for $p=q=0$, and also for
$p=q=N/2$ when $N$ is even. The cosine sum is then $N$,
whereas the sine sum is zero. When exactly one selector is
$1$, the cosine sum is $N/2$; the sine sum is $N/2$ for
$p=q$ and $-N/2$ for $p+q=N$ with $p\ne q$.
When neither selector is $1$, both sums vanish. These cases
prove the corrected formulas, including $N=1$ and $N=2$.
\erratum{The red cases in the statement include the pairs with
$p+q=N$, $p\ne q$, for which $\Delta_N(p+q)=1$: for example,
$N=4$, $p=1$, $q=3$ gives the cosine sum $2=N/2$ and the sine sum
$-2=-N/2$, although $p\ne q$. The middle cosine case is written
with $p+q=0$ or $N$, so that $p=q=0$, where both of its
conditions hold, falls under the first case only; with $p+q=N$
alone the first two cases would overlap at $p=q=0$.}
\erratum{The final product-to-sum line in the Instructor's
Manual's proof on p.~344 has a cosine--cosine product on its left. It must
be a sine--sine product to agree with the difference of cosines
on its right. The selector derivation above gives that identity
directly.}

\Needspace{5\baselineskip}
\par\bigskip
\noindent\textbf{20.6.2}\quad Show in detail how to go from
\[
F_{p}=\frac{1}{N^{1/2}}\sum_{k=0}^{N-1}
f_{k}e^{2\pi ipk\erratum{/N}}
\qquad\text{to}\qquad
f_{k}=\frac{1}{N^{1/2}}\sum_{p=0}^{N-1}
F_{p}e^{-2\pi ipk\erratum{/N}}.
\]
\solution
Insert the defining expression for $F_p$ into the proposed
inverse, using a different dummy index $j$:
\begin{align*}
\frac1{\sqrt N}\sum_{p=0}^{N-1}F_pe^{-2\pi ipk/N}
&=\frac1N\sum_{p=0}^{N-1}\sum_{j=0}^{N-1}
f_je^{2\pi ipj/N}e^{-2\pi ipk/N}\\
&=\frac1N\sum_{j=0}^{N-1}f_j
\sum_{p=0}^{N-1}e^{2\pi ip(j-k)/N}.
\end{align*}
The inner sum is $N\Delta_N(j-k)$ by the finite geometric
series of Exercise~20.6.1 (with $m=j-k$). For $0\le j,k<N$ we
have $|j-k|<N$, so the only surviving index is $j=k$, and its
factor $N$ cancels the $1/N$.
Thus the expression is $f_k$, which proves
\[
\boxed{f_k=\frac1{\sqrt N}\sum_{p=0}^{N-1}F_pe^{-2\pi ipk/N}}.
\]
All sums are finite, so their interchange requires no
convergence assumption. \erratum{The factors $1/N$ in both
exponents are essential: without them every exponential is $1$
for integer indices, and the transform cannot distinguish the
input entries.}

\Needspace{5\baselineskip}
\par\bigskip
\noindent\textbf{20.6.3}\quad The $N$-membered point sets $f_{k}$
and $F_{p}$ are discrete Fourier transforms of each other. Derive
the following symmetry relations:
\begin{enumerate}
\item[(a)] If $f_{k}$ is real, $F_{p}$ is Hermitian symmetric;
that is, $F_{p}=F_{N-p}^{*}$.
\item[(b)] If $f_{k}$ is pure imaginary, then
$F_{p}=-F_{N-p}^{*}$.
\end{enumerate}
\par\emph{Note.}~The symmetry of part (a) is an illustration of
aliasing. If $F_{p}$ describes an amplitude at a frequency
proportional to $p$, we necessarily predict an equal amplitude at
the frequency proportional to $N-p$.
\par\erratum{Indices are understood modulo $N$, in particular $F_N=F_0$. Equal paired amplitudes in the note mean equal magnitudes, with conjugate phases.}
\solution
Regard the frequency index modulo $N$, so $F_N=F_0$.
Complex conjugation gives
\begin{align*}
F_{N-p}^{*}
&=\frac1{\sqrt N}\sum_{k=0}^{N-1}f_k^*
 e^{-2\pi i(N-p)k/N}\\
&=\frac1{\sqrt N}\sum_{k=0}^{N-1}f_k^*
 e^{-2\pi ik}e^{2\pi ipk/N}
=\frac1{\sqrt N}\sum_{k=0}^{N-1}f_k^*e^{2\pi ipk/N},
\end{align*}
since $e^{-2\pi ik}=1$ for integer $k$.
If $f_k$ is real, $f_k^*=f_k$, so $F_{N-p}^*=F_p$.
If $f_k$ is purely imaginary, $f_k^*=-f_k$, so
$F_{N-p}^*=-F_p$. These are precisely (a) and (b).
For a real input, the zero-frequency coefficient and, for even
$N$, the coefficient at $p=N/2$ are real because they are their
own conjugates. For a purely imaginary input those same
coefficients are purely imaginary. In general the paired
coefficients have equal magnitudes; their phases are related
by conjugation (and by an extra minus sign in part (b)). The
index $N-p$ represents the same sampled frequency as $-p$.

\Needspace{9\baselineskip}
\subsection*{20.7 Laplace Transforms}

\Needspace{5\baselineskip}
\par\bigskip
\noindent\textbf{20.7.1}\quad Prove that
\[
\lim_{s\to\infty}s f(s)=\lim_{t\to+0}F(t).
\]
\par\emph{Hint.}~Assume that $F(t)$ can be expressed as
$F(t)=\displaystyle\sum_{n=0}^{\infty}a_{n}t^{n}$.
\par\erratum{Let $s\to+\infty$ along the real axis. Assume a finite $F(0+)$ and an exponential bound $|F(t)|\le Me^{ct}$; these hypotheses suffice without an unjustified termwise integration.}
\solution
\erratum{Take $s\to+\infty$ along the positive real axis.
A sufficient hypothesis is that $F$ has the finite right limit
$F(0+)$, is locally integrable, and has an exponential bound
$|F(t)|\le M e^{ct}$ for $t\ge0$.}
In the defining integral put $u=st$, so $dt=du/s$ and, for
$s>0$, the limits stay $0$ and $\infty$:
\[
sf(s)=s\int_0^\infty e^{-st}F(t)dt
=\int_0^\infty e^{-u}F(u/s)du.
\]
For each fixed $u>0$, $u/s\to0^+$ and so $F(u/s)\to F(0+)$.
For sufficiently large $s$, $c/s\le1/2$, and
$|e^{-u}F(u/s)|\le Me^{-u}e^{cu/s}\le M e^{-u/2}$, which is
integrable (if $c\le0$, the bound $Me^{-u}$ holds for every
$s>0$).
Dominated convergence therefore gives
\[
\lim_{s\to+\infty}sf(s)
=F(0+)\int_0^\infty e^{-u}du=\boxed{F(0+)}.
\]
The hint also describes the coefficients: if a Taylor expansion
can be integrated termwise, then, with
$\int_0^\infty e^{-st}t^ndt=\Gamma(n+1)/s^{n+1}$ (substitute
$u=st$ as above; see also Exercise~20.7.7),
\[
f(s)=\sum_{n\ge0}\frac{a_n\Gamma(n+1)}{s^{n+1}},\qquad
sf(s)=a_0+\frac{a_1}{s}+\frac{2a_2}{s^2}+\cdots.
\]
Even when this infinite series is only asymptotic, the direct
limit proof above remains valid and identifies $a_0=F(0+)$.

\Needspace{5\baselineskip}
\par\bigskip
\noindent\textbf{20.7.2}\quad Show that
\[
\frac1\pi\lim_{s\to0}L\{\cos(xt)\}=\delta(x).
\]
\par\erratum{This formula leaves the transform variable, the
approach to zero, and the meaning of convergence unspecified; the
published errata therefore say to disregard it. Prove instead the
following precise replacement:}
\[
\erratum{\frac1\pi\lim_{s\downarrow0}L_t\{\cos xt\}(s)=\delta(x)
\quad\text{as distributions in }x.}
\]
\solution
\erratum{The published errata withdraw the statement as
ill-defined. It has a useful, precise replacement: let $s$
approach zero through positive real values, take the Laplace
transform in $t$, and interpret convergence as convergence of
distributions in $x$.}
For $s>0$ and real $x$, write $\cos xt=\Re e^{ixt}$; direct
integration of $e^{-(s-ix)t}$ over $0<t<\infty$ gives $1/(s-ix)$,
so
\[
L_t\{\cos xt\}(s)=\Re\frac1{s-ix}=\Re\frac{s+ix}{s^2+x^2}
=\frac{s}{s^2+x^2}.
\]
The resulting functions do not converge to an ordinary
function representing a delta: they approach zero at every
$x\ne0$ and diverge at $x=0$. To prove the correct weak limit,
take a smooth test function $\phi$ with compact support and
substitute $x=su$, so $dx=s\,du$ and $s^2+x^2=s^2(1+u^2)$:
\[
\int_{-\infty}^{\infty}\frac{s\phi(x)}{\pi(s^2+x^2)}dx
=\frac1\pi\int_{-\infty}^{\infty}\frac{\phi(su)}{1+u^2}du.
\]
The integrand is bounded by the integrable
$\|\phi\|_\infty/(1+u^2)$, and $\phi(su)\to\phi(0)$ for each $u$
by continuity. Consequently, by dominated convergence, the limit is
\[
\frac{\phi(0)}\pi[\tan^{-1}u]_{-\infty}^{\infty}=\phi(0).
\]
Thus the repaired assertion is
\[
\boxed{\lim_{s\downarrow0}\frac1\pi L_t\{\cos xt\}(s)
=\delta(x)\quad\text{in the distributional sense}.}
\]
The one-sided real approach is part of the assertion; an
unspecified complex limit is not the prescription proved here.

\Needspace{5\baselineskip}
\par\bigskip
\noindent\textbf{20.7.3}\quad Verify that
\[
L\left\{\frac{\cos at-\cos bt}{b^{2}-a^{2}}\right\}
=\frac{s}{(s^{2}+a^{2})(s^{2}+b^{2})},\qquad a^{2}\ne b^{2}.
\]
\solution
For real $a,b$ and $\Re s>0$, write
$\cos at=\frac12(e^{iat}+e^{-iat})$ and use
$L\{e^{\pm iat}\}=\int_0^\infty e^{-(s\mp ia)t}dt=1/(s\mp ia)$:
\[
L\{\cos at\}=\frac12\left(\frac1{s-ia}+\frac1{s+ia}\right)
=\frac12\,\frac{(s+ia)+(s-ia)}{s^2+a^2}
=\frac{s}{s^2+a^2}.
\]
Linearity and a common denominator now yield
\begin{align*}
L\left\{\frac{\cos at-\cos bt}{b^2-a^2}\right\}
&=\frac{s}{b^2-a^2}\left(\frac1{s^2+a^2}-\frac1{s^2+b^2}\right)\\
&=\frac{s}{b^2-a^2}
\frac{(s^2+b^2)-(s^2+a^2)}{(s^2+a^2)(s^2+b^2)}\\
&=\boxed{\frac{s}{(s^2+a^2)(s^2+b^2)}}.
\end{align*}
The stipulated $a^2\ne b^2$ permits the cancellation. For
complex frequencies the same algebra applies in a common
half-plane of convergence.

\Needspace{5\baselineskip}
\par\bigskip
\noindent\textbf{20.7.4}\quad Using partial fraction expansions,
show that
\begin{enumerate}
\item[(a)] $L^{-1}\left\{\dfrac{1}{(s+a)(s+b)}\right\}
=\dfrac{e^{-at}-e^{-bt}}{b-a}$,\quad $a\ne b$.
\item[(b)] $L^{-1}\left\{\dfrac{s}{(s+a)(s+b)}\right\}
=\dfrac{ae^{-at}-be^{-bt}}{a-b}$,\quad $a\ne b$.
\end{enumerate}
\solution
(a) Set $1/[(s+a)(s+b)]=A/(s+a)+B/(s+b)$.
After multiplication by the denominator,
$1=A(s+b)+B(s+a)$. Setting $s=-a$ gives $1=A(b-a)$ and
setting $s=-b$ gives $1=B(a-b)$, so
$A=1/(b-a)$ and $B=1/(a-b)=-1/(b-a)$. Since
$L\{e^{-at}\}=1/(s+a)$,
\[
\boxed{L^{-1}\left\{\frac1{(s+a)(s+b)}\right\}
=\frac{e^{-at}-e^{-bt}}{b-a}}.
\]
(b) Now $s=A(s+b)+B(s+a)$, so the same substitutions give
$-a=A(b-a)$ and $-b=B(a-b)$, i.e.,
$A=a/(a-b)$, $B=-b/(a-b)$. Hence
\[
\boxed{L^{-1}\left\{\frac{s}{(s+a)(s+b)}\right\}
=\frac{ae^{-at}-be^{-bt}}{a-b}}.
\]
The defining transforms converge for
$\Re s>\max(-\Re a,-\Re b)$. Although excluded in the
question, the equal-parameter limits are regular. By
l'H\^opital's rule in $b$ (differentiate numerator and
denominator with respect to $b$, then let $b\to a$),
\begin{align*}
\lim_{b\to a}\frac{e^{-at}-e^{-bt}}{b-a}
&=\lim_{b\to a}\frac{te^{-bt}}1=te^{-at},\\
\lim_{b\to a}\frac{ae^{-at}-be^{-bt}}{a-b}
&=\lim_{b\to a}\frac{-e^{-bt}+bte^{-bt}}{-1}=(1-at)e^{-at};
\end{align*}
these are the inverse transforms of $1/(s+a)^2$ and
$s/(s+a)^2$.

\Needspace{5\baselineskip}
\par\bigskip
\noindent\textbf{20.7.5}\quad Using partial fraction expansions,
show that for $a^{2}\ne b^{2}$,
\begin{enumerate}
\item[(a)] $L^{-1}\left\{\dfrac{1}{(s^{2}+a^{2})(s^{2}+b^{2})}
\right\}=-\dfrac{1}{a^{2}-b^{2}}\left(\dfrac{\sin at}{a}
-\dfrac{\sin bt}{b}\right)$.
\item[(b)] $L^{-1}\left\{\dfrac{s^{2}}{(s^{2}+a^{2})(s^{2}+b^{2})}
\right\}=\dfrac{1}{a^{2}-b^{2}}\{a\sin at-b\sin bt\}$.
\end{enumerate}
\par\erratum{For real $a,b$, use $\sin(at)/a\to t$ if $a=0$, and the corresponding limit if $b=0$.}
\solution
For real parameters use
$L\{\sin at/a\}=1/(s^2+a^2)$ in $\Re s>0$.
Write the transforms as $A/(s^2+a^2)+B/(s^2+b^2)$ and
multiply by $(s^2+a^2)(s^2+b^2)$.
(a) Equating coefficients of $s^2$ and $s^0$ in
\[
1=A(s^2+b^2)+B(s^2+a^2)
\]
gives $A+B=0$, $Ab^2+Ba^2=1$. Substituting $B=-A$ into the
second equation gives $A(b^2-a^2)=1$. Thus
$A=1/(b^2-a^2)$ and $B=-A$, and inversion gives
\[
\boxed{\frac1{b^2-a^2}\left(\frac{\sin at}{a}-\frac{\sin bt}{b}\right)}
=-\frac1{a^2-b^2}\left(\frac{\sin at}{a}-\frac{\sin bt}{b}\right),
\]
the printed form.
(b) Equating coefficients in
$s^2=A(s^2+b^2)+B(s^2+a^2)$ gives
$A+B=1$, $Ab^2+Ba^2=0$. Substituting $B=1-A$ into the second
equation gives $A(b^2-a^2)+a^2=0$, so
$A=a^2/(a^2-b^2)$ and $B=1-A=-b^2/(a^2-b^2)$. Inverting
$A/(s^2+a^2)+B/(s^2+b^2)$ gives
$A\sin(at)/a+B\sin(bt)/b$, in which $a^2\sin(at)/a=a\sin at$
and $b^2\sin(bt)/b=b\sin bt$. Therefore
\[
\boxed{L^{-1}\left\{\frac{s^2}{(s^2+a^2)(s^2+b^2)}\right\}
=\frac{a\sin at-b\sin bt}{a^2-b^2}}.
\]
\erratum{If either parameter is zero, the quotient
$\sin(at)/a$ in part (a) means its continuous limit $t$.
The condition $a^2\ne b^2$ alone does not exclude that case.}

\Needspace{5\baselineskip}
\par\bigskip
\noindent\textbf{20.7.6}\quad Show that
\begin{enumerate}
\item[(a)] $\displaystyle\int_{0}^{\infty}\frac{\cos s}{s^{\nu}}\,ds
=\frac{\pi}{2\erratum{\Gamma(\nu)}\cos(\nu\pi/2)}$,
\qquad$0<\nu<1$.
\item[(b)] $\displaystyle\int_{0}^{\infty}\frac{\sin s}{s^{\nu}}\,ds
=\frac{\pi}{2\erratum{\Gamma(\nu)}\sin(\nu\pi/2)}$,
\qquad$0<\nu<2$.
\end{enumerate}
Why is $\nu$ restricted to $(0,1)$ for (a), to $(0,2)$ for (b)?
These integrals may be interpreted as Fourier transforms of
$s^{-\nu}$ and as Mellin transforms of $\sin s$ and $\cos s$.
\par\emph{Hint.}~Replace $s^{-\nu}$ by a Laplace transform
integral: $L\big\{t^{\nu-1}/\Gamma(\nu)\big\}$. Then integrate
with respect to $s$. The resulting integral can be treated as a
beta function (Section~13.3).
\solution
To avoid confusing the transform parameter with the
integration variable, write the latter as $x$. For $\nu>0$,
substitution $v=ux$ into the gamma integral gives
\[
x^{-\nu}=\frac1{\Gamma(\nu)}\int_0^\infty
u^{\nu-1}e^{-ux}du.
\]
Insert a convergence factor $e^{-\epsilon x}$, $\epsilon>0$,
before changing the order of integration: let
\[
C_\epsilon=\int_0^\infty e^{-\epsilon x}x^{-\nu}\cos x\,dx,
\qquad
S_\epsilon=\int_0^\infty e^{-\epsilon x}x^{-\nu}\sin x\,dx,
\]
and substitute the representation of $x^{-\nu}$. In the ranges
stated in the question the resulting double integral is absolutely
integrable: after integrating its absolute value in $u$, it is
$\int_0^\infty e^{-\epsilon x}x^{-\nu}|\cos x|dx$ or its sine
counterpart, finite because $|\cos x|\le1$ and $|\sin x|\le x$
near $x=0$. Hence the order may be interchanged, and the inner
$x$ integrals are those of Exercise~20.2.3 with $a=u+\epsilon$,
$\omega=1$:
$\int_0^\infty e^{-(u+\epsilon)x}\cos x\,dx=(u+\epsilon)/[(u+\epsilon)^2+1]$
and
$\int_0^\infty e^{-(u+\epsilon)x}\sin x\,dx=1/[(u+\epsilon)^2+1]$.
Thus
\begin{align*}
C_\epsilon
&=\frac1{\Gamma(\nu)}\int_0^\infty
\frac{u^{\nu-1}(u+\epsilon)}{(u+\epsilon)^2+1}du,\\
S_\epsilon
&=\frac1{\Gamma(\nu)}\int_0^\infty
\frac{u^{\nu-1}}{(u+\epsilon)^2+1}du.
\end{align*}
For $0<\epsilon\le1$, on $0<u<1$ both integrands are bounded
by $2u^{\nu-1}$, since $(u+\epsilon)/[(u+\epsilon)^2+1]\le
u+\epsilon<2$ and $1/[(u+\epsilon)^2+1]\le1$. For $u\ge1$,
since $(u+\epsilon)/[(u+\epsilon)^2+1]<1/(u+\epsilon)\le1/u$
and $1/[(u+\epsilon)^2+1]<1/u^2$, the bounds are
$u^{\nu-2}$ and $u^{\nu-3}$, respectively. These dominating
functions are integrable because $\nu>0$ (at $u=0$) and
$\nu<1$, respectively $\nu<2$ (at $u=\infty$). Dominated
convergence as $\epsilon\downarrow0$ gives
\[
C_0=\frac1{\Gamma(\nu)}\int_0^\infty\frac{u^\nu}{1+u^2}du,
\qquad
S_0=\frac1{\Gamma(\nu)}\int_0^\infty\frac{u^{\nu-1}}{1+u^2}du.
\]
For $0<q<2$, put $v=u^2$, so $u^{q-1}du=\frac12v^{q/2-1}dv$,
and use the beta integral
$\int_0^\infty v^{\alpha-1}(1+v)^{-\alpha-\beta}dv
=B(\alpha,\beta)=\Gamma(\alpha)\Gamma(\beta)/\Gamma(\alpha+\beta)$
(Section~13.3) with $\alpha=q/2$, $\beta=1-q/2$, followed by the
reflection formula $\Gamma(z)\Gamma(1-z)=\pi/\sin\pi z$:
\[
\int_0^\infty\frac{u^{q-1}}{1+u^2}du
=\frac12\int_0^\infty\frac{v^{q/2-1}}{1+v}dv
=\frac12\Gamma(q/2)\Gamma(1-q/2)
=\frac\pi{2\sin(\pi q/2)}.
\]
Take $q=\nu+1$ for the cosine, with
$\sin[\pi(\nu+1)/2]=\cos(\pi\nu/2)$, and $q=\nu$ for the sine;
the condition $0<q<2$ is then $-1<\nu<1$, respectively
$0<\nu<2$:
\[
\boxed{C_0=\frac\pi{2\Gamma(\nu)\cos(\pi\nu/2)}},\quad0<\nu<1;
\qquad
\boxed{S_0=\frac\pi{2\Gamma(\nu)\sin(\pi\nu/2)}},\quad0<\nu<2.
\]
At the origin the cosine integrand behaves as $x^{-\nu}$,
requiring $\nu<1$, whereas the sine integrand behaves as
$x^{1-\nu}$, requiring $\nu<2$. At infinity Dirichlet's test
applies for $\nu>0$; for $\nu\le0$ the ordinary oscillatory
improper integral fails to converge. Abel's theorem for
integrals (Exercise~20.2.2(c), with $h(x)=x^{-\nu}\cos x$ or
$x^{-\nu}\sin x$) therefore identifies the regulated limits
$C_0=\lim_{\epsilon\downarrow0}C_\epsilon$ and
$S_0=\lim_{\epsilon\downarrow0}S_\epsilon$ above with the
ordinary improper integrals in precisely the stated ranges.
The gamma notation in the corrected statement accommodates
noninteger $\nu$.

\Needspace{5\baselineskip}
\par\bigskip
\noindent\textbf{20.7.7}\quad A function $F(t)$ can be expanded in
a Maclaurin series,
\[
F(t)=\sum_{n=0}^{\infty}a_{n}t^{n}.
\]
\erratum{If this series represents $F$ on the half-line and
termwise Laplace integration is justified, then}
\[
L\{F(t)\}=\int_{0}^{\infty}e^{-st}\sum_{n=0}^{\infty}a_{n}t^{n}\,dt
=\sum_{n=0}^{\infty}a_{n}\int_{0}^{\infty}e^{-st}t^{n}\,dt.
\]
Show that $f(s)$, the Laplace transform of $F(t)$, contains no
powers of $s$ greater than $s^{-1}$. Check your result by
calculating $L\{\delta(t)\}$, and comment on this fiasco.
\par\erratum{For the impulse comparison use the causal convention in which a delta at $t=0$ has full unit mass.}
\solution
Whenever termwise integration is justified, $u=st$ (so
$dt=du/s$ and $t^n=u^n/s^n$) gives
\[
\int_0^\infty e^{-st}t^n dt
=\frac1{s^{n+1}}\int_0^\infty e^{-u}u^n du
=\frac{n!}{s^{n+1}},\qquad s>0,
\]
since $\int_0^\infty e^{-u}u^ndu=\Gamma(n+1)=n!$.
Thus
\[
\boxed{f(s)=\sum_{n=0}^\infty\frac{a_n n!}{s^{n+1}}},
\]
which contains only $s^{-1},s^{-2},\ldots$.
\erratum{Existence of a local Maclaurin series does not by
itself justify the displayed interchange over the entire
half-line. For example, $F(t)=1/(1+t)$ has a Maclaurin series
near zero and a Laplace transform for $s>0$, but the resulting
formal series $\sum_{n\ge0}(-1)^n n!/s^{n+1}$ diverges for
every finite $s$. It is an asymptotic expansion at infinity,
not a convergent identity.}
(The Maclaurin series $\sum(-1)^nt^n$ converges only for
$|t|<1$, and in the formal series the ratio of consecutive terms
has modulus $(n+1)/s\to\infty$, so its terms do not even tend to
zero.)
With the causal Laplace convention in which a delta impulse
at $0$ has unit mass, the distributional pairing gives
\[
L\{\delta(t)\}(s)=\langle\delta,e^{-st}\rangle=e^{-s\cdot0}=1.
\]
This $1=s^0$ is a higher power than $s^{-1}$, but it is no
contradiction: a delta distribution is
not an ordinary function with a Maclaurin series. Its
transform lies outside the hypotheses of the argument. The
same distinction permits $L\{\delta^{(n)}\}=s^n$ for causal
distributions, since
$\langle\delta^{(n)},e^{-st}\rangle
=(-1)^n\,d^ne^{-st}/dt^n\big|_{t=0}=(-1)^n(-s)^n=s^n$.
\erratum{The convention at the lower endpoint
must be specified: integrating a symmetric regularization
over $[0,\infty)$ assigns half the mass, whereas the causal
$0^-$ convention used here assigns the full mass.}

\Needspace{5\baselineskip}
\par\bigskip
\noindent\textbf{20.7.8}\quad Show that the Laplace transform of
the confluent hypergeometric function $M(a, c\erratum{,}\, x)$ is
\[
L\big\{M(a, c\erratum{,}\, x)\big\}
=\frac{1}{s}\,{}_{2}F_{1}\left(a, 1; c; \frac{1}{s}\right).
\]
\par\erratum{For the ordinary series definitions take $c\notin\{0,-1,-2,\ldots\}$ and $\Re s>1$; exceptional terminating cases may admit larger domains.}
\solution
Use the defining confluent hypergeometric series
\[
M(a,c,t)=\sum_{n=0}^\infty\frac{(a)_n}{(c)_n}\frac{t^n}{n!}.
\]
Each term is integrated with
$\int_0^\infty e^{-st}t^ndt=n!/s^{n+1}$, $\Re s>0$ (Exercise~20.7.7
for real $s$; both sides are analytic in $s$ for $\Re s>0$, so
the identity extends to that half-plane).
For $\Re s>1$ the series may be integrated absolutely:
since $|e^{-st}|=e^{-t\Re s}$, the sum of the integrals of the
absolute terms is
\[
\sum_{n=0}^\infty\left|\frac{(a)_n}{(c)_n}\right|
\frac1{n!}\frac{n!}{(\Re s)^{n+1}}
=\frac1{\Re s}\sum_{n=0}^\infty
\left|\frac{(a)_n}{(c)_n}\right|(\Re s)^{-n},
\]
which converges by the ratio test (or terminates): the ratio
of consecutive terms is $|(a+n)/(c+n)|/\Re s\to1/\Re s<1$.
Hence, by dominated convergence (every partial sum is bounded in
modulus by the sum of the absolute terms, whose integral, by
monotone convergence, is the finite sum above), termwise
integration is allowed and
\begin{align*}
L\{M(a,c,t)\}(s)
&=\sum_{n=0}^\infty\frac{(a)_n}{(c)_n n!}
\frac{n!}{s^{n+1}}\\
&=\frac1s\sum_{n=0}^\infty
\frac{(a)_n(1)_n}{(c)_n n!}\left(\frac1s\right)^n
=\boxed{\frac1s\,{}_2F_1\left(a,1;c;\frac1s\right)},
\end{align*}
since $(1)_n=n!$.
\erratum{The ordinary definitions require
$c\notin\{0,-1,-2,\ldots\}$, unless a specified terminating
or limiting case removes the singularity. The half-plane
$\Re s>1$ is sufficient for the proof; continuation beyond it
must stay within the transform's convergence domain and use
consistent analytic branches. Terminating cases can have a
larger domain.}

\Needspace{9\baselineskip}
\subsection*{20.8 Properties of Laplace Transforms}

\Needspace{5\baselineskip}
\par\bigskip
\noindent\textbf{20.8.1}\quad Use the expression for the transform
of a second derivative to obtain the transform of $\cos kt$.
\solution
Put $F(t)=\cos kt$ and $f(s)=L\{F\}(s)$. The differential
identity $F''=-k^2F$ and initial values $F(0)=1$, $F'(0)=0$
give
\[
L\{F''\}=s^2f(s)-sF(0)-F'(0)=s^2f(s)-s=-k^2f(s),
\]
where the last equality transforms $F''=-k^2F$ by linearity.
The derivative formula applies for $\Re s>0$, since $F$ and
$F'=-k\sin kt$ are continuous and bounded. Collecting the $f$
terms gives $(s^2+k^2)f(s)=s$. For $\Re s>0$ the factor
$s^2+k^2$ does not vanish (its zeros $\pm ik$ lie on the
imaginary axis), so
\[
\boxed{L\{\cos kt\}(s)=\frac{s}{s^2+k^2}},\qquad \Re s>0
\quad(k\in\mathbb R).
\]
For $k=0$ this reduces to $L\{1\}=1/s$, as required.

\Needspace{5\baselineskip}
\par\bigskip
\noindent\textbf{20.8.2}\quad A mass $m$ is attached to one end of
an unstretched spring, spring constant $k$ (Fig.~20.17). Starting
at time $t=0$, the free end of the spring experiences a constant
acceleration $a$, away from the mass. Using Laplace transforms,
\begin{enumerate}
\item[(a)] find the position $x$ of $m$ as a function of time.
\item[(b)] determine the limiting form of $x(t)$ for small $t$.
\end{enumerate}
\par\nobreak\hfill\emph{ANS.}\enspace
(a) $x=\dfrac{1}{2}at^{2}-\dfrac{a}{\omega^{2}}(1-\cos\omega t)$,
$\omega^{2}=\dfrac{k}{m}$,\par\nobreak\hfill
(b) $x=\dfrac{a\omega^{2}}{4!}\,t^{4}$, $\omega t\ll1$.
\par\erratum{Use the initially resting preparation $x(0)=x'(0)=0$, with the driven end initially at rest and $m,k>0$; the textbook states no initial conditions, and its answer (a) presumes these.}
\solution
(a) Measure $x$ from the initial mass position, positive in
the direction of the acceleration $a$, and subtract
the unstretched spring length from the driven endpoint's
position. The latter coordinate is $y(t)=at^2/2$, so the spring
extension is $y-x$. Newton's law is
\[
m x''=k(y-x),\qquad
x''+\omega^2x=\frac{a\omega^2t^2}{2},\qquad \omega^2=k/m.
\]
The mass is initially at rest at $x=0$, so with
$\widehat x=L\{x\}$ we have
$L\{x''\}=s^2\widehat x-sx(0)-x'(0)=s^2\widehat x$; together
with $L\{t^2/2\}=1/s^3$, the transformed equation is
\[
(s^2+\omega^2)\widehat x=\frac{a\omega^2}{s^3},\qquad
\widehat x=\frac{a\omega^2}{s^3(s^2+\omega^2)}
=\frac a{s^3}-\frac a{s(s^2+\omega^2)}.
\]
The last step writes the numerator as
$\omega^2=(s^2+\omega^2)-s^2$. The same splitting gives
$1/[s(s^2+\omega^2)]=\omega^{-2}[1/s-s/(s^2+\omega^2)]$, and
the pairs $L\{t^2/2\}=1/s^3$, $L\{1\}=1/s$,
$L\{\cos\omega t\}=s/(s^2+\omega^2)$ then invert $\widehat x$
term by term:
\[
\boxed{x(t)=\frac12at^2-\frac a{\omega^2}(1-\cos\omega t)}.
\]
(b) Expanding the cosine shows explicitly the cancellation of the
quadratic terms:
\[
x(t)=\frac12at^2-\frac a{\omega^2}
\left(\frac{\omega^2t^2}{2}-\frac{\omega^4t^4}{24}
+\frac{\omega^6t^6}{720}-\cdots\right)
=\boxed{\frac{a\omega^2t^4}{24}+O(t^6)}.
\]
The small-time condition is $|\omega t|\ll1$. Differentiating
the exact answer gives $x''=a(1-\cos\omega t)$, which agrees
with $\omega^2(y-x)$ and verifies the equation of motion.

\Needspace{5\baselineskip}
\par\bigskip
\noindent\textbf{20.8.3}\quad Radioactive nuclei decay according
to the law
\[
\frac{dN}{dt}=-\lambda N,
\]
with $N$ the concentration of a given nuclide and $\lambda$ its
particular decay constant. This equation may be interpreted as
stating that the rate of decay is proportional to the number of
these radioactive nuclei present. They all decay independently.

Consider now a radioactive series of $n$ different nuclides, with
Nuclide 1 decaying into Nuclide 2, Nuclide 2 into Nuclide 3, etc.,
until reaching Nuclide $n$, which is stable. The concentrations of
the various nuclides satisfy the system of ODEs
\[
\frac{dN_{1}}{dt}=-\lambda_{1}N_{1},\qquad
\frac{dN_{2}}{dt}=\lambda_{1}N_{1}-\lambda_{2}N_{2},\qquad
\cdots,\qquad
\frac{dN_{n}}{dt}=\lambda_{n-1}N_{n-1}.
\]
\begin{enumerate}
\item[(a)] For the case $n=3$ find $N_{1}(t)$, $N_{2}(t)$, and
$N_{3}(t)$, with $N_{1}(0)=N_{0}$ and $N_{2}(0)=N_{3}(0)=0$.
\item[(b)] Find an approximate expression for $N_{2}$ and $N_{3}$,
valid for small $t$ when $\lambda_{1}\approx\lambda_{2}$.
\item[(c)] Find approximate expressions for $N_{2}$ and $N_{3}$,
valid for large $t$, when (1) $\lambda_{1}\gg\lambda_{2}$,
(2) $\lambda_{1}\ll\lambda_{2}$.
\end{enumerate}
\par\nobreak\hfill\emph{ANS.}\enspace
(a) $N_{1}(t)=N_{0}e^{-\lambda_{1}t}$,\quad
$N_{2}(t)=N_{0}\dfrac{\lambda_{1}}{\lambda_{2}-\lambda_{1}}
\big(e^{-\lambda_{1}t}-e^{-\lambda_{2}t}\big)$,\par\nobreak\hfill
$N_{3}(t)=N_{0}\left(1-\dfrac{\lambda_{2}}{\lambda_{2}-\lambda_{1}}
e^{-\lambda_{1}t}+\dfrac{\lambda_{1}}{\lambda_{2}-\lambda_{1}}
e^{-\lambda_{2}t}\right)$.\par\nobreak\hfill
(b) $N_{2}\approx N_{0}\lambda_{1}t$,\quad
$N_{3}\approx\dfrac{N_{0}}{2}\lambda_{1}\lambda_{2}t^{2}$.
\par\nobreak\hfill
(c) (1) $N_{2}\approx N_{0}e^{\erratum{-}\lambda_{2}t}$\par
\nobreak\hfill$N_{3}\approx N_{0}\big(1-e^{-\lambda_{2}t}\big)$,
\quad$\lambda_{1}t\gg1$.\par\nobreak\hfill
(2) $N_{2}\approx N_{0}(\lambda_{1}/\lambda_{2})
e^{-\lambda_{1}t}$,\par\nobreak\hfill
$N_{3}\approx N_{0}\big(1-e^{-\lambda_{1}t}\big)$,
\quad$\lambda_{2}t\gg1$.
\par\erratum{The quotient answers in part (a) are written for distinct rates; at equal rates take the continuous limits derived below.}
\solution
(a) Denote Laplace transforms by hats. By
$L\{N_i'\}=s\widehat N_i-N_i(0)$, the initial values
and the three rate equations give
\[
(s+\lambda_1)\widehat N_1=N_0,\quad
(s+\lambda_2)\widehat N_2=\lambda_1\widehat N_1,\quad
s\widehat N_3=\lambda_2\widehat N_2.
\]
Hence $\widehat N_1=N_0/(s+\lambda_1)$ and
$\widehat N_2=N_0\lambda_1/[(s+\lambda_1)(s+\lambda_2)]$.
For $\lambda_1\ne\lambda_2$ the partial fractions
$1/[(s+\lambda_1)(s+\lambda_2)]
=[1/(s+\lambda_1)-1/(s+\lambda_2)]/(\lambda_2-\lambda_1)$
and the pair $L\{e^{-\lambda t}\}=1/(s+\lambda)$ give
\begin{align*}
N_1(t)&=N_0e^{-\lambda_1t},\\
\widehat N_2&=\frac{N_0\lambda_1}{\lambda_2-\lambda_1}
\left(\frac1{s+\lambda_1}-\frac1{s+\lambda_2}\right),\\
N_2(t)&=\frac{N_0\lambda_1}{\lambda_2-\lambda_1}
(e^{-\lambda_1t}-e^{-\lambda_2t}).
\end{align*}
Adding the three original equations gives
$(N_1+N_2+N_3)'=0$, so their sum remains $N_0$. Subtraction
therefore gives
\[
N_3=N_0-N_1-N_2=N_0\left[1-\left(1+\frac{\lambda_1}
{\lambda_2-\lambda_1}\right)e^{-\lambda_1t}
+\frac{\lambda_1e^{-\lambda_2t}}{\lambda_2-\lambda_1}\right],
\]
and $1+\lambda_1/(\lambda_2-\lambda_1)=\lambda_2/(\lambda_2-\lambda_1)$,
so
\[
\boxed{N_3(t)=N_0\left[1-
\frac{\lambda_2e^{-\lambda_1t}-\lambda_1e^{-\lambda_2t}}
{\lambda_2-\lambda_1}\right]}.
\]
\erratum{The quotient formulas presume $\lambda_1\ne\lambda_2$.
At the removable equal-rate case $\lambda_1=\lambda_2=\lambda$,
$N_2=N_0\lambda t e^{-\lambda t}$ and
$N_3=N_0[1-(1+\lambda t)e^{-\lambda t}]$.}
They follow either from $\widehat N_2=N_0\lambda/(s+\lambda)^2$
and $L\{te^{-\lambda t}\}=1/(s+\lambda)^2$, or from the limit
$(e^{-\lambda_1t}-e^{-\lambda_2t})/(\lambda_2-\lambda_1)\to
-\partial_\lambda e^{-\lambda t}=te^{-\lambda t}$ of the
exponential difference; then $N_3=N_0-N_1-N_2$.

(b) Insert $e^{-\lambda t}=1-\lambda t+\lambda^2t^2/2
-\lambda^3t^3/6+O(t^4)$ into the formulas of (a):
\begin{align*}
e^{-\lambda_1t}-e^{-\lambda_2t}
&=(\lambda_2-\lambda_1)t-\frac{\lambda_2^2-\lambda_1^2}2t^2
+O(t^3),\\
\lambda_2e^{-\lambda_1t}-\lambda_1e^{-\lambda_2t}
&=(\lambda_2-\lambda_1)\left[1-\frac{\lambda_1\lambda_2}2t^2
+\frac{\lambda_1\lambda_2(\lambda_1+\lambda_2)}6t^3\right]
+O(t^4);
\end{align*}
in the second line the terms linear in $t$ cancel. Dividing by
$\lambda_2-\lambda_1$, with
$(\lambda_2^2-\lambda_1^2)/(\lambda_2-\lambda_1)=\lambda_1+\lambda_2$,
gives
\begin{align*}
N_2&=N_0\lambda_1\left[t-
\frac{\lambda_1+\lambda_2}{2}t^2+O(t^3)\right],\\
N_3&=N_0\lambda_1\lambda_2\left[
\frac{t^2}{2}-\frac{\lambda_1+\lambda_2}{6}t^3+O(t^4)\right].
\end{align*}
The leading terms are those stated. They hold when
$\max(\lambda_1,\lambda_2)t\ll1$, including equal or nearly
equal rates; closeness of the rates is not needed for this
small-time expansion. Indeed, Taylor's theorem and the rate
equations at $t=0$ give the same leading terms for all rates:
$N_2(0)=0$ and $N_2'(0)=\lambda_1N_1(0)=\lambda_1N_0$, while
$N_3(0)=0$, $N_3'(0)=\lambda_2N_2(0)=0$ and
$N_3''(0)=\lambda_2N_2'(0)=\lambda_1\lambda_2N_0$.

(c1) If $\lambda_1\gg\lambda_2$ and $\lambda_1t\gg1$, then
$e^{-\lambda_1t}=e^{-(\lambda_1-\lambda_2)t}e^{-\lambda_2t}$ is
negligible beside $e^{-\lambda_2t}$, because
$(\lambda_1-\lambda_2)t\simeq\lambda_1t\gg1$. Discarding it in
the exact formulas of (a) leaves
\[
N_2\simeq N_0\frac{\lambda_1}{\lambda_1-\lambda_2}e^{-\lambda_2t},
\quad
N_3\simeq N_0\left[1-
\frac{\lambda_1}{\lambda_1-\lambda_2}e^{-\lambda_2t}\right].
\]
Since $\lambda_1/(\lambda_1-\lambda_2)=1/(1-\lambda_2/\lambda_1)
=1+O(\lambda_2/\lambda_1)$,
these become $N_0e^{-\lambda_2t}$ and
$N_0(1-e^{-\lambda_2t})$.
\erratum{The positive exponent printed for $N_2$ in this
case must be negative. A positive exponent would violate both
the rate equation and $0\le N_2\le N_0$.}
The final $N_3$ stays accurate even while $N_3$ itself is
still small. By $N_3=N_0-N_1-N_2$ and part (a), it exceeds
the exact $N_3$ by
$N_1+N_2-N_0e^{-\lambda_2t}=N_0(e^{-\lambda_2t}-e^{-\lambda_1t})
[\lambda_1/(\lambda_1-\lambda_2)-1]$, that is, by
$N_0\lambda_2(e^{-\lambda_2t}-e^{-\lambda_1t})/(\lambda_1-\lambda_2)$.
Relative to $N_0(1-e^{-\lambda_2t})$ this is at most
$[\lambda_2/(\lambda_1-\lambda_2)]/(e^{\lambda_2t}-1)
\le1/[(\lambda_1-\lambda_2)t]\ll1$, by $e^x-1\ge x$.

(c2) If $\lambda_1\ll\lambda_2$ and $\lambda_2t\gg1$,
discard $e^{-\lambda_2t}$ instead, by the same argument with
the rates interchanged:
\[
N_2\simeq N_0\frac{\lambda_1}{\lambda_2-\lambda_1}e^{-\lambda_1t}
\simeq N_0\frac{\lambda_1}{\lambda_2}e^{-\lambda_1t},
\]
\[
N_3\simeq N_0\left[1-
\frac{\lambda_2}{\lambda_2-\lambda_1}e^{-\lambda_1t}\right]
\simeq N_0(1-e^{-\lambda_1t}).
\]
The last approximations keep the leading terms in the small
ratio $\lambda_1/\lambda_2$, using
$\lambda_2-\lambda_1=\lambda_2(1-\lambda_1/\lambda_2)$.
Since $N_3=N_0-N_1-N_2$ exactly, the final form of $N_3$ drops
precisely $N_2\simeq N_0(\lambda_1/\lambda_2)e^{-\lambda_1t}$.
Relative to $N_0(1-e^{-\lambda_1t})$ this is
$(\lambda_1/\lambda_2)/(e^{\lambda_1t}-1)\le1/(\lambda_2t)\ll1$,
by $e^x-1\ge x$; so the approximation stays accurate even while
$N_3$ itself is still small.

\Needspace{5\baselineskip}
\par\bigskip
\noindent\textbf{20.8.4}\quad The rate of formation of an isotope
in a nuclear reactor is given by
\[
\frac{dN_{2}}{dt}=\varphi\big[\sigma_{1}N_{1}(0)
-\sigma_{2}N_{2}(t)\big]-\lambda_{2}N_{2}(t).
\]
Here $N_{1}(0)$ is the concentration of the original isotope
(assumed constant), and $N_{2}$ is that of the newly formed
isotope. The first two terms on the right-hand side describe the
production and destruction of the new isotope via neutron
absorption; $\varphi$ is the neutron flux (units
cm$^{-2}$\,s$^{-1}$); $\sigma_{1}$ and $\sigma_{2}$ (units
cm$^{2}$) are neutron absorption cross sections. The final term
describes the radioactive decay of the new isotope, with decay
constant $\lambda_{2}$.
\begin{enumerate}
\item[(a)] Find the concentration $N_{2}$ of the new isotope as a
function of time.
\item[(b)] For original isotope $^{153}$Eu, $\sigma_{1}=400$
barns $=400\times10^{-24}$ cm$^{2}$, $\sigma_{2}=1000$ barns
$=1000\times10^{-24}$ cm$^{2}$, and
$\lambda_{2}=1.4\times10^{-9}$ s$^{-1}$. If $N_{1}(0)=10^{20}$
and $\varphi=10^{9}$ cm$^{-2}$\,s$^{-1}$, find $N_{2}$, the
concentration of $^{154}$Eu, after 1 year of continuous
irradiation. Is the assumption that $N_{1}$ is constant
justified?
\end{enumerate}
\par\erratum{In (a), retain $N_2(0)=N_{20}$ as an initial datum; for the numerical irradiation in (b), take $N_{20}=0$.}
\solution
(a) Set $A=\varphi\sigma_1N_1(0)$ and
$b=\varphi\sigma_2+\lambda_2$. The equation is
$N_2'+bN_2=A$. With $N_2(0)=N_{20}$, the transforms
$L\{N_2'\}=s\widehat N_2-N_{20}$ and $L\{A\}=A/s$ give
\[
(s+b)\widehat N_2=N_{20}+A/s,
\quad
\widehat N_2=\frac{N_{20}}{s+b}
+\frac A b\left(\frac1s-\frac1{s+b}\right),
\]
by the partial fractions $1/[s(s+b)]=[1/s-1/(s+b)]/b$.
Inverting with $L\{1\}=1/s$ and $L\{e^{-bt}\}=1/(s+b)$,
\[
\boxed{N_2(t)=N_{20}e^{-bt}+\frac A b(1-e^{-bt})}.
\]
\erratum{An initial concentration is needed. For irradiation
of the original isotope with none of the new isotope initially
present, use $N_{20}=0$; the formula above also covers a
nonzero initial inventory.} If $b=0$, the continuous limit is
$N_{20}+At$, since $(1-e^{-bt})/b\to t$ as $b\to0$.

(b) For the stated data,
\begin{align*}
A&=(10^9)(400\times10^{-24})(10^{20})
=4.00\times10^7\ \mathrm{s}^{-1},\\
b&=(10^9)(1000\times10^{-24})+1.4\times10^{-9}
=1.401\times10^{-9}\ \mathrm{s}^{-1}.
\end{align*}
Taking one year to mean $365$ days gives
$T=365\cdot86\,400=31\,536\,000$ s. Then
$bT=(1.401\times10^{-9})(31\,536\,000)=0.0441819$, so
$1-e^{-bT}=0.0432201$, while $A/b=2.85510\times10^{16}$.
For $N_{20}=0$,
\[
N_2(T)=\frac Ab\left(1-e^{-bT}\right)
=(2.85510\times10^{16})(0.0432201)
=\boxed{1.2340\times10^{15}}.
\]
To the precision of the given data (two significant figures
in $\lambda_2$) this is $1.2\times10^{15}$.
The parent ${}^{153}$Eu is stable, so it is consumed only by
neutron absorption: $N_1'=-\varphi\sigma_1N_1$ and
$N_1(T)=N_1(0)e^{-\varphi\sigma_1T}$. The depletion rate is
$\varphi\sigma_1=4.00\times10^{-13}$ s$^{-1}$ and
$\varphi\sigma_1T=1.26144\times10^{-5}$, so the fraction consumed is
\[
1-e^{-\varphi\sigma_1T}=1.26143\times10^{-5}
=0.00126143\%.
\]
Thus the assumption is justified: treating $N_1$ as constant
is accurate to about thirteen parts per million over this
interval. The calculation uses
the same number or concentration units as $N_1(0)$.

\Needspace{5\baselineskip}
\par\bigskip
\noindent\textbf{20.8.5}\quad In a nuclear reactor $^{135}$Xe is
formed as both a direct fission product of $^{235}$U and by decay
of $^{135}$I (another fission product), half-life 6.7 hours. The
half-life of $^{135}$Xe is 9.2 hours. Because $^{135}$Xe strongly
absorbs thermal neutrons, thereby ``poisoning'' the nuclear
reactor, its concentration is a matter of great interest. The
relevant equations are
\begin{align*}
\frac{dN_{\mathrm{I}}}{dt}
&=\varphi\gamma_{\mathrm{I}}(\sigma_{f}N_{\mathrm{U}})
-\lambda_{\mathrm{I}}N_{\mathrm{I}},\\
\frac{dN_{\mathrm{Xe}}}{dt}
&=\varphi\big[\gamma_{\mathrm{Xe}}(\sigma_{f}N_{\mathrm{U}})
-\sigma_{\mathrm{Xe}}N_{\mathrm{Xe}}\big]
+\lambda_{\mathrm{I}}N_{\mathrm{I}}
-\lambda_{\mathrm{Xe}}N_{\mathrm{Xe}}.
\end{align*}
Here $N_{\mathrm{I}}$, $N_{\mathrm{Xe}}$, $N_{\mathrm{U}}$ are the
concentrations of $^{135}$I, $^{135}$Xe, $^{235}$U, with
$N_{\mathrm{U}}$ assumed to be constant. The neutron flux
$\varphi$ in the reactor causes fission of $^{235}$U with cross
section $\sigma_{f}$ and removes $^{135}$Xe by neutron absorption
with cross section $\sigma_{\mathrm{Xe}}=3.5\times10^{6}$ barns
$=3.5\times10^{-18}$ cm$^{2}$. Neutron absorption by $^{135}$I is
negligible. The yield of $^{135}$I and $^{135}$Xe per fission
are, respectively, $\gamma_{\mathrm{I}}=0.060$ and
$\gamma_{\mathrm{Xe}}=0.003$.
\begin{enumerate}
\item[(a)] Find $N_{\mathrm{Xe}}(t)$ in terms of neutron flux
$\varphi$ and the product $\sigma_{f}N_{\mathrm{U}}$.
\item[(b)] Find $N_{\mathrm{Xe}}(t\to\infty)$.
\item[(c)] After $N_{\mathrm{Xe}}$ has reached equilibrium, the
reactor is shut down: $\varphi=0$. Find $N_{\mathrm{Xe}}(t)$
following shutdown. Notice the increase in $N_{\mathrm{Xe}}$,
which may for a few hours interfere with starting the reactor
up again. \erratum{The increase is not automatic: determine
when $N_{\mathrm{Xe}}$ increases immediately after shutdown;
this depends on the pre-shutdown flux.}
\end{enumerate}
\par\emph{Hint.}~The half-life $t_{1/2}$ of a radioactive
isotope is the time required for decay of half of the nuclides in
a sample. For a decay rate $dN/dt=-\lambda N$, the half-life has
the value $t_{1/2}=\ln 2/\lambda$, so $\lambda$ can be computed
as $\lambda=\ln 2/t_{1/2}=0.693/t_{1/2}$.
\par\erratum{In (a), retain initial inventories $N_{\rm I}(0)=I_0$ and $N_{\rm Xe}(0)=X_0$, then specialize to $I_0=X_0=0$.}
\solution
Keep the neutron flux and the uranium concentration
constant during operation, and abbreviate
\[
P=\varphi\sigma_fN_{\rm U},\quad
\alpha=\lambda_{\rm I},\quad\beta=\lambda_{\rm Xe},\quad
\mu=\beta+\varphi\sigma_{\rm Xe},\quad
I=N_{\rm I},\quad X=N_{\rm Xe}.
\]
The system becomes $I'+\alpha I=P\gamma_{\rm I}$ and
$X'+\mu X=P\gamma_{\rm Xe}+\alpha I$.
\erratum{The textbook's part (a) did not specify initial inventories.
We first keep $I(0)=I_0$, $X(0)=X_0$ arbitrary, and then give
the usual initially clean case $I_0=X_0=0$.}
The equilibrium values, obtained by setting both derivatives
to zero ($\alpha I_e=P\gamma_{\rm I}$ and
$\mu X_e=P\gamma_{\rm Xe}+\alpha I_e
=P(\gamma_{\rm Xe}+\gamma_{\rm I})$), are
\[
I_e=\frac{P\gamma_{\rm I}}\alpha,\qquad
X_e=\frac{P(\gamma_{\rm I}+\gamma_{\rm Xe})}{\mu}.
\]
(a) With $L\{I'\}=s\widehat I-I_0$, the first equation
transforms into
$(s+\alpha)\widehat I=I_0+P\gamma_{\rm I}/s=I_0+\alpha I_e/s$.
The partial fractions $\alpha/[s(s+\alpha)]=1/s-1/(s+\alpha)$
then give
\[
\widehat I=\frac{I_e}s+\frac{I_0-I_e}{s+\alpha},\qquad
I(t)=I_e+(I_0-I_e)e^{-\alpha t}.
\]
Next subtract $\mu X_e=P\gamma_{\rm Xe}+\alpha I_e$ from both
sides of the second equation:
\[
(X-X_e)'+\mu(X-X_e)=\alpha(I-I_e)
=\alpha(I_0-I_e)e^{-\alpha t}.
\]
Since $(X-X_e)(0)=X_0-X_e$, its transform is
\[
(s+\mu)L\{X-X_e\}=X_0-X_e+
\frac{\alpha(I_0-I_e)}{s+\alpha}.
\]
Divide by $s+\mu$, use
$1/[(s+\alpha)(s+\mu)]
=[1/(s+\alpha)-1/(s+\mu)]/(\mu-\alpha)$ (for $\mu\ne\alpha$),
and invert to obtain
\[
\boxed{X(t)=X_e+(X_0-X_e)e^{-\mu t}
+\frac{\alpha(I_0-I_e)}{\mu-\alpha}
(e^{-\alpha t}-e^{-\mu t})}.
\]
In particular, for zero initial inventories ($I_0=X_0=0$,
with $\alpha I_e=P\gamma_{\rm I}$),
\[
\boxed{X(t)=\frac{P(\gamma_{\rm I}+\gamma_{\rm Xe})}{\mu}
(1-e^{-\mu t})
-\frac{P\gamma_{\rm I}}{\mu-\alpha}
(e^{-\alpha t}-e^{-\mu t})}.
\]
\erratum{When $\mu=\alpha$, replace
$(e^{-\alpha t}-e^{-\mu t})/(\mu-\alpha)$ by $t e^{-\alpha t}$.
The populations remain finite at this coincident rate.}
This is the limit of the quotient as $\mu\to\alpha$, namely
$-\partial_\mu e^{-\mu t}=te^{-\mu t}$ at $\mu=\alpha$.
For example, in the initially clean case this gives
$X=X_e(1-e^{-\alpha t})-P\gamma_{\rm I}t e^{-\alpha t}$.

(b) Since the rates $\alpha$ and $\mu$ are positive, all
exponentials in (a) decay, so $X\to X_e$:
\[
\boxed{N_{\rm Xe}(\infty)=
\frac{\varphi(\gamma_{\rm I}+\gamma_{\rm Xe})\sigma_fN_{\rm U}}
{\lambda_{\rm Xe}+\varphi\sigma_{\rm Xe}}}.
\]
(c) Reset $t=0$ at shutdown. The initial values are now
$I_e,X_e$, but production and neutron absorption cease:
switching off the flux removes the terms $P\gamma_{\rm I}$,
$P\gamma_{\rm Xe}$ and $\varphi\sigma_{\rm Xe}X$, so the
equations become $I'=-\alpha I$ and $X'=-\beta X+\alpha I$.
Below, $\varphi$, $P$ and $\mu$ keep their pre-shutdown values,
which determine $I_e$ and $X_e$.
Thus $I=I_e e^{-\alpha t}$, and the second equation, with
$X(0)=X_e$, transforms into
\[
(s+\beta)\widehat X=X_e+\frac{\alpha I_e}{s+\alpha}.
\]
Here $\alpha>\beta$, because the iodine half-life is the
shorter. Dividing by $s+\beta$ and inverting with
$1/[(s+\alpha)(s+\beta)]=[1/(s+\beta)-1/(s+\alpha)]/(\alpha-\beta)$
gives
\[
\boxed{X(t)=X_e e^{-\beta t}
+\frac{\alpha I_e}{\alpha-\beta}
(e^{-\beta t}-e^{-\alpha t})}.
\]
\erratum{In the Instructor's Manual's shutdown formula the
coefficient of
the exponential difference uses $N_{\rm Xe}(0)$; it must use
$N_{\rm I}(0)=I_e$. The source of xenon after shutdown is
$\alpha I$, as the differential equation directly shows.}
The initial slope follows from the shutdown equation at $t=0$,
$X'(0)=\alpha I_e-\beta X_e$. With $\alpha I_e=P\gamma_{\rm I}$,
the value of $X_e$, and $\mu-\beta=\varphi\sigma_{\rm Xe}$,
\begin{align*}
X'(0)&=P\gamma_{\rm I}-\beta X_e
=\frac{P[\gamma_{\rm I}\mu-\beta(\gamma_{\rm I}+\gamma_{\rm Xe})]}
{\mu}\\
&=\frac{P(\gamma_{\rm I}\varphi\sigma_{\rm Xe}
-\beta\gamma_{\rm Xe})}{\mu}.
\end{align*}
\erratum{A short-term increase is conditional, not automatic:
for a positive pre-shutdown flux it occurs exactly when
$\varphi\sigma_{\rm Xe}>\beta\gamma_{\rm Xe}/\gamma_{\rm I}$.}
Here $\alpha=\ln2/(6.7\cdot3600)=2.87374\times10^{-5}$ s$^{-1}$
and $\beta=\ln2/(9.2\cdot3600)=2.09284\times10^{-5}$ s$^{-1}$.
Since $\gamma_{\rm Xe}/\gamma_{\rm I}=0.003/0.060=1/20$,
the threshold is
\[
\varphi>\frac{\beta}{20\sigma_{\rm Xe}}
=2.98977\times10^{11}\ \mathrm{cm}^{-2}\,\mathrm{s}^{-1}.
\]
For an explicit peak time put $A=\alpha I_e/(\alpha-\beta)>0$.
Then $X=(X_e+A)e^{-\beta t}-Ae^{-\alpha t}$, and
$X'=-\beta(X_e+A)e^{-\beta t}+\alpha Ae^{-\alpha t}=0$ gives
$e^{(\alpha-\beta)t}=\alpha A/[\beta(X_e+A)]$, that is,
\[
t_{\rm peak}=\frac1{\alpha-\beta}
\ln\frac{\alpha A}{\beta(X_e+A)}.
\]
This is positive exactly when $\alpha A>\beta(X_e+A)$, i.e.
$(\alpha-\beta)A=\alpha I_e>\beta X_e$, which is the
preceding increase condition $X'(0)>0$. Because
$e^{\beta t}X'=-\beta(X_e+A)+\alpha Ae^{-(\alpha-\beta)t}$
decreases in $t$, $X$ then has a single maximum, at
$t_{\rm peak}$; otherwise $X'<0$ for all $t>0$ and the maximum
after shutdown is at $t=0$.
Power reactors operate at thermal fluxes of about $10^{13}$ to
$10^{14}$ cm$^{-2}$\,s$^{-1}$, some $30$ to $300$ times this
threshold, so the increase noted in the textbook is the usual
case. For example, at $\varphi=10^{13}$ cm$^{-2}$\,s$^{-1}$
($\mu=5.59\times10^{-5}$ s$^{-1}$; $P$ cancels from the ratios)
the peak-time formula gives $t_{\rm peak}=6.4$ h with
$X(t_{\rm peak})=1.31X_e$, and solving $X(t)=X_e$ numerically
shows that $X$ falls back to $X_e$ about $17$ h after shutdown.

\Needspace{5\baselineskip}
\par\bigskip
\noindent\textbf{20.8.6}\quad Solve Eq.~(20.160), which describes
a damped simple harmonic oscillator, for $X(0)=X_{0}$,
$X'(0)=0$, and
\begin{enumerate}
\item[(a)] $b^{2}=4mk$ (critically damped),
\item[(b)] $b^{2}>4mk$ (overdamped).
\end{enumerate}
\par\nobreak\hfill\emph{ANS.}\enspace
(a) $X(t)=X_{0}e^{-(b/2m)t}\left(1+\dfrac{b}{2m}t\right)$.
\solution
The equation is $mX''+bX'+kX=0$. Put
$\alpha=b/(2m)$ and $\omega_0^2=k/m$, and divide by $m$:
$X''+2\alpha X'+\omega_0^2X=0$. With $X(0)=X_0$ and $X'(0)=0$,
the derivative formulas give $L\{X''\}=s^2\widehat X-sX_0$
and $L\{X'\}=s\widehat X-X_0$, so
\[
(s^2+2\alpha s+\omega_0^2)\widehat X=(s+2\alpha)X_0,
\quad
\widehat X=X_0\frac{(s+\alpha)+\alpha}
{(s+\alpha)^2+\omega_0^2-\alpha^2},
\]
where the denominator has been completed to a square,
$s^2+2\alpha s+\omega_0^2=(s+\alpha)^2+\omega_0^2-\alpha^2$.
Each case is inverted with the shift theorem
$L\{e^{-\alpha t}F(t)\}=f(s+\alpha)$.

(a) Critical damping means $\omega_0^2=\alpha^2$, so
the pairs $L\{1\}=1/s$ and $L\{t\}=1/s^2$, shifted, give
\[
\widehat X=X_0\left(\frac1{s+\alpha}
+\frac{\alpha}{(s+\alpha)^2}\right),\qquad
\boxed{X=X_0e^{-\alpha t}(1+\alpha t)}.
\]
(b) In the overdamped case define
$\kappa=\sqrt{\alpha^2-\omega_0^2}>0$, so the denominator is
$(s+\alpha)^2-\kappa^2$. The hyperbolic pairs
$L\{\cosh\kappa t\}=s/(s^2-\kappa^2)$ and
$L\{\sinh\kappa t\}=\kappa/(s^2-\kappa^2)$, shifted by
$s\mapsto s+\alpha$, give
\[
\widehat X=X_0\left[\frac{s+\alpha}{(s+\alpha)^2-\kappa^2}
+\frac\alpha\kappa\,\frac{\kappa}{(s+\alpha)^2-\kappa^2}\right],
\]
\[
\boxed{X=X_0e^{-\alpha t}
\left(\cosh\kappa t+\frac\alpha\kappa\sinh\kappa t\right)}.
\]
Writing $\cosh$ and $\sinh$ in exponentials, this is
\[
\frac{X_0}{2}\left[
\left(1+\frac\alpha\kappa\right)e^{(-\alpha+\kappa)t}
+\left(1-\frac\alpha\kappa\right)e^{(-\alpha-\kappa)t}\right].
\]
Both exponents $-\alpha\pm\kappa$ are the roots of the
characteristic equation $r^2+2\alpha r+\omega_0^2=0$; the
coefficients sum to $X_0$, and the initial derivative is
$(X_0/2)[(1+\alpha/\kappa)(\kappa-\alpha)
-(1-\alpha/\kappa)(\kappa+\alpha)]
=(X_0/2)[(\kappa-\alpha^2/\kappa)-(\kappa-\alpha^2/\kappa)]=0$.
Taking $\kappa\to0$ gives the critical result because
$\sinh\kappa t/\kappa\to t$.

\Needspace{5\baselineskip}
\par\bigskip
\noindent\textbf{20.8.7}\quad Again solve Eq.~(20.160), which
describes a damped simple harmonic oscillator, but this time for
$X(0)=0$, $X'(0)=v_{0}$, and
\begin{enumerate}
\item[(a)] $b^{2}<4mk$ (underdamped),
\item[(b)] $b^{2}=4mk$ (critically damped),
\item[(c)] $b^{2}>4mk$ (overdamped).
\end{enumerate}
\par\nobreak\hfill\emph{ANS.}\enspace
(a) $X(t)=\dfrac{v_{0}}{\omega_{1}}\,e^{-(b/2m)t}\sin\omega_{1}t$,
\par\nobreak\hfill
(b) $X(t)=v_{0}te^{-(b/2m)t}$.
\solution
With $\alpha=b/(2m)$, $\omega_0^2=k/m$, and initial data
$X(0)=0$, $X'(0)=v_0$, we have $L\{X''\}=s^2\widehat X-v_0$
and $L\{X'\}=s\widehat X$, so transformation of the equation
(divided by $m$) gives
\[
(s^2+2\alpha s+\omega_0^2)\widehat X=v_0,
\qquad
\widehat X=\frac{v_0}{(s+\alpha)^2+\omega_0^2-\alpha^2}.
\]
(a) If $\omega_1=\sqrt{\omega_0^2-\alpha^2}
=\sqrt{k/m-b^2/(4m^2)}>0$, then
$\widehat X=(v_0/\omega_1)\,\omega_1/[(s+\alpha)^2+\omega_1^2]$,
and the shifted sine transform
$L\{e^{-\alpha t}\sin\omega_1t\}=\omega_1/[(s+\alpha)^2+\omega_1^2]$
gives
\[
\boxed{X(t)=\frac{v_0}{\omega_1}e^{-\alpha t}\sin\omega_1t}.
\]
(b) At $\omega_0^2=\alpha^2$ the transform is
$v_0/(s+\alpha)^2$, whose inverse, by the shifted pair
$L\{te^{-\alpha t}\}=1/(s+\alpha)^2$, is
\[
\boxed{X(t)=v_0t e^{-\alpha t}}.
\]
(c) If $\kappa=\sqrt{\alpha^2-\omega_0^2}>0$, the
denominator factors as
$(s+\alpha)^2-\kappa^2=(s+\alpha-\kappa)(s+\alpha+\kappa)$, and
partial fractions give
\[
\widehat X=\frac{v_0}{2\kappa}
\left(\frac1{s+\alpha-\kappa}-\frac1{s+\alpha+\kappa}\right),
\]
so
\[
X(t)=\frac{v_0}{2\kappa}\left(e^{-(\alpha-\kappa)t}
-e^{-(\alpha+\kappa)t}\right)
=\boxed{\frac{v_0}{\kappa}e^{-\alpha t}\sinh\kappa t}.
\]
All three expressions have $X(0)=0$, $X'(0)=v_0$; for example,
in (a) $X'=v_0e^{-\alpha t}[\cos\omega_1t-(\alpha/\omega_1)
\sin\omega_1t]$.
The critical formula is the continuous limit of both others,
using $\sin z/z\to1$ and $\sinh z/z\to1$.

\Needspace{5\baselineskip}
\par\bigskip
\noindent\textbf{20.8.8}\quad The motion of a body falling in a
resisting medium may be described by
\[
m\,\frac{d^{2}X(t)}{dt^{2}}=mg-b\,\frac{dX(t)}{dt}
\]
when the retarding force is proportional to the velocity. Find
$X(t)$ and $dX(t)/dt$ for the initial conditions
\[
X(0)=\left.\frac{dX}{dt}\right|_{t=0}=0.
\]
\solution
Choose the downward direction positive and set
$\Gamma=b/m>0$. The equation is $X''+\Gamma X'=g$.
With both initial values zero, $L\{X''\}=s^2\widehat X$,
$L\{X'\}=s\widehat X$, and $L\{g\}=g/s$, so
\[
(s^2+\Gamma s)\widehat X=g/s,\qquad
\widehat X=\frac g{s^2(s+\Gamma)}.
\]
The partial fractions are
\[
\frac1{s^2(s+\Gamma)}
=\frac1{\Gamma s^2}-\frac1{\Gamma^2s}
+\frac1{\Gamma^2(s+\Gamma)}.
\]
Multiplication by $s^2(s+\Gamma)$ verifies the identity: its
right numerator is $(s+\Gamma)/\Gamma-s(s+\Gamma)/\Gamma^2
+s^2/\Gamma^2=1$. Inverting term by term with
$L^{-1}\{1/s^2\}=t$, $L^{-1}\{1/s\}=1$, and
$L^{-1}\{1/(s+\Gamma)\}=e^{-\Gamma t}$,
\[
X(t)=g\left[\frac t\Gamma-\frac1{\Gamma^2}
+\frac{e^{-\Gamma t}}{\Gamma^2}\right]
=\frac g\Gamma t-\frac g{\Gamma^2}(1-e^{-\Gamma t}),
\]
and differentiation gives $X'=(g/\Gamma)(1-e^{-\Gamma t})$.
With $\Gamma=b/m$,
\[
\boxed{X(t)=\frac{mg}{b}t-\frac{m^2g}{b^2}(1-e^{-bt/m})},
\qquad
\boxed{X'(t)=\frac{mg}{b}(1-e^{-bt/m})}.
\]
The velocity approaches the terminal value $mg/b$, since
$e^{-bt/m}\to0$. Expanding
$e^{-\Gamma t}=1-\Gamma t+\Gamma^2t^2/2-\Gamma^3t^3/6+\cdots$
gives
\[
X=\frac g\Gamma t-\frac g{\Gamma^2}
\left(\Gamma t-\frac{\Gamma^2t^2}2+\frac{\Gamma^3t^3}6
-\cdots\right)
=\frac{gt^2}2-\frac{g\Gamma t^3}6+\cdots,
\]
and similarly $X'=gt-g\Gamma t^2/2+\cdots$. Hence the
zero-drag limits are $X\to gt^2/2$ and $X'\to gt$: the
apparent singularities as $b\to0$ cancel.

\Needspace{5\baselineskip}
\par\bigskip
\noindent\textbf{20.8.9}\quad Ringing circuit. In certain
electronic devices, resistance, inductance, and capacitance are
placed in a circuit as shown in Fig.~20.18. A constant voltage is
maintained across the capacitance, keeping it charged. At time
$t=0$ the circuit is disconnected from the voltage source. Find
the voltages across each of the elements $R$, $L$, and $C$ as a
function of time. Assume $R$ to be small.
\par\emph{Hint.}~By Kirchhoff's laws
\[
I_{RL}+I_{C}=0\quad\text{and}\quad E_{R}+E_{L}=E_{C},
\]
where
\[
E_{R}=I_{RL}R,\qquad
E_{L}=L\,\frac{dI_{RL}}{dt},\qquad
E_{C}=\frac{q_{0}}{C}+\frac{1}{C}\int_{0}^{t}I_{C}\,dt,
\]
$q_{0}=$ initial charge of capacitor.
\par\erratum{Take $R,L,C>0$ and $R^2<4L/C$. Keep the initial branch current $I_{RL}(0)=I_0$ explicit; if the source was connected long enough for DC steady state, use $I_0=q_0/(CR)$.}
\solution
The drawing has a series $R,L$ branch in parallel with
$C$. Define $V_0=q_0/C$, write $I=I_{RL}$ for the current
through that branch and $I_0=I(0)$, and use the voltage
polarities in the hint.
Kirchhoff's laws, with the hint's $E_C$ differentiated
($E_C'=I_C/C$), imply
\[
I_C=-I,\qquad E_C'=-I/C,\qquad LI'+RI=E_C.
\]
Substituting $I=-CE_C'$ into the last relation gives
\[
LC E_C''+RC E_C'+E_C=0,\qquad
E_C(0)=V_0,\quad E_C'(0)=-I_0/C.
\]
Put $\alpha=R/(2L)$ and
$\omega=\sqrt{1/(LC)-\alpha^2}>0$; ``small $R$'' means
$R^2<4L/C$. With a hat denoting the Laplace transform, the
derivative formulas give the transforms
$s^2\widehat E_C-sV_0+I_0/C$ of $E_C''$ and
$s\widehat E_C-V_0$ of $E_C'$, so transformation and division by
$LC$ (note $R/L=2\alpha$) give
\[
\Big(s^2+2\alpha s+\frac1{LC}\Big)\widehat E_C
=(s+2\alpha)V_0-\frac{I_0}C,
\]
and $s^2+2\alpha s+1/(LC)=(s+\alpha)^2+\omega^2$. Hence
\[
\widehat E_C(s)=
\frac{(s+2\alpha)V_0-I_0/C}{(s+\alpha)^2+\omega^2}
=V_0\frac{s+\alpha}{(s+\alpha)^2+\omega^2}
+\frac{\alpha V_0-I_0/C}{(s+\alpha)^2+\omega^2}.
\]
Since $(s+\alpha)/[(s+\alpha)^2+\omega^2]$ and
$\omega/[(s+\alpha)^2+\omega^2]$ are the transforms of
$e^{-\alpha t}\cos\omega t$ and $e^{-\alpha t}\sin\omega t$
(shift theorem), inversion gives
\[
\boxed{E_C=e^{-\alpha t}\left[
V_0\cos\omega t+\frac{\alpha V_0-I_0/C}{\omega}\sin\omega t
\right]}.
\]
Write $B=(\alpha V_0-I_0/C)/\omega$ for the sine coefficient.
The product rule gives
\[
E_C'=e^{-\alpha t}\big[(\omega B-\alpha V_0)\cos\omega t
-(\alpha B+\omega V_0)\sin\omega t\big],
\]
where $\omega B-\alpha V_0=-I_0/C$ and, using
$\omega^2+\alpha^2=1/(LC)$,
$\alpha B+\omega V_0=[(\alpha^2+\omega^2)V_0-\alpha I_0/C]/\omega
=[V_0/(LC)-\alpha I_0/C]/\omega$. Multiplying by $-C$, the
cosine coefficient of $I=-C E_C'$ is $I_0$, while its sine
coefficient is $(V_0/L-\alpha I_0)/\omega$. Thus
\[
I=e^{-\alpha t}\left[I_0\cos\omega t+
\frac{V_0/L-\alpha I_0}{\omega}\sin\omega t\right],
\]
\[
\boxed{E_R=Re^{-\alpha t}\left[I_0\cos\omega t+
\frac{V_0/L-\alpha I_0}{\omega}\sin\omega t\right]},
\]
\[
\boxed{E_L=e^{-\alpha t}\left[(V_0-RI_0)\cos\omega t+
\frac{-\alpha V_0+(R\alpha-1/C)I_0}{\omega}\sin\omega t\right]}.
\]
The last expression is $E_C-E_R$ (Kirchhoff's voltage law):
its cosine coefficient is $V_0-RI_0$, and its sine coefficient
is $[\alpha V_0-I_0/C-R(V_0/L-\alpha I_0)]/\omega
=[-\alpha V_0+(R\alpha-1/C)I_0]/\omega$, because
$RV_0/L=2\alpha V_0$. It also equals $LI'$.
These formulas display both initial energy stores explicitly.

For the circuit exactly as drawn, if the constant source has
been connected long enough to establish a DC steady state,
the inductor has zero DC voltage but the resistor does not:
\[
\boxed{I_0=V_0/R}.
\]
Substitution gives, since then $I_0/C=V_0/(RC)$,
$R(V_0/L-\alpha I_0)=(2\alpha-\alpha)V_0=\alpha V_0$,
$V_0-RI_0=0$, and
$-\alpha V_0+(R\alpha-1/C)I_0=-V_0/(RC)$,
\begin{align*}
E_C&=V_0e^{-\alpha t}\left[
\cos\omega t+\frac{\alpha-1/(RC)}{\omega}\sin\omega t\right],\\
E_R&=V_0e^{-\alpha t}\left[
\cos\omega t+\frac\alpha\omega\sin\omega t\right],\\
E_L&=-\frac{V_0}{RC\omega}e^{-\alpha t}\sin\omega t.
\end{align*}
These have $E_R(0)=E_C(0)=V_0$, $E_L(0)=0$, as a prior
DC steady state requires.
\erratum{The Instructor's Manual instead uses damping $1/(2RC)$ and
$E_C(0)=q_0/C=0$. Those are not the equations or initial state
of the charged-capacitor circuit in Fig.~20.18. Its damping
is $R/(2L)$. An ideal inductor in series with $R$ does not
short the capacitor.}
\erratum{If a different switching preparation leaves the
inductor initially unenergized, set $I_0=0$ in the general
formulas; this is an additional initial condition, not the DC
steady state of the displayed circuit.}
In that case $E_C=V_0e^{-\alpha t}(\cos\omega t+
\alpha\sin\omega t/\omega)$,
$E_R=RV_0e^{-\alpha t}\sin\omega t/(L\omega)$, and
$E_L=V_0e^{-\alpha t}(\cos\omega t-
\alpha\sin\omega t/\omega)$.
The $R=0$ idealization admits this separately specified initial
state, but no finite steady DC current $V_0/R$.

\Needspace{5\baselineskip}
\par\bigskip
\noindent\textbf{20.8.10}\quad With $J_{0}(t)$ expressed as a
contour integral, apply the Laplace transform operation, reverse
the order of integration, and thus show that
\[
L\{J_{0}(t)\}=(s^{2}+1)^{-1/2},\quad\text{for }s>0.
\]
\solution
The generating function
$\exp[(t/2)(z-z^{-1})]=\sum_nJ_n(t)z^n$ (Section~14.1) is a
Laurent series in $z$, so the coefficient formula for its
$z^0$ term gives the contour representation
\[
J_0(t)=\frac1{2\pi i}\oint_{|z|=1}
\exp\left[\frac t2(z-z^{-1})\right]\frac{dz}{z},
\]
with counterclockwise orientation. On this contour
$z=e^{i\theta}$ and $z-z^{-1}=2i\sin\theta$, so the exponential
has modulus $1$.
For real $s>0$, multiplication by $e^{-st}$ makes the double
integral absolutely convergent (the integrand is bounded by
$e^{-st}$), and interchange is valid. The inner integral is
elementary, because the real part of
$s-(z-z^{-1})/2=s-i\sin\theta$ is $s>0$:
\begin{align*}
L\{J_0\}(s)
&=\frac1{2\pi i}\oint\frac{dz}{z}
\int_0^\infty e^{-t[s-(z-z^{-1})/2]}dt
=\frac1{2\pi i}\oint\frac{dz}{z[s-(z-z^{-1})/2]}\\
&=\frac1{2\pi i}\oint\frac{2\,dz}{1+2sz-z^2},
\end{align*}
since $z[s-(z-z^{-1})/2]=(2sz-z^2+1)/2$.
Let $D=\sqrt{s^2+1}>0$. The denominator vanishes at the roots
$z_\pm=s\pm D$ of $z^2-2sz-1=0$. Here $z_+>1$ and, since
$(D-s)(D+s)=D^2-s^2=1$,
$|z_-|=D-s=1/(D+s)<1$, so just $z_-$ lies inside.
At this simple pole the residue is the numerator divided by
the derivative $2s-2z$ of the denominator:
\[
\frac{2}{2s-2z_-}=\frac{2}{2D}=\frac1D.
\]
The residue theorem therefore gives
\[
\boxed{L\{J_0(t)\}(s)=\frac1{\sqrt{s^2+1}}},\qquad s>0.
\]
The transform extends analytically to $\Re s>0$ with the
branch whose square root is positive on the positive real axis.
Indeed, for $s=x+iy$ with $x>0$, $s^2+1=x^2-y^2+1+2ixy$ is real
only when $y=0$, and then it is positive; so the principal root
$\sqrt{s^2+1}$ is analytic in $\Re s>0$. The Laplace integral of
the bounded function $J_0$ is analytic there too, and the two
agree for $s>0$, hence throughout $\Re s>0$ by the identity
theorem.

\Needspace{5\baselineskip}
\par\bigskip
\noindent\textbf{20.8.11}\quad Develop the Laplace transform of
$J_{n}(t)$ from $L\{J_{0}(t)\}$ by using the Bessel function
recurrence relations.
\par\emph{Hint.}~Here is a chance to use mathematical induction
(Section~1.4).
\par\erratum{The formula derived by the indicated induction is for integer $n\ge0$.}
\solution
Write $g_n(s)=L\{J_n(t)\}(s)$,
$D=\sqrt{s^2+1}$, and $r=D-s$ for $s>0$.
Exercise~20.8.10 gives $g_0=1/D$. Since $J_0'=-J_1$ and
$J_0(0)=1$, the derivative formula
$L\{F'\}=sL\{F\}-F(0)$ gives
\[
g_1=-L\{J_0'\}=-[sg_0-J_0(0)]=1-s/D=(D-s)/D=r/D.
\]
For integer $n\ge1$, $2J_n'=J_{n-1}-J_{n+1}$ and
$J_n(0)=0$, so $L\{J_n'\}=sg_n$; transforming the recurrence
gives
\[
2sg_n=g_{n-1}-g_{n+1},\qquad
g_{n+1}=g_{n-1}-2sg_n.
\]
Assume $g_{n-1}=r^{n-1}/D$ and $g_n=r^n/D$ for some $n\ge1$;
this holds for $n=1$ by the two results above. Then
\[
g_{n+1}=\frac{r^{n-1}}D-2s\frac{r^n}D
=\frac{r^{n-1}}D(1-2sr).
\]
The identity (with $D^2=s^2+1$)
\[
r^2=(D-s)^2=D^2-2sD+s^2
=1+2s^2-2sD=1-2s(D-s)=1-2sr
\]
makes this $r^{n-1}r^2/D=r^{n+1}/D$. Induction establishes
\[
\boxed{L\{J_n(t)\}(s)=
\frac{(\sqrt{s^2+1}-s)^n}{\sqrt{s^2+1}}},\qquad n=0,1,2,\ldots.
\]
For negative integer order use $J_{-n}=(-1)^nJ_n$;
the displayed power formula is intended for nonnegative order.
Analytic continuation to $\Re s>0$ uses the branch specified
in Exercise~20.8.10.

\Needspace{5\baselineskip}
\par\bigskip
\noindent\textbf{20.8.12}\quad A calculation of the magnetic field
of a circular current loop in circular cylindrical coordinates
leads to the integral
\[
\int_{0}^{\infty}e^{-kz}kJ_{1}(ka)\,dk,\qquad\erratum{a>0,\quad\Re z>0}.
\]
Show that this integral is equal to $a/(z^{2}+a^{2})^{3/2}$.
\solution
\erratum{For an ordinary improper integral the boundary
$\Re z=0$, which the textbook's printed condition $\Re(z)\ge0$
includes, must
be excluded when $a>0$: the stated
integrand has oscillations with amplitude growing like
$\sqrt{k}$. The formula holds ordinarily for $\Re z>0$ and
has boundary values in an Abel or distributional sense.}
Scale Exercise~20.8.10: for $a>0$ and $\Re z>0$, substitute
$t=ka$ ($dk=dt/a$) and use that result, continued
analytically to $s=z/a$ with $\Re s>0$:
\[
\int_0^\infty e^{-kz}J_0(ka)dk
=\frac1a\int_0^\infty e^{-(z/a)t}J_0(t)dt
=\frac1a\,\frac1{\sqrt{(z/a)^2+1}}
=\frac1{\sqrt{z^2+a^2}}.
\]
By the chain rule and $J_0'=-J_1$, the $a$-derivative of the
integrand is $e^{-kz}kJ_0'(ka)=-e^{-kz}kJ_1(ka)$. Since
$|J_1(x)|\le1$ for real $x$ (Exercise~14.1.1), it is bounded
by the integrable $ke^{-k\Re z}$, so differentiation under the
integral sign with respect to $a$ is allowed:
\[
-\int_0^\infty e^{-kz}kJ_1(ka)dk
=\frac\partial{\partial a}(z^2+a^2)^{-1/2}
=-a(z^2+a^2)^{-3/2}.
\]
Therefore
\[
\boxed{\int_0^\infty e^{-kz}kJ_1(ka)dk
=\frac a{(z^2+a^2)^{3/2}}},\qquad \Re z>0.
\]
The square root is chosen positive for real $z>0$ and
continued throughout the right half-plane. To see the
boundary issue directly, at $z=0$,
$kJ_1(ka)\sim\sqrt{2k/(\pi a)}\cos(ka-3\pi/4)$;
its oscillatory primitive does not have a limit.
Nevertheless inserting $e^{-\epsilon k}$, that is, taking
$z=\epsilon>0$ in the boxed formula, and letting
$\epsilon\downarrow0$ gives the Abel value $a/(a^2)^{3/2}=1/a^2$.
For $z=iy$, the analogous limiting formula is singular at
$y=\pm a$; it is not a finite ordinary value there.
\erratum{The Instructor's Manual, p.~351, first integrates
$\int_0^\infty e^{-kz}J_1(ka)\,dk$ by parts, with
$J_1(ka)=-a^{-1}\,dJ_0(ka)/dk$, and then applies
$-\partial/\partial z$. In that intermediate integral it prints
$K_1(ka)$ where $J_1(ka)$ is required, and it replaces the
boundary term $[e^{-kz}J_0(ka)]_0^\infty=-1$ by $+1$. The
resulting constant, $-a^{-1}$ times this term, is independent
of $z$ and disappears on differentiation with respect to $z$,
but the intermediate equality is incorrect: the integral equals
$[1-z/\sqrt{z^2+a^2}]/a$.}

\Needspace{5\baselineskip}
\par\bigskip
\noindent\textbf{20.8.13}\quad Show that
\[
L\{I_{0}(at)\}=(s^{2}-a^{2})^{-1/2},\qquad\erratum{a\in\mathbb R,\quad s>|a|}.
\]
\solution
For real $a$ use the angular representation
(Exercise~14.5.14 with $n=0$; see also Exercise~14.6.3)
\[
I_0(at)=\frac1\pi\int_0^\pi e^{at\cos\theta}d\theta.
\]
\erratum{If $a$ is allowed to be negative, the sufficient
real convergence condition is $s>|a|$, not merely $s>a$.
For a positive scale $a$, the textbook's condition $s>a$ is correct.}
For $s>|a|$ the integrand of the double integral is bounded by
$e^{-st}e^{|a|t}=e^{-(s-|a|)t}$, so it converges absolutely and
the order of integration may be exchanged. The inner integral
is $\int_0^\infty e^{-(s-a\cos\theta)t}dt=1/(s-a\cos\theta)$,
so
\[
L\{I_0(at)\}(s)=\frac1\pi\int_0^\pi
\frac{d\theta}{s-a\cos\theta}.
\]
Set $u=\tan(\theta/2)$, so
$\cos\theta=(1-u^2)/(1+u^2)$ and
$d\theta=2du/(1+u^2)$; the range $0\le\theta<\pi$ becomes
$0\le u<\infty$, and
$s-a\cos\theta=[(s-a)+(s+a)u^2]/(1+u^2)$. Then
\begin{align*}
L\{I_0(at)\}
&=\frac2\pi\int_0^\infty
\frac{du}{(s-a)+(s+a)u^2}\\
&=\frac2{\pi\sqrt{(s-a)(s+a)}}
\left[\tan^{-1}\left(u\sqrt{\frac{s+a}{s-a}}\right)\right]_0^\infty\\
&=\frac2{\pi\sqrt{s^2-a^2}}\cdot\frac\pi2
=\boxed{\frac1{\sqrt{s^2-a^2}}}.
\end{align*}
The second line uses
$\int du/(p+qu^2)=(pq)^{-1/2}\tan^{-1}(u\sqrt{q/p})$ with
$p=s-a>0$ and $q=s+a>0$.
For complex $s$ in $\Re s>|a|$, use the branch positive on
$s>|a|$, equivalently the branch asymptotic to $s$ at positive
infinity in the denominator.

\Needspace{5\baselineskip}
\par\bigskip
\noindent\textbf{20.8.14}\quad Verify the following Laplace
transforms:
\begin{enumerate}
\item[(a)] $L\{j_{0}(at)\}=L\left\{\dfrac{\sin at}{at}\right\}
=\dfrac{1}{a}\cot^{-1}\left(\dfrac{s}{a}\right)$,
\item[(b)] $L\{y_{0}(at)\}$ does not exist.
\item[(c)] $L\{i_{0}(at)\}=L\left\{\dfrac{\sinh at}{at}\right\}
=\dfrac{1}{2a}\ln\dfrac{s+a}{s-a}
=\dfrac{1}{a}\coth^{-1}\left(\dfrac{s}{a}\right)$,
\item[(d)] $L\{k_{0}(at)\}$ does not exist.
\end{enumerate}
\par\erratum{Use $a>0$, with real $s>0$ in (a) and $s>a$ in (c), and the real inverse-function branches there.}
\solution
Take $a>0$; use real $s>0$ for (a) and $s>a$ for (c).
(a) Since $L\{\sin at\}(u)=a/(u^2+a^2)$, integrate the
transform parameter from $s$ to infinity, inserting the
defining integral and using $\int_s^\infty e^{-ut}du=e^{-st}/t$:
\[
\int_s^\infty\frac a{u^2+a^2}du
=\int_0^\infty\sin(at)\left(\int_s^\infty e^{-ut}du\right)dt
=\int_0^\infty e^{-st}\frac{\sin at}{t}dt.
\]
Absolute integrability after the factor $e^{-st}$ justifies
the interchange, because $|\sin at|/t\le a$. The left side is
$[\tan^{-1}(u/a)]_s^\infty=\pi/2-\tan^{-1}(s/a)$, and
$\pi/2-\tan^{-1}x=\cot^{-1}x$ for $x>0$. Since
$j_0(at)=\sin(at)/(at)$, dividing by $a$ yields
\[
\boxed{L\{j_0(at)\}=
\frac1a\left[\frac\pi2-\tan^{-1}(s/a)\right]
=\frac1a\cot^{-1}(s/a)}.
\]
(b) The spherical Neumann function is
$y_0(at)=-\cos(at)/(at)=-1/(at)+O(t)$ near zero, since
$\cos at=1+O(t^2)$.
Thus its Laplace integral contains a nonzero multiple of
$\int_0^\epsilon dt/t$ and diverges for every $s$.

(c) In the convergent half-plane $u>a$,
$L\{\sinh at\}(u)=a/(u^2-a^2)$. The same integration, now
justified by $0\le\sinh(at)/t\le a\cosh at\le ae^{at}$ and
$s>a$, followed by division by $a$ (since
$i_0(at)=\sinh(at)/(at)$), and the partial fractions
$1/(u^2-a^2)=[1/(u-a)-1/(u+a)]/(2a)$, give
\begin{align*}
L\{i_0(at)\}
&=\int_s^\infty\frac{du}{u^2-a^2}
=\frac1{2a}\left[\ln\frac{u-a}{u+a}\right]_s^\infty\\
&=\frac1{2a}\left[0-\ln\frac{s-a}{s+a}\right]
=\boxed{\frac1{2a}\ln\frac{s+a}{s-a}
=\frac1a\coth^{-1}(s/a)}.
\end{align*}
The two logarithms $\ln(u-a)$ and $\ln(u+a)$ must be kept
together at the infinite endpoint, where $(u-a)/(u+a)\to1$;
the individual logarithms do not have finite limits. The last
form uses $\coth^{-1}x=\frac12\ln[(x+1)/(x-1)]$ for $x>1$, with
$x=s/a$.
\erratum{The Instructor's Manual's final logarithmic expression for (c),
on p.~351, omits the factor $1/a$. Its coefficient must be
$1/(2a)$, as in the textbook answer and the partial fractions
shown here.}

(d) In the book's spherical normalization
$k_0(at)=e^{-at}/(at)$. It also has a nonintegrable $1/(at)$
singularity at zero, so its ordinary Laplace transform does
not exist.
\erratum{The real inverse-trigonometric and inverse-hyperbolic
forms are written for the positive scale and the domains
specified above; the $a=0$ limits of (a) and (c) are $1/s$.
Negative $a$ leaves $j_0(at)$ and $i_0(at)$ unchanged (they are
even in $a$). Both forms in (c) are even in $a$ too and remain
valid for $s>|a|$. In (a), however, $\cot^{-1}(s/a)$ at a
negative argument depends on the branch (values in $(0,\pi)$
give the wrong sign), so use $\tan^{-1}(a/s)/a$, which is even
in $a$ and valid for $s>0$.}

\Needspace{5\baselineskip}
\par\bigskip
\noindent\textbf{20.8.15}\quad Develop a Laplace transform
solution of Laguerre's equation,
\[
tF''(t)+(1-t)F'(t)+nF(t)=0.
\]
Note that you need a derivative of a transform and a transform of
derivatives. Go as far as you can with a general value of $n$;
then (and only then) set $n=0$.
\par\erratum{Find the solution regular at the origin, with finite $F(0)=C$ and $F'(0)$. A logarithmically singular second solution is outside these initial-value transform rules.}
\solution
Seek the solution regular at $t=0$, with finite
$F(0)=C$ and $F'(0)$. Write $f=L\{F\}$. The derivative
formulas and the derivative of a transform,
$L\{tG(t)\}=-\frac d{ds}L\{G\}$, give the needed transforms
\begin{align*}
L\{F'\}&=sf-C, &L\{F''\}&=s^2f-sC-F'(0),\\
L\{tF'\}&=-\frac d{ds}(sf-C)=-f-sf',
&L\{tF''\}&=-\frac d{ds}(s^2f-sC-F'(0))=-2sf-s^2f'+C.
\end{align*}
In the sum $L\{tF''\}+L\{F'\}-L\{tF'\}+nf$ the $f'$ terms give
$(-s^2+s)f'$, the $f$ terms give $(-2s+s+1+n)f$, and the
constants $C-C$ cancel, so the transformed equation is
\[
s(1-s)f'+(n+1-s)f=0.
\]
For a nonzero solution, divide by $s(1-s)f$ and decompose in
partial fractions; covering up $s$ at $s=0$ and $1-s$ at $s=1$
gives $(n+1-s)/[s(1-s)]=(n+1)/s+n/(1-s)$, so
\[
\frac{f'}f=-\frac{n+1-s}{s(1-s)}
=-\frac{n+1}{s}+\frac n{s-1}.
\]
Integrating each logarithmic derivative gives
\[
\ln f=\ln A+n\ln(s-1)-(n+1)\ln s,
\qquad f(s)=A\frac{(s-1)^n}{s^{n+1}}.
\]
For general $n$, take the powers positive for real $s>1$ and
continue them analytically into the half-plane $\Re s>1$.
The initial-value theorem (Exercise~20.7.1)
$\lim_{s\to\infty}sf(s)=F(0)$ gives $A=C$, because
$sf(s)=A(1-1/s)^n\to A$.
In this half-plane $|1/s|\le1/\Re s<1$, so expand by the
generalized
binomial theorem,
$(1-1/s)^n=\sum_j\binom nj(-1/s)^j$, where
$(-1)^j\binom nj=(-n)(-n+1)\cdots(-n+j-1)/j!=(-n)_j/j!$:
\[
f(s)=\frac C s\left(1-\frac1s\right)^n
=C\sum_{j=0}^\infty\frac{(-n)_j}{j!s^{j+1}}.
\]
Since $L\{t^j\}=j!/s^{j+1}$, inverting term by term identifies
the regular solution
\[
\boxed{F(t)=C\sum_{j=0}^\infty\frac{(-n)_j}{(j!)^2}t^j
=C\,{}_1F_1(-n;1;t)}.
\]
The termwise inversion is justified by Exercise~20.7.8 with
$a=-n$, $c=1$: for $\Re s>1$ it gives
$L\{M(-n,1,t)\}=s^{-1}{}_2F_1(-n,1;1;1/s)=s^{-1}(1-1/s)^n$,
because ${}_2F_1(-n,1;1;x)=\sum_j(-n)_jx^j/j!=(1-x)^n$; by
uniqueness of the inverse transform,
$F=C\,M(-n,1,t)=C\,{}_1F_1(-n;1;t)$.
If $n$ is a nonnegative
integer, the series terminates and $F=C L_n(t)$ (the
coefficients $(-n)_j/(j!)^2=(-1)^j\binom nj/j!$ are those of
the series (18.53) of $L_n$, in the ascending form derived in
Exercise~18.3.1), with
$f=C(s-1)^n/s^{n+1}$ already valid for $\Re s>0$.
Only now set $n=0$: $f=C/s$, so
\[
\boxed{F(t)=C,\qquad L_0(t)=1\text{ when }C=1}.
\]
\erratum{This is the regular solution selected by the finite
initial values used in transforming derivatives. The second
local solution is logarithmically singular at $t=0$ and is
not included by that derivation. At $n=0$, direct integration
gives $F'=B e^t/t$; its integral is singular unless $B=0$.}

\Needspace{5\baselineskip}
\par\bigskip
\noindent\textbf{20.8.16}\quad Show that the Laplace transform of
the Laguerre polynomial $L_{n}(at)$ is given by
\[
L\{L_{n}(at)\}=\frac{(s-a)^{n}}{s^{n+1}},\qquad s>0.
\]
\solution
For integer $n\ge0$, use the finite polynomial expansion
(Eq.~(18.53), in the ascending form derived in
Exercise~18.3.1, where
$n!/[(n-j)!(j!)^2]=\binom nj/j!$)
\[
L_n(x)=\sum_{j=0}^n(-1)^j\binom nj\frac{x^j}{j!}.
\]
There is no infinite-series interchange: transform its $n+1$
terms directly in $\Re s>0$, with
$L\{(at)^j\}=a^jj!/s^{j+1}$, and sum by the binomial theorem:
\begin{align*}
L\{L_n(at)\}(s)
&=\sum_{j=0}^n(-1)^j\binom nj\frac{a^j}{j!}
\frac{j!}{s^{j+1}}\\
&=\frac1s\sum_{j=0}^n\binom nj\left(-\frac a s\right)^j
=\frac1s\left(1-\frac a s\right)^n\\
&=\boxed{\frac{(s-a)^n}{s^{n+1}}}.
\end{align*}
This also includes $n=0$ and $a=0$. Because the integrand is
an exponential times a polynomial, no condition $s>a$ is
needed here; any real $s>0$ suffices.

\Needspace{5\baselineskip}
\par\bigskip
\noindent\textbf{20.8.17}\quad Show that
\[
L\{E_{1}(t)\}=\frac{1}{s}\ln(s+1),\qquad s>0,
\]
where
\[
E_{1}(t)=\int_{t}^{\infty}\frac{e^{-\tau}}{\tau}\,d\tau
=\int_{1}^{\infty}\frac{e^{-xt}}{x}\,dx.
\]
$E_{1}(t)$ is the exponential integral function, first
encountered in this book in Table~1.2.
\solution
For real $s>0$, the given representation has a
nonnegative integrand, so Tonelli's theorem permits the
interchange:
\begin{align*}
L\{E_1\}(s)
&=\int_0^\infty e^{-st}\int_1^\infty\frac{e^{-xt}}x\,dx\,dt\\
&=\int_1^\infty\frac{dx}{x}\int_0^\infty e^{-(s+x)t}dt
=\int_1^\infty\frac{dx}{x(s+x)}.
\end{align*}
Use $1/[x(s+x)]=s^{-1}[1/x-1/(x+s)]$ and a common upper
cutoff $R$:
\[
\int_1^R\frac{dx}{x(s+x)}
=\frac1s\big[\ln R-\ln(R+s)+\ln(1+s)\big].
\]
The first two logarithms approach $0$ together, since
$\ln R-\ln(R+s)=-\ln(1+s/R)\to0$ as $R\to\infty$. Therefore
\[
\boxed{L\{E_1(t)\}(s)=\frac{\ln(1+s)}s}.
\]
The apparent singularity at $s=0$ is removable: since
$\ln(1+s)=s-s^2/2+\cdots$, the limit is $1$. This is
consistent with
\[
\int_0^\infty E_1(t)dt
=\int_1^\infty\frac{dx}{x}\int_0^\infty e^{-xt}dt
=\int_1^\infty\frac{dx}{x^2}=1,
\]
obtained by reversing the order in its defining integral.

\Needspace{5\baselineskip}
\par\bigskip
\noindent\textbf{20.8.18}
\begin{enumerate}
\item[(a)] From Eq.~(20.184) show that
\[
\int_{0}^{\infty}f(x)\,dx=\int_{0}^{\infty}\frac{F(t)}{t}\,dt,
\]
provided the integrals exist.
\item[(b)] From the preceding result show that
\[
\int_{0}^{\infty}\frac{\sin t}{t}\,dt=\frac{\pi}{2},
\]
in agreement with Eqs.~(20.146) and (11.107).
\end{enumerate}
\par\erratum{In (a), justify the interchange, for example by $\int_0^\infty|F(t)|\,dt/t<\infty$; in (b), use an Abel limiting argument for the conditionally convergent integral.}
\solution
(a) Let $f(x)=\int_0^\infty e^{-xt}F(t)dt$.
\erratum{An interchange of improper integrals requires a
convergence justification. A simple sufficient condition for
the direct calculation is
$\int_0^\infty |F(t)|\,dt/t<\infty$; conditionally convergent
cases need a limiting argument, as in part (b).}
Under this absolute condition the double integral of
$e^{-xt}|F(t)|$ over $x,t>0$ equals
$\int_0^\infty|F(t)|\,dt/t<\infty$, so Fubini's theorem gives
\begin{align*}
\int_0^\infty f(x)dx
&=\int_0^\infty F(t)\left(\int_0^\infty e^{-xt}dx\right)dt\\
&=\boxed{\int_0^\infty\frac{F(t)}t\,dt},
\end{align*}
because the inner integral is $1/t$ for $t>0$.
Equivalently, Eq.~(20.184),
$L\{F(t)/t\}(s)=\int_s^\infty f(x)dx$ (the same interchange
with $\int_s^\infty e^{-xt}dx=e^{-st}/t$), yields the same
relation as $s\downarrow0$. Under the absolute condition this
limit is justified on both sides: the left side tends to
$\int_0^\infty F(t)\,dt/t$ by dominated convergence (dominating
function $|F(t)|/t$), and the right side to
$\int_0^\infty f(x)dx$ because
$\int_0^\infty|f(x)|dx\le\int_0^\infty|F(t)|\,dt/t<\infty$.

(b) For $F(t)=\sin t$, $f(x)=1/(1+x^2)$ for $x>0$. Insert a
positive lower cutoff $s$ in the $x$ integration. The double
integral of $e^{-xt}|\sin t|$ over $x>s$, $t>0$ is
$\int_0^\infty e^{-st}|\sin t|\,dt/t\le\int_0^\infty e^{-st}dt
=1/s$, since $|\sin t|\le t$; this absolute integrability
permits the interchange of (a), now with
$\int_s^\infty e^{-xt}dx=e^{-st}/t$:
\[
\int_0^\infty e^{-st}\frac{\sin t}{t}dt
=\int_s^\infty\frac{dx}{1+x^2}
=\frac\pi2-\tan^{-1}s.
\]
The unregulated sine integral exists: its integrand tends to
$1$ as $t\to0$, and integration by parts bounds its tail,
\[
\int_R^\infty\frac{\sin t}t\,dt
=\frac{\cos R}R-\int_R^\infty\frac{\cos t}{t^2}dt,
\qquad
\left|\int_R^\infty\frac{\sin t}t\,dt\right|
\le\frac1R+\int_R^\infty\frac{dt}{t^2}=\frac2R.
\]
Abel's theorem for integrals (Exercise~20.2.2(c)) thus lets
$s\downarrow0$ in the regulated integral without changing its
improper value, and $\tan^{-1}s\to0$.
Consequently
\[
\boxed{\int_0^\infty\frac{\sin t}{t}dt=\frac\pi2}.
\]
This proof uses a regulator precisely because the absolute
condition of part (a) fails for $|\sin t|/t$.

\Needspace{5\baselineskip}
\par\bigskip
\noindent\textbf{20.8.19}
\begin{enumerate}
\item[(a)] Show that
\[
L\left\{\frac{\sin kt}{t}\right\}=\cot^{-1}\left(\frac{s}{k}\right).
\]
\item[(b)] Using this result (with $k=1$), prove that
\[
L\{\mathrm{si}(t)\}=-\frac{1}{s}\tan^{-1}s,
\]
where
\[
\mathrm{si}(t)=-\int_{t}^{\infty}\frac{\sin x}{x}\,dx,
\qquad\text{the sine integral.}
\]
\end{enumerate}
\par\erratum{The real inverse-cotangent expression in (a) assumes $k>0$, $s>0$, and values in $(0,\pi/2)$. Use $s>0$ also in (b).}
\solution
(a) For real $s>0$, integrate the sine transform
$L\{\sin kt\}(u)=k/(u^2+k^2)$ in its transform parameter, by
the interchange of Exercise~20.8.14(a):
\[
L\left\{\frac{\sin kt}{t}\right\}(s)
=\int_s^\infty\frac{k}{u^2+k^2}du
=\boxed{\tan^{-1}(k/s)}.
\]
For $k>0$ the integral is
$[\tan^{-1}(u/k)]_s^\infty=\pi/2-\tan^{-1}(s/k)$, which equals
$\tan^{-1}(k/s)$ by $\tan^{-1}x+\tan^{-1}(1/x)=\pi/2$ for $x>0$;
for $k<0$ the same steps hold with $-\pi/2$ in place of $\pi/2$,
and for $k=0$ both sides vanish.
For $k>0$ this is $\cot^{-1}(s/k)$ with the usual value in
$(0,\pi/2)$.
\erratum{The printed inverse-cotangent form presumes $k>0$
and this branch. The arctangent form extends continuously to
$k=0$ and all real $k$.}

(b) From the defining tail integral,
$\mathrm{si}'(t)=\sin t/t$ (differentiate at the lower limit)
and $\mathrm{si}(0+)=-\int_0^\infty(\sin x/x)\,dx=-\pi/2$ by
Exercise~20.8.18(b). If $h(s)=L\{\mathrm{si}\}(s)$,
the derivative formula $L\{\mathrm{si}'\}=sh(s)-\mathrm{si}(0+)$
($\mathrm{si}$ is continuous and bounded, with a finite limit at
$0$) and part (a) with $k=1$ give
\[
s h(s)-\mathrm{si}(0+)
=s h(s)+\pi/2=\tan^{-1}(1/s).
\]
For $s>0$, $\tan^{-1}(1/s)=\pi/2-\tan^{-1}s$, so
$sh(s)=-\tan^{-1}s$ and
\[
\boxed{L\{\mathrm{si}(t)\}(s)=-\frac{\tan^{-1}s}{s}}.
\]
The finite value at the origin and the decaying oscillatory
tail of $\mathrm{si}$ ensure convergence in $\Re s>0$;
analytic continuation there uses the corresponding logarithmic
branch of the arctangent.

\Needspace{5\baselineskip}
\par\bigskip
\noindent\textbf{20.8.20}\quad If $F(t)$ is periodic
(Fig.~20.19) with a period $a$ so that $F(t+a)=F(t)$ for all
$t\ge0$, show that
\[
L\{F(t)\}=\frac{1}{1-e^{-as}}\int_{0}^{a}e^{-st}F(t)\,dt.
\]
Note that the integration is now over only the first period of
$F(t)$.
\par\erratum{Assume $a>0$, that $F$ is integrable over a period, and $\Re s>0$.}
\solution
Assume $a>0$ and $F$ integrable over one period.
For $\Re s>0$ split the half-line into successive periods,
putting $t=na+u$ ($0\le u\le a$) in the $n$th interval, so that
$e^{-st}=e^{-nas}e^{-su}$, and $F(na+u)=F(u)$ by $n$
applications of $F(t+a)=F(t)$:
\begin{align*}
L\{F\}(s)
&=\sum_{n=0}^\infty\int_{na}^{(n+1)a}e^{-st}F(t)dt\\
&=\sum_{n=0}^\infty e^{-nas}
\int_0^a e^{-su}F(na+u)du\\
&=\left(\sum_{n=0}^\infty e^{-nas}\right)
\int_0^a e^{-su}F(u)du.
\end{align*}
The geometric series converges because
$|e^{-as}|=e^{-a\Re s}<1$, and its sum is
$\sum_{n\ge0}(e^{-as})^n=1/(1-e^{-as})$.
It also bounds the absolute integrals,
$\sum_n\int_{na}^{(n+1)a}|e^{-st}F(t)|dt
\le\int_0^a|F(u)|du/(1-e^{-a\Re s})<\infty$, justifying the
splitting and summation. Hence
\[
\boxed{L\{F(t)\}(s)=\frac{\displaystyle\int_0^a e^{-st}F(t)dt}
{1-e^{-as}}},\qquad\Re s>0.
\]
Values assigned at isolated period boundaries do not affect
the transform.

\Needspace{5\baselineskip}
\par\bigskip
\noindent\textbf{20.8.21}\quad Find the Laplace transform of the
square wave (period $a$) defined by
\[
F(t)=\begin{cases}
1, & 0<t<a/2,\\
0, & a/2<t<a.
\end{cases}
\]
\par\nobreak\hfill\emph{ANS.}\enspace
$f(s)=\dfrac{1}{s}\,\dfrac{1-e^{-as/2}}{1-e^{-as}}$.
\par\erratum{Take $a>0$ and $\Re s>0$.}
\solution
Apply Exercise~20.8.20. Since $F=0$ on $(a/2,a)$, only the
first half of a period contributes to the numerator:
\[
\int_0^a e^{-st}F(t)dt
=\int_0^{a/2}e^{-st}dt=\left[-\frac{e^{-st}}s\right]_0^{a/2}
=\frac{1-e^{-as/2}}s.
\]
Dividing by $1-e^{-as}$ gives
\[
\boxed{f(s)=\frac{1-e^{-as/2}}{s(1-e^{-as})}
=\frac1{s(1+e^{-as/2})}},\qquad\Re s>0.
\]
The simplification follows from
$1-e^{-as}=(1-e^{-as/2})(1+e^{-as/2})$.
For real $s\to\infty$, $sf(s)=1/(1+e^{-as/2})\to1$, which
matches $F(0+)$ (the initial-value theorem,
Exercise~20.7.1).
For $s\downarrow0$, $sf(s)\to1/(1+1)=1/2$, the average value
over a period; this is not a pointwise final value of the
oscillating square wave.

\Needspace{5\baselineskip}
\par\bigskip
\noindent\textbf{20.8.22}\quad Show that
\begin{enumerate}
\item[(a)] $L\{\cosh at\cos at\}=\dfrac{s^{3}}{s^{4}+4a^{4}}$,
\item[(b)] $L\{\cosh at\sin at\}=\dfrac{as^{2}+2a^{3}}{s^{4}+4a^{4}}$,
\item[(c)] $L\{\sinh at\cos at\}=\dfrac{as^{2}-2a^{3}}{s^{4}+4a^{4}}$,
\item[(d)] $L\{\sinh at\sin at\}=\dfrac{2a^{2}s}{s^{4}+4a^{4}}$.
\end{enumerate}
\par\erratum{Take $a\in\mathbb R$ and $\Re s>|a|$.}
\solution
Take real $a$ and $\Re s>|a|$. Define
\[
D_-=(s-a)^2+a^2=s^2-2as+2a^2,\qquad
D_+=(s+a)^2+a^2=s^2+2as+2a^2.
\]
The substitution (frequency-shifting) theorem
$L\{e^{bt}F(t)\}=f(s-b)$, applied with $b=\pm a$ to
$L\{\cos at\}=s/(s^2+a^2)$ and $L\{\sin at\}=a/(s^2+a^2)$,
gives, since $(s\mp a)^2+a^2=D_\mp$,
\[
L\{e^{\pm at}\cos at\}=\frac{s\mp a}{D_\mp},\qquad
L\{e^{\pm at}\sin at\}=\frac a{D_\mp}.
\]
Also $D_\mp=(s^2+2a^2)\mp2as$, so
\[
D_-D_+=(s^2+2a^2)^2-(2as)^2=s^4+4a^4,\qquad
D_++D_-=2s^2+4a^2,\qquad D_+-D_-=4as.
\]
Using $\cosh at=(e^{at}+e^{-at})/2$ and
$\sinh at=(e^{at}-e^{-at})/2$, each transform is half the sum
or difference of two shifted transforms, brought to the common
denominator $D_-D_+$:
\begin{align*}
\text{(a)}\quad L\{\cosh at\cos at\}
&=\frac{(s-a)D_++(s+a)D_-}{2D_-D_+}
=\boxed{\frac{s^3}{s^4+4a^4}},\\
\text{(b)}\quad L\{\cosh at\sin at\}
&=\frac{a(D_++D_-)}{2D_-D_+}
=\boxed{\frac{as^2+2a^3}{s^4+4a^4}},\\
\text{(c)}\quad L\{\sinh at\cos at\}
&=\frac{(s-a)D_+-(s+a)D_-}{2D_-D_+}
=\boxed{\frac{as^2-2a^3}{s^4+4a^4}},\\
\text{(d)}\quad L\{\sinh at\sin at\}
&=\frac{a(D_+-D_-)}{2D_-D_+}
=\boxed{\frac{2a^2s}{s^4+4a^4}}.
\end{align*}
In (b) and (d) the numerators are
$a(2s^2+4a^2)$ and $a\cdot4as$, which halve to the boxed ones.
For the two less immediate numerators, direct expansion gives
$(s-a)D_+=s^3+as^2-2a^3$ and
$(s+a)D_-=s^3-as^2+2a^3$; adding gives $2s^3$ and subtracting
gives $2as^2-4a^3$, which proves the simplifications in (a)
and (c).

\Needspace{5\baselineskip}
\par\bigskip
\noindent\textbf{20.8.23}\quad Show that
\begin{enumerate}
\item[(a)] $L^{-1}\left\{(s^{2}+a^{2})^{-2}\right\}
=\dfrac{1}{2a^{3}}\sin at-\dfrac{t}{2a^{2}}\cos at$,
\item[(b)] $L^{-1}\left\{s(s^{2}+a^{2})^{-2}\right\}
=\dfrac{t}{2a}\sin at$,
\item[(c)] $L^{-1}\left\{s^{2}(s^{2}+a^{2})^{-2}\right\}
=\dfrac{1}{2a}\sin at+\dfrac{t}{2}\cos at$,
\item[(d)] $L^{-1}\left\{s^{3}(s^{2}+a^{2})^{-2}\right\}
=\cos at-\dfrac{at}{2}\sin at$.
\end{enumerate}
\par\erratum{Take real $a\ne0$ in the displayed quotients and use their continuous limits at $a=0$.}
\solution
Begin with
$L\{\sin at/a\}=1/(s^2+a^2)$ and
$L\{\cos at\}=s/(s^2+a^2)$ for real $a\ne0$, $\Re s>0$.
(a) Differentiate the first pair with respect to $a$ and
multiply by $-1/(2a)$. Differentiation under the transform
integral is allowed because the $a$-derivative of the
original,
$\partial_a(\sin at/a)=t\cos at/a-\sin at/a^2$, is bounded by
$t/|a|+1/a^2$; and
$\partial_a(s^2+a^2)^{-1}=-2a(s^2+a^2)^{-2}$. Hence
\[
-\frac1{2a}\frac\partial{\partial a}\frac1{s^2+a^2}
=\frac1{(s^2+a^2)^2},\quad
-\frac1{2a}\frac\partial{\partial a}\frac{\sin at}{a}
=\boxed{\frac{\sin at}{2a^3}-\frac{t\cos at}{2a^2}}.
\]
(b) Since $L\{tF\}=-dL\{F\}/ds$, take $F=\sin at/a$, whose
transform is $1/(s^2+a^2)$; then $\frac12tF$ has the transform
\[
-\frac12\frac d{ds}\frac1{s^2+a^2}
=\frac{s}{(s^2+a^2)^2},
\qquad
\boxed{L^{-1}\left\{\frac{s}{(s^2+a^2)^2}\right\}
=\frac t{2a}\sin at}.
\]
(c) Split the numerator as $s^2=(s^2+a^2)-a^2$ and use (a):
\[
\frac{s^2}{(s^2+a^2)^2}
=\frac1{s^2+a^2}-\frac{a^2}{(s^2+a^2)^2}.
\]
Its inverse is
$\sin at/a-a^2[\sin at/(2a^3)-t\cos at/(2a^2)]
=\sin at/a-\sin at/(2a)+(t/2)\cos at$, or
\[
\boxed{\frac{\sin at}{2a}+\frac t2\cos at}.
\]
(d) Similarly, from $s^3=s(s^2+a^2)-a^2s$,
$s^3/(s^2+a^2)^2=s/(s^2+a^2)-a^2s/(s^2+a^2)^2$.
The pair $L\{\cos at\}=s/(s^2+a^2)$ and part (b) then give
$\cos at-a^2\,t\sin at/(2a)$, that is,
\[
\boxed{\cos at-\frac{at}{2}\sin at}.
\]
\erratum{At $a=0$, interpret all quotient formulas by their
continuous limits. The four answers reduce to $t^3/6$,
$t^2/2$, $t$, and $1$, the inverses of $s^{-4}$ through
$s^{-1}$.}
For (a), $\sin at-at\cos at=a^3t^3/3+O(a^5)$ by the Maclaurin
series, so $(\sin at-at\cos at)/(2a^3)\to t^3/6$; the other
three limits follow from $\sin at/a\to t$ and $\cos at\to1$.

\Needspace{5\baselineskip}
\par\bigskip
\noindent\textbf{20.8.24}\quad Show that
\[
L\big\{(t^{2}-k^{2})^{-1/2}u(t-k)\big\}=K_{0}(ks).
\]
\par\erratum{(Here $u(t-k)$ is the unit step function.) This
clarification is from the published errata, which list it
under the misprinted number 20.8.34.}
\par\emph{Hint.}~Try transforming an integral representation of
$K_{0}(ks)$ into the Laplace transform integral.
\par\erratum{Take $k>0$ and $\Re s>0$; the step restricts the real square root to $t>k$.}
\solution
\erratum{Take $k>0$ and $\Re s>0$. The lower endpoint
$t=k$ then has an integrable square-root singularity. At
$k=0$ the kernel is $1/t$ and the ordinary transform diverges.}
The unit step $u(t-k)$ ($0$ for $t<k$, $1$ for $t>k$)
restricts the transform to
\[
L\{(t^2-k^2)^{-1/2}u(t-k)\}
=\int_k^\infty\frac{e^{-st}}{\sqrt{t^2-k^2}}dt.
\]
Set $t=k\cosh v$, $v\ge0$; the limits $t=k$ and $t\to\infty$
become $v=0$ and $v\to\infty$. Since
$dt=k\sinh v\,dv$ and, by $\cosh^2v-1=\sinh^2v$ with
$\sinh v\ge0$,
$\sqrt{t^2-k^2}=k\sinh v$, the Jacobian cancels the denominator:
\[
\int_k^\infty\frac{e^{-st}}{\sqrt{t^2-k^2}}dt
=\int_0^\infty e^{-ks\cosh v}dv=\boxed{K_0(ks)}.
\]
The last expression is the standard hyperbolic integral for
$K_0$ on $\Re(ks)>0$ (Exercise~14.5.15, or Exercise~14.6.3 with
$\nu=0$). The value assigned to $u(0)$ does not
change this integrable transform.

\Needspace{9\baselineskip}
\subsection*{20.9 Laplace Convolution Theorem}

\Needspace{5\baselineskip}
\par\bigskip
\noindent\textbf{20.9.1}\quad From the convolution theorem show
that
\[
\frac{1}{s}\,f(s)=L\left\{\int_{0}^{t}F(x)\,dx\right\},
\]
where $f(s)=L\{F(t)\}$.
\solution
In the one-sided Laplace convolution, choose $G(t)=1$
for $t\ge0$. Its transform is
$g(s)=\int_0^\infty e^{-st}dt=1/s$ for $\Re s>0$. Since
$G(t-x)=1$ for $0\le x\le t$, the definition gives
\[
(F*G)(t)=\int_0^tF(x)G(t-x)dx=\int_0^tF(x)dx.
\]
The convolution theorem, $L\{F*G\}=f(s)g(s)$, therefore yields
\[
\boxed{L\left\{\int_0^tF(x)dx\right\}(s)
=f(s)g(s)=\frac{f(s)}s}.
\]
For locally integrable $F$ of exponential order this
holds in the common half-plane where $\Re s>0$ and the
integral defining $f(s)$ converges. As a check, the primitive
$P(t)=\int_0^tF(x)dx$ has $P'=F$ and $P(0)=0$, so the
derivative formula $L\{P'\}=sL\{P\}-P(0)$ gives
$f(s)=sL\{P\}$, which is the same result.

\Needspace{5\baselineskip}
\par\bigskip
\noindent\textbf{20.9.2}\quad If $F(t)=t^{a}$ and $G(t)=t^{b}$,
$a>-1$, $b>-1$,
\begin{enumerate}
\item[(a)] Show that the convolution $F*G$ is given by
\[
F*G=t^{a+b+1}\int_{0}^{1}y^{a}(1-y)^{b}\,dy.
\]
\item[(b)] By using the convolution theorem, show that
\[
\int_{0}^{1}y^{a}(1-y)^{b}\,dy
=\erratum{\frac{\Gamma(a+1)\Gamma(b+1)}{\Gamma(a+b+2)}}=B(a+1,b+1),
\]
where $B$ is the beta function.
\end{enumerate}
\solution
(a) For $t>0$, substitute $x=ty$ into the convolution. Then
$dx=t\,dy$, $x^a=t^ay^a$, $(t-x)^b=t^b(1-y)^b$, and the limits
$x=0,t$ become $y=0,1$:
\[
(F*G)(t)=\int_0^t x^a(t-x)^b dx
=t^{a+b+1}\int_0^1 y^a(1-y)^bdy.
\]
The endpoint conditions $a,b>-1$ make this integral finite:
near $y=0$ the integrand behaves as $y^a$, and near $y=1$ as
$(1-y)^b$.

(b) Call the integral of (a) $C$. For $\nu>-1$ and $s>0$, the
substitution $u=st$ in the gamma integral gives
\[
L\{t^\nu\}=\int_0^\infty e^{-st}t^\nu dt
=\frac1{s^{\nu+1}}\int_0^\infty e^{-u}u^\nu du
=\frac{\Gamma(\nu+1)}{s^{\nu+1}}.
\]
With $\nu=a+b+1>-1$, $\nu=a$ and $\nu=b$ in turn,
\[
L\{F*G\}=C\frac{\Gamma(a+b+2)}{s^{a+b+2}},\qquad
f(s)=\frac{\Gamma(a+1)}{s^{a+1}},\qquad
g(s)=\frac{\Gamma(b+1)}{s^{b+1}}.
\]
The convolution theorem $L\{F*G\}=fg$ now reads
$C\,\Gamma(a+b+2)s^{-(a+b+2)}=\Gamma(a+1)\Gamma(b+1)s^{-(a+b+2)}$.
Cancelling $s^{-(a+b+2)}$ and dividing by $\Gamma(a+b+2)$ yields
\[
\boxed{\int_0^1y^a(1-y)^bdy
=\frac{\Gamma(a+1)\Gamma(b+1)}{\Gamma(a+b+2)}
=B(a+1,b+1)}.
\]
\erratum{The textbook writes this ratio as $a!\,b!/(a+b+1)!$;
since $a,b>-1$ need not be integers, the factorials mean gamma
functions, $a!=\Gamma(a+1)$, which is the form displayed above.}

\Needspace{5\baselineskip}
\par\bigskip
\noindent\textbf{20.9.3}\quad Using the convolution integral,
calculate
\[
L^{-1}\left\{\frac{s}{(s^{2}+a^{2})(s^{2}+b^{2})}\right\},
\qquad a^{2}\ne b^{2}.
\]
\solution
Factor the transform into
$[1/(s^2+a^2)]\,[s/(s^2+b^2)]$. For real $b$ and real $a\ne0$,
the factors are the transforms of $\sin(at)/a$ and $\cos bt$, so
the convolution theorem gives
\[
H(t)=\frac1a\int_0^t\sin(ax)\cos[b(t-x)]dx.
\]
Use $\sin u\cos v=[\sin(u+v)+\sin(u-v)]/2$ with $u=ax$,
$v=b(t-x)$, so that $u+v=(a-b)x+bt$ and $u-v=(a+b)x-bt$.
At $x=t$ both arguments equal $at$; at $x=0$ they are $\pm bt$,
and $\cos(-bt)=\cos bt$. Hence
\begin{align*}
H(t)&=\frac1{2a}\int_0^t
\{\sin[(a-b)x+bt]+\sin[(a+b)x-bt]\}dx\\
&=\frac1{2a}\left[-\frac{\cos[(a-b)x+bt]}{a-b}
-\frac{\cos[(a+b)x-bt]}{a+b}\right]_{x=0}^{x=t}\\
&=\frac1{2a}\left[
\frac{\cos bt-\cos at}{a-b}
+\frac{\cos bt-\cos at}{a+b}\right]\\
&=\frac{\cos bt-\cos at}{2a}\,
\frac{(a+b)+(a-b)}{a^2-b^2}\\
&=\boxed{\frac{\cos at-\cos bt}{b^2-a^2}}.
\end{align*}
The assumption $a^2\ne b^2$ makes $a\pm b$ nonzero.
\erratum{If $a=0$ (allowed when $b\ne0$), the first factor is
$1/s^2$, the transform of $x$ rather than $\sin(ax)/a$. Then
the substitution $w=t-x$ and an integration by parts give
$H(t)=\int_0^t(t-w)\cos bw\,dw=(1-\cos bt)/b^2$, which is the
boxed formula at $a=0$.}
This agrees with the independently transformed result of
Exercise~20.7.3.

\Needspace{5\baselineskip}
\par\bigskip
\noindent\textbf{20.9.4}\quad An undamped oscillator is driven by
a force $F_{0}\sin\omega t$. Find the displacement $X(t)$ as a
function of time, subject to initial conditions $X(0)=X'(0)=0$.
\erratum{For $\omega^2\ne\omega_0^2$}, the solution is a linear
combination of two simple harmonic motions, one with the frequency
of the driving force and one with the frequency $\omega_0>0$ of
the free oscillator. \erratum{Find the resonant limit separately.}
\par\nobreak\hfill\emph{ANS.}\enspace
$X(t)=\dfrac{F_{0}/m}{\omega^{2}-\omega_{0}^{2}}
\left(\dfrac{\omega}{\omega_{0}}\sin\omega_{0}t-\sin\omega t\right)$.
\par\erratum{The displayed quotient answer applies for $\omega^2\ne\omega_0^2$.}
\solution
Let $\omega_0^2=k/m>0$. The equation $mX''+kX=F_0\sin\omega t$
and the initial data, divided by $m$, are
\[
X''+\omega_0^2X=\frac{F_0}m\sin\omega t,\qquad X(0)=X'(0)=0.
\]
With $L\{X''\}=s^2\widehat X-sX(0)-X'(0)=s^2\widehat X$ and
$L\{\sin\omega t\}=\omega/(s^2+\omega^2)$, transformation gives
$(s^2+\omega_0^2)\widehat X=(F_0/m)\,\omega/(s^2+\omega^2)$, so
\[
\widehat X(s)=\frac{F_0}m\,
\frac\omega{(s^2+\omega^2)(s^2+\omega_0^2)}.
\]
Equivalently, since $1/(s^2+\omega_0^2)$ and
$\omega/(s^2+\omega^2)$ are the transforms of
$\sin(\omega_0t)/\omega_0$ and $\sin\omega t$, the convolution
theorem gives the response
\[
X(t)=\frac{F_0}{m\omega_0}\int_0^t
\sin[\omega_0(t-u)]\sin(\omega u)du.
\]
Partial fractions evaluate it more quickly. For
$\omega^2\ne\omega_0^2$,
$1/(s^2+\omega^2)-1/(s^2+\omega_0^2)
=(\omega_0^2-\omega^2)/[(s^2+\omega^2)(s^2+\omega_0^2)]$, so
\[
\widehat X=\frac{F_0\omega}{m(\omega_0^2-\omega^2)}
\left(\frac1{s^2+\omega^2}-\frac1{s^2+\omega_0^2}\right).
\]
Inverting the two terms by $L^{-1}\{1/(s^2+c^2)\}=\sin(ct)/c$
yields
\[
X(t)=\frac{F_0\omega}{m(\omega_0^2-\omega^2)}
\left(\frac{\sin\omega t}{\omega}
-\frac{\sin\omega_0t}{\omega_0}\right),
\]
that is, after moving $\omega$ into the bracket and changing
the sign of both factors,
\[
\boxed{X(t)=\frac{F_0/m}{\omega^2-\omega_0^2}
\left(\frac\omega{\omega_0}\sin\omega_0t-\sin\omega t\right)}.
\]
Its value vanishes at $t=0$, and so does its derivative, which
is proportional to $(\omega/\omega_0)\omega_0-\omega=0$ there:
the two sine slopes cancel.
\erratum{The printed answer requires a separate resonant
limit. At $\omega=\omega_0$ it is a removable $0/0$ form and
cannot be used as written.}
Apply l'H\^opital's rule to the ratio as $\omega\to\omega_0$;
for every $t$ the numerator
$(\omega/\omega_0)\sin\omega_0t-\sin\omega t$ and the
denominator $\omega^2-\omega_0^2$ both vanish at
$\omega=\omega_0$. Their $\omega$-derivatives are
$\sin(\omega_0t)/\omega_0-t\cos(\omega t)$ and $2\omega$, whose
ratio at $\omega=\omega_0$ is
$(\sin\omega_0t-\omega_0t\cos\omega_0t)/(2\omega_0^2)$. Thus
\[
\boxed{X_{\rm res}(t)=\frac{F_0}{2m\omega_0^2}
(\sin\omega_0t-\omega_0t\cos\omega_0t)}.
\]
This limit is indeed the solution at resonance. For
$\omega=\omega_0$ the transform is
$\widehat X=(F_0/m)\,\omega_0/(s^2+\omega_0^2)^2$, while the
rule $L\{tG(t)\}=-g'(s)$ for $g=L\{G\}$ gives
$L\{t\cos\omega_0t\}=-\frac d{ds}\frac s{s^2+\omega_0^2}
=\frac{s^2-\omega_0^2}{(s^2+\omega_0^2)^2}$ and hence
\[
L\{\sin\omega_0t-\omega_0t\cos\omega_0t\}
=\frac{\omega_0(s^2+\omega_0^2)-\omega_0(s^2-\omega_0^2)}
{(s^2+\omega_0^2)^2}
=\frac{2\omega_0^3}{(s^2+\omega_0^2)^2},
\]
so $X_{\rm res}=(F_0/m)\,\omega_0\cdot
(\sin\omega_0t-\omega_0t\cos\omega_0t)/(2\omega_0^3)$, as above.
The growing $t\cos\omega_0t$ term is the resonant response;
it is no longer a sum of two bounded oscillations at distinct
frequencies. For $\omega=0$ the driving force is identically
zero and the solution is $X=0$.

\Needspace{9\baselineskip}
\subsection*{20.10 Inverse Laplace Transform}

\Needspace{5\baselineskip}
\par\bigskip
\noindent\textbf{20.10.1}\quad Derive the Bromwich integral from
Cauchy's integral formula.
\par\emph{Hint.}~Apply the inverse transform $L^{-1}$ to
\[
f(s)=\frac{1}{2\pi i}\lim_{\alpha\to\infty}
\int_{\beta-i\alpha}^{\beta+i\alpha}\frac{f(z)}{s-z}\,dz,
\]
where $f(z)$ is analytic for $\Re e\,z\ge\beta$.
\par\erratum{Assume decay sufficient to make the right closing arc vanish and the usual Laplace inversion hypotheses; analyticity alone does not imply the identity in the hint. Take $\Re s>\beta$.}
\solution
Choose $\beta$ to the right of all singularities of the
Laplace transform, and fix $\Re s>\beta$. Close the upward
vertical segment $\beta-iR$ to $\beta+iR$ by a semicircle to
its right. This closed contour is clockwise, and it encloses
$s$ once $R>|s-\beta|$. The kernel
$f(z)/(s-z)=-f(z)/(z-s)$ has residue $-f(s)$ at $z=s$, so
Cauchy's formula with this orientation gives
\[
\int_{\beta-iR}^{\beta+iR}\frac{f(z)}{s-z}dz
+\int_{\text{right arc}}\frac{f(z)}{s-z}dz
=-2\pi i[-f(s)]=2\pi i f(s).
\]
If, for example, $|f(z)|\le M|z|^{-\epsilon}$ on these arcs for
some $\epsilon>0$, the arc integral is $O(R^{-\epsilon})$: the
arc $z=\beta+Re^{i\theta}$, $|\theta|\le\pi/2$, has length $\pi R$,
and on it $|s-z|\ge R-|s-\beta|\ge R/2$ once $R\ge2|s-\beta|$.
Letting $R\to\infty$ then gives the identity in the hint.
For each $z$ on the line $\Re z=\beta$,
\[
\frac1{s-z}=\int_0^\infty e^{-st}e^{zt}dt,
\qquad L_s^{-1}\{(s-z)^{-1}\}(t)=e^{zt},
\]
since $|e^{-(s-z)t}|=e^{-(\Re s-\beta)t}$ decays. Insert this
integral into the Cauchy representation and interchange the
$z$ and $t$ integrations (formally; see the remarks below):
\begin{align*}
f(s)&=\frac1{2\pi i}\lim_{R\to\infty}
\int_{\beta-iR}^{\beta+iR}f(z)
\left[\int_0^\infty e^{-st}e^{zt}dt\right]dz\\
&=\int_0^\infty e^{-st}\left[\frac1{2\pi i}\lim_{R\to\infty}
\int_{\beta-iR}^{\beta+iR}e^{zt}f(z)dz\right]dt.
\end{align*}
Thus $f$ is the Laplace transform of the last bracket, and
applying $L^{-1}$ (uniqueness of the inverse transform) gives
\[
\boxed{F(t)=\frac1{2\pi i}\lim_{R\to\infty}
\int_{\beta-iR}^{\beta+iR}e^{zt}f(z)dz},\qquad t>0,
\]
which is the Bromwich inversion formula at a continuity point
of $F$.
\erratum{Analyticity alone is not enough to discard the arc
or interchange a limiting inverse transform: $f(z)\equiv1$ is
analytic, yet for real $s>\beta$ (put $z=\beta+i\omega$,
$c=s-\beta$, so $1/(s-z)=(c+i\omega)/(c^2+\omega^2)$) the right
side of the hint is
$(2\pi)^{-1}\lim_{\alpha\to\infty}\int_{-\alpha}^{\alpha}
(c+i\omega)\,d\omega/(c^2+\omega^2)
=(2\pi)^{-1}\int_{-\infty}^\infty c\,d\omega/(c^2+\omega^2)
=1/2$, not $1$ (the odd imaginary part cancels between
$\pm\omega$). Decay on the
closing contours and inversion hypotheses are required. A
standard sufficient setting is a causal, piecewise smooth
function of exponential order, with the integral interpreted
by symmetric truncation (or Abel regularization). At a jump
it recovers the mean of the two one-sided limits.}
Writing $z=\beta+i\omega$ turns this into the
Fourier inversion of $e^{-\beta t}F(t)$ extended by zero to
$t<0$; Exercise~20.10.2 supplies a regularized kernel proof
of the limiting step.

\Needspace{5\baselineskip}
\par\bigskip
\noindent\textbf{20.10.2}\quad Starting with
\[
\frac{1}{2\pi i}\int_{\beta-i\infty}^{\beta+i\infty}
e^{st}f(s)\,ds,
\]
show that by introducing
\[
f(s)=\int_{0}^{\infty}e^{-sz}F(z)\,dz
\]
we can convert our integral into the Fourier representation of a
Dirac delta function. From this derive the inverse Laplace
transform.
\par\erratum{Use a Bromwich line to the right of the convergence abscissa, and recover $F(t)$ at positive continuity points; interpret the delta integral distributionally.}
\solution
Parameterize the line by $s=\beta+i\omega$, so
$ds/(2\pi i)=d\omega/(2\pi)$ and
$e^{st}e^{-sz}=e^{\beta(t-z)}e^{i\omega(t-z)}$. On inserting
the Laplace integral and interchanging the integrations, the
proposed inverse becomes formally
\begin{align*}
\frac1{2\pi i}\int_{\beta-i\infty}^{\beta+i\infty}e^{st}f(s)ds
&=\int_0^\infty F(z)e^{\beta(t-z)}
\left[\frac1{2\pi}\int_{-\infty}^{\infty}
e^{i\omega(t-z)}d\omega\right]dz\\
&=\int_0^\infty F(z)e^{\beta(t-z)}\delta(t-z)dz=F(t),
\qquad t>0,
\end{align*}
by the Fourier representation
$\delta(u)=\frac1{2\pi}\int_{-\infty}^\infty e^{i\omega u}d\omega$
and the sifting property; the factor $e^{\beta(t-z)}$ equals $1$
at $z=t$.

Here is a precise justification of the delta. Multiply the
frequency integrand by $e^{-\epsilon\omega^2}$, with
$\epsilon>0$. Completing the square,
$-\epsilon\omega^2+i\omega u
=-\epsilon[\omega-iu/(2\epsilon)]^2-u^2/(4\epsilon)$, and
$\int_{-\infty}^\infty e^{-\epsilon v^2}dv=\sqrt{\pi/\epsilon}$
(the shift of the line of integration is allowed by
Cauchy's theorem) give the Fourier integral
\[
K_\epsilon(u)=\frac1{2\pi}\int_{-\infty}^{\infty}
e^{-\epsilon\omega^2+i\omega u}d\omega
=\frac{e^{-u^2/(4\epsilon)}}{\sqrt{4\pi\epsilon}}.
\]
It is positive, even and decreasing in $|u|$, and has mass $1$
(substitute $u=2\sqrt\epsilon\,v$); for fixed $\delta>0$ both
$K_\epsilon(\delta)$ and $\int_{|u|\ge\delta}K_\epsilon(u)du$
tend to zero as $\epsilon\to0$. Let
$H(z)=e^{-\beta z}F(z)$ for $z\ge0$ and $H(z)=0$ for $z<0$,
with $H\in L^1$ ($\beta$ beyond the abscissa of absolute
convergence). Since $e^{-\epsilon\omega^2}|H(z)|$ is integrable
over the $(\omega,z)$ plane, Fubini's theorem justifies the
interchange in the regularized inverse:
\[
\frac1{2\pi}\int_{-\infty}^{\infty}e^{-\epsilon\omega^2}
e^{st}f(s)\,d\omega
=e^{\beta t}\int_{-\infty}^{\infty}K_\epsilon(t-z)H(z)dz.
\]
Let $t>0$ be a continuity point of $F$, hence of $H$. Given
$\eta>0$, choose $\delta>0$ with $|H(z)-H(t)|<\eta$ for
$|z-t|<\delta$. Because $K_\epsilon$ has mass $1$ and
$K_\epsilon(t-z)\le K_\epsilon(\delta)$ for $|z-t|\ge\delta$,
\[
\left|\int_{-\infty}^{\infty}K_\epsilon(t-z)H(z)dz-H(t)\right|
\le\eta+K_\epsilon(\delta)\int_{-\infty}^\infty|H|\,dz
+|H(t)|\int_{|u|\ge\delta}K_\epsilon(u)du.
\]
The last two terms tend to zero as $\epsilon\to0$, and $\eta$
is arbitrary. Hence the regularized inverse tends to
$e^{\beta t}H(t)=F(t)$ at every positive continuity
point. This proves
\[
\boxed{F(t)=\frac1{2\pi i}\int_{\beta-i\infty}^{\beta+i\infty}
e^{st}f(s)ds},\qquad t>0,
\]
with the integral understood as the limit $\epsilon\to0$ of
its regularized form.
\erratum{The unsmoothed frequency integral is a distribution,
not an absolutely convergent integral. At an interior jump
symmetric inversion gives $[F(t-)+F(t+)]/2$; at $t=0$ it gives
$F(0+)/2$ for the causal extension. For $t<0$ it gives zero.}

\Needspace{5\baselineskip}
\par\bigskip
\noindent\textbf{20.10.3}\quad Derive the Laplace transformation
convolution theorem by use of the Bromwich integral.
\solution
Let $f=L\{F\}$ and $g=L\{G\}$, with both functions
extended by zero for negative arguments. In a common
half-plane of convergence choose a Bromwich line $\Re s=\beta$.
Substitute the defining integral for $g$ into the inverse of
the product, combine $e^{st}e^{-su}=e^{s(t-u)}$, and
interchange the integrations:
\begin{align*}
L^{-1}\{fg\}(t)
&=\frac1{2\pi i}\int_{\beta-i\infty}^{\beta+i\infty}
e^{st}f(s)\left[\int_0^\infty e^{-su}G(u)du\right]ds\\
&=\int_0^\infty G(u)\left[
\frac1{2\pi i}\int_{\beta-i\infty}^{\beta+i\infty}
e^{s(t-u)}f(s)ds\right]du.
\end{align*}
The bracket is the Bromwich integral of Exercise~20.10.2
evaluated at the time $t-u$. It therefore equals $F(t-u)$ for
$u<t$ and vanishes for $u>t$, where $t-u<0$ (the inverse is
causal). Hence for $t\ge0$ the $u$ integral stops at
$u=t$, and the result is
\[
\boxed{L^{-1}\{f(s)g(s)\}(t)=\int_0^tF(t-u)G(u)du}.
\]
Changing variables $x=t-u$ ($du=-dx$; the limits $u=0,t$
become $x=t,0$) also gives
$\int_0^tF(x)G(t-x)dx$, the Laplace convolution.
For a direct rigorous interchange, it suffices that the
vertical integral of $|f(\beta+i\omega)|$ and
$\int_0^\infty e^{-\beta u}|G(u)|du$ are finite: on the line
$|e^{s(t-u)}f(s)G(u)|=e^{\beta t}|f(\beta+i\omega)|\,
e^{-\beta u}|G(u)|$, so the double integral converges
absolutely and Fubini's theorem applies.
For the usual broader class of locally integrable functions
of exponential order, insert the Gaussian frequency regulator
of Exercise~20.10.2, interchange the absolutely convergent
integrals, and then take its approximate-identity limit at
continuity points of the convolution. Values at isolated
endpoints do not alter the convolution integral.

\Needspace{5\baselineskip}
\par\bigskip
\noindent\textbf{20.10.4}\quad Find
\[
L^{-1}\left\{\frac{s}{s^{2}-k^{2}}\right\}
\]
\begin{enumerate}
\item[(a)] by a partial fraction expansion.
\item[(b)] Repeat, using the Bromwich integral.
\end{enumerate}
\solution
(a) For $k\ne0$, since $(s+k)+(s-k)=2s$,
\[
\frac{s}{s^2-k^2}
=\frac12\left(\frac1{s-k}+\frac1{s+k}\right).
\]
Thus, by $L\{e^{\pm kt}\}=1/(s\mp k)$,
\[
\boxed{L^{-1}\left\{\frac{s}{s^2-k^2}\right\}
=\frac{e^{kt}+e^{-kt}}2=\cosh kt}.
\]
(b) Take $\beta>|k|$ for real $k$. For $t>0$, close the
Bromwich segment from $\beta-iT$ to $\beta+iT$ by the arc of
the circle $|s|=R=\sqrt{\beta^2+T^2}$ to its left; the closed
contour is positively oriented. The following arc estimate
is used throughout this section. Suppose the factor $\varphi(s)$
multiplying $e^{st}$ satisfies $|\varphi(s)|\le M_R$ on the
arc, with $M_R\to0$ as $R\to\infty$. On the left semicircle
$s=Re^{i\theta}$, $\pi/2\le\theta\le3\pi/2$, put
$\theta=\pi/2+\phi$; then $|e^{st}|=e^{-tR\sin\phi}$, and
Jordan's inequality $\sin\phi\ge2\phi/\pi$ on $[0,\pi/2]$ gives
\[
\int|e^{st}|\,|ds|=2R\int_0^{\pi/2}e^{-tR\sin\phi}d\phi
\le2R\int_0^{\pi/2}e^{-2tR\phi/\pi}d\phi<\frac\pi t.
\]
The two remaining pieces (present for $\beta>0$), where
$0<\Re s<\beta$, have length
$R\arcsin(\beta/R)\le\pi\beta/2$ each, and there
$|e^{st}|\le e^{\beta t}$. Hence the arc contributes at most
$M_R(\pi/t+\pi\beta e^{\beta t})\to0$. Here
$\varphi(s)=s/(s^2-k^2)$ and $|s^2-k^2|\ge R^2-k^2$, so
$M_R=R/(R^2-k^2)$. The Bromwich integral therefore equals
the sum of the residues at the simple poles $s=\pm k$, each
found by cancelling its factor $s\mp k$ and setting $s=\pm k$:
\[
\mathop{\rm Res}_{s=k}\frac{se^{st}}{(s-k)(s+k)}
=\frac{ke^{kt}}{2k}=\frac{e^{kt}}2,
\]
\[
\mathop{\rm Res}_{s=-k}\frac{se^{st}}{(s-k)(s+k)}
=\frac{-ke^{-kt}}{-2k}=\frac{e^{-kt}}2.
\]
Their sum $(e^{kt}+e^{-kt})/2$ reproduces $\cosh kt$. At $k=0$
the transform reduces directly to $1/s$, with inverse $1$, also
the continuous limit of the answer. At $t=0$ the formula gives
the right-hand value $F(0+)=\cosh0=1$.

\Needspace{5\baselineskip}
\par\bigskip
\noindent\textbf{20.10.5}\quad Find
\[
L^{-1}\left\{\frac{k^{2}}{s(s^{2}+k^{2})}\right\}
\]
\begin{enumerate}
\item[(a)] by using a partial fraction expansion.
\item[(b)] Repeat using the convolution theorem.
\item[(c)] Repeat using the Bromwich integral.
\end{enumerate}
\par\nobreak\hfill\emph{ANS.}\enspace$F(t)=1-\cos kt$.
\solution
Assume real $k\ne0$ during the algebra.
(a) Writing the numerator as $k^2=(s^2+k^2)-s^2$ gives the
decomposition
\[
\frac{k^2}{s(s^2+k^2)}=\frac1s-\frac{s}{s^2+k^2},
\]
and the pairs $L\{1\}=1/s$, $L\{\cos kt\}=s/(s^2+k^2)$ give
$F(t)=1-\cos kt$.

(b) Factor the transform as
$[1/s]\,[k^2/(s^2+k^2)]$. The inverse factors are $1$ and
$k\sin kt$, so convolution gives
\[
F(t)=\int_0^t 1\cdot k\sin(ku)du=[-\cos(ku)]_0^t=1-\cos kt.
\]
(c) For a Bromwich line to the right of $0$, close to the
left for $t>0$. The rational factor is $O(s^{-3})$: on
$|s|=R>|k|$ it is at most $M_R=k^2/[R(R^2-k^2)]$, so the closing
arc vanishes by the arc estimate of Exercise~20.10.4(b). At
the three simple poles, cancelling each pole factor and
evaluating the rest there, with $(ik)(2ik)=(-ik)(-2ik)=-2k^2$,
\begin{align*}
\mathop{\rm Res}_{s=0}\frac{k^2e^{st}}{s(s-ik)(s+ik)}
&=\frac{k^2}{(-ik)(ik)}=1,\\
\mathop{\rm Res}_{s=ik}\frac{k^2e^{st}}{s(s-ik)(s+ik)}
&=\frac{k^2e^{ikt}}{(ik)(2ik)}=-\frac12e^{ikt},\\
\mathop{\rm Res}_{s=-ik}\frac{k^2e^{st}}{s(s-ik)(s+ik)}
&=\frac{k^2e^{-ikt}}{(-ik)(-2ik)}=-\frac12e^{-ikt}.
\end{align*}
Summing gives
\[
\boxed{F(t)=1-\frac12(e^{ikt}+e^{-ikt})=1-\cos kt}.
\]
At $k=0$ the transform and its inverse are both zero, which
agrees with the continuous limit. All three methods give
$F(0)=0$.

\Needspace{5\baselineskip}
\par\bigskip
\noindent\textbf{20.10.6}\quad Use the Bromwich integral to find
the function whose transform is $f(s)=s^{-1/2}$. Note that $f(s)$
has a branch point at $s=0$. The negative $x$-axis may be taken
as a cut line. See Fig.~20.23.
\par\emph{Hint.}~A portion of the path needed to close the
contour will yield nonzero contributions to the contour integral.
These will need to be taken into account to get the proper value
for the Bromwich integral.
\par\nobreak\hfill\emph{ANS.}\enspace$F(t)=(\pi t)^{-1/2}$.
\solution
Choose the principal branch
$s^{-1/2}=|s|^{-1/2}e^{-i\arg(s)/2}$ with
$-\pi<\arg(s)<\pi$. For $t>0$, close the Bromwich segment
to the left as in Exercise~20.10.4(b), but leave out the cut:
the closed contour follows the arc down to the upper bank of
the cut, runs along the upper bank from $-R$ to $-\rho$,
circles the origin clockwise on $|s|=\rho$, returns along the
lower bank from $-\rho$ to $-R$, and finishes on the arc. It
encloses no singularity, so its integral vanishes by Cauchy's
theorem. On the arc $|f|=R^{-1/2}\to0$, so the arc estimate of
Exercise~20.10.4(b) applies; on the small circle (length
$2\pi\rho$) $|e^{st}f|\le e^{\rho t}\rho^{-1/2}$, so it
contributes $O(\sqrt\rho)$ as $\rho\to0$. Hence the Bromwich
integral equals the reversed bank integrals: along the lower
bank from $-\infty$ to $0$ and along the upper bank from $0$
to $-\infty$, a counterclockwise path around the cut.
There are no poles; the two banks supply the entire inverse.
For $r>0$, $\arg s\to\pm\pi$ on the two banks, so
\[
f(-r+i0)=r^{-1/2}e^{-i\pi/2}=-i r^{-1/2},\qquad
f(-r-i0)=r^{-1/2}e^{i\pi/2}=+i r^{-1/2}.
\]
With $s=-r$, $ds=-dr$, the upper bank ($r$ from $0$ to
$\infty$) contributes $-\int_0^\infty e^{-rt}f(-r+i0)dr$ and
the lower bank ($r$ from $\infty$ to $0$) contributes
$+\int_0^\infty e^{-rt}f(-r-i0)dr$. Their sum is therefore
\[
\int_0^\infty e^{-rt}
\big[f(-r-i0)-f(-r+i0)\big]dr
=2i\int_0^\infty e^{-rt}r^{-1/2}dr.
\]
Dividing by $2\pi i$ and substituting $u=rt$ (so that
$r^{-1/2}dr=t^{-1/2}u^{-1/2}du$) gives, with
$\Gamma(1/2)=\sqrt\pi$,
\[
F(t)=\frac1\pi\int_0^\infty e^{-rt}r^{-1/2}dr
=\frac{\Gamma(1/2)}{\pi\sqrt t}
=\boxed{\frac1{\sqrt{\pi t}}},\qquad t>0.
\]
Its integrable singularity at zero is compatible with a
Laplace transform: direct substitution gives
$\int_0^\infty e^{-st}(\pi t)^{-1/2}dt=s^{-1/2}$ for
real $s>0$, fixing the branch and normalization independently.

\Needspace{5\baselineskip}
\par\bigskip
\noindent\textbf{20.10.7}\quad Show that
\[
L^{-1}\big\{(s^{2}+1)^{-1/2}\big\}=J_{0}(t)
\]
by evaluation of the Bromwich integral.
\par\emph{Hint.}~Convert your Bromwich integral into an integral
representation of $J_{0}(t)$. Figure~20.24 shows a possible
contour.
\solution
Take a cut along the segment from $-i$ to $i$ and choose
$D(s)=\sqrt{s^2+1}$ analytic outside this segment with
$D(s)\sim s$ at infinity and $D(s)>0$ for positive real $s$.
For $s=x+iy$ with $x>0$, the number $s^2+1=x^2-y^2+1+2ixy$
never lies on $(-\infty,0]$ (a real value needs $xy=0$, so
$y=0$, and then $s^2+1=x^2+1>0$). Hence the principal root
$\sqrt{s^2+1}$ is analytic on $\Re s>0$; its ratio to $D$ is a
continuous function with values $\pm1$ on this connected
half-plane, so it is constant, and it equals $1$ for $s>0$.
Thus $D$ is the principal root there. Letting $x\to0+$ gives,
on the right bank
$s=iy+0$, $-1<y<1$,
\[
D(iy+0)=+\sqrt{1-y^2}.
\]
On the left bank its sign is reversed:
$-D(-s)$ is also analytic outside the cut, squares to
$s^2+1$ and behaves as $s$ at infinity, so $D(-s)=-D(s)$ and
$D(iy-0)=-D(-iy+0)=-\sqrt{1-y^2}$.
For $t>0$, close the Bromwich segment by the arc of
Exercise~20.10.4(b). On $|s|=R>1$,
$|1/D(s)|\le(R^2-1)^{-1/2}\to0$, so the arc contributes
nothing in the limit. The closed contour encloses the cut,
and by Cauchy's theorem it may be shrunk onto a
counterclockwise contour hugging the cut: upward along the
right bank, around a small circle about $i$, downward along
the left bank, and around a small circle about $-i$. The
endpoint singularities are inverse square roots
($|D(s)|^2=|s-i|\,|s+i|$), so the circles, of radius $\rho$
and length $2\pi\rho$, contribute $O(\sqrt\rho)$. With
$s=iy$, $ds=i\,dy$, and
$1/D_{\rm right}-1/D_{\rm left}=2/\sqrt{1-y^2}$,
\begin{align*}
F(t)&=\frac1{2\pi i}
\int_{-i}^{i}e^{st}\left[\frac1{D(s)_{\rm right}}
-\frac1{D(s)_{\rm left}}\right]ds
=\frac1{2\pi i}\int_{-1}^1e^{iyt}
\frac{2\,i\,dy}{\sqrt{1-y^2}}\\
&=\frac1\pi\int_{-1}^1\frac{e^{iyt}}{\sqrt{1-y^2}}dy
=\frac2\pi\int_0^1\frac{\cos(yt)}{\sqrt{1-y^2}}dy.
\end{align*}
The sine part of $e^{iyt}$ cancels by oddness. Finally
$y=\cos\theta$ converts the symmetric integral
$\frac1\pi\int_{-1}^1\cos(yt)(1-y^2)^{-1/2}dy$: here
$dy=-\sin\theta\,d\theta$, $\sqrt{1-y^2}=\sin\theta\ge0$ on
$[0,\pi]$, and $y=-1,1$ correspond to $\theta=\pi,0$, so the
reversed limits absorb the minus sign. By the integral
representation of Exercise~14.1.17(b) with $\nu=0$,
\[
F(t)=\frac1\pi\int_0^\pi\cos(t\cos\theta)d\theta
=\boxed{J_0(t)}.
\]
This also has the right-hand initial value $J_0(0)=1$, in
agreement with the initial-value relation
$F(0+)=\lim_{s\to\infty}sf(s)=\lim_{s\to\infty}s/\sqrt{s^2+1}=1$.
The branch must
be continued through the exterior of the chosen cut; taking
an unrelated square-root sign on the left would spoil the
jump and the factor of two.

\Needspace{5\baselineskip}
\par\bigskip
\noindent\textbf{20.10.8}\quad Evaluate the inverse Laplace
transform
\[
L^{-1}\big\{(s^{2}-a^{2})^{-1/2}\big\}
\]
by each of the following methods:
\begin{enumerate}
\item[(a)] Expansion in a series and term-by-term inversion.
\item[(b)] Direct evaluation of the Bromwich integral.
\item[(c)] Change of variable in the Bromwich integral:
$s=(a/2)(z+z^{-1})$.
\end{enumerate}
\solution
Take $a>0$ and the branch of
$\sqrt{s^2-a^2}$ asymptotic to $s$, with cut $[-a,a]$:
$D(s)=\sqrt{s-a}\,\sqrt{s+a}$ with principal square roots.
The first factor changes sign across $(-\infty,a)$ and the
second across $(-\infty,-a)$; on $(-\infty,-a)$ the two sign
changes cancel, so $D$ is analytic outside $[-a,a]$, and
$D(s)=s\sqrt{1-a^2/s^2}\sim s$. Write $f=1/D$.
The Bromwich line lies at $\Re s=\beta>a$.

(a) For $|s|>a$ put $w=a^2/s^2$, so $|w|<1$. The binomial
series $(1-w)^{-1/2}=\sum_n\binom{-1/2}{n}(-w)^n$ has the
coefficients
\[
(-1)^n\binom{-1/2}{n}=\frac{1\cdot3\cdots(2n-1)}{2^nn!}
=\frac{(2n)!}{4^n(n!)^2},
\]
since $1\cdot3\cdots(2n-1)=(2n)!/(2^nn!)$. Hence on this branch
\[
(s^2-a^2)^{-1/2}=\frac1s(1-w)^{-1/2}
=\frac1s\sum_{n=0}^\infty
\frac{(2n)!}{4^n(n!)^2}\frac{a^{2n}}{s^{2n}}.
\]
The inverse of $s^{-2n-1}$ is $t^{2n}/(2n)!$, which cancels
the $(2n)!$, and $a^{2n}t^{2n}/4^n=(at/2)^{2n}$. Thus, by the
power series of $I_\nu$ at $\nu=0$,
\[
F(t)=\sum_{n=0}^\infty
\frac{(at/2)^{2n}}{(n!)^2}=\boxed{I_0(at)}.
\]
The resulting series converges for every finite $t$.
To justify the termwise inversion, transform it back. For
real $s>a$ every term $e^{-st}(at/2)^{2n}/(n!)^2$ is
nonnegative, so by monotone convergence the series may be
integrated term by term:
\[
L\{I_0(at)\}(s)=\sum_{n=0}^\infty\frac{a^{2n}}{4^n(n!)^2}\,
\frac{(2n)!}{s^{2n+1}}=(s^2-a^2)^{-1/2},
\]
the convergent binomial expansion again. By uniqueness of
the inverse transform, $F=I_0(at)$; no uniform convergence on
a Bromwich line is needed.

(b) For $-a<x<a$, $\sqrt{s+a}\to\sqrt{a+x}>0$ on both banks,
while $\arg(s-a)\to\pm\pi$ gives $\sqrt{s-a}\to\pm i\sqrt{a-x}$
on the upper and lower bank. Hence
$D(x\pm i0)=\pm i\sqrt{a^2-x^2}$, and the bank values of $f$
are
\[
f(x+i0)=-\frac i{\sqrt{a^2-x^2}},\qquad
f(x-i0)=+\frac i{\sqrt{a^2-x^2}}.
\]
For $t>0$, close the Bromwich segment by the arc of
Exercise~20.10.4(b). On $|s|=R>a$,
$|f|\le(R^2-a^2)^{-1/2}\to0$, so the arc drops out, and by
Cauchy's theorem the closed contour may be shrunk onto a
counterclockwise contour around the cut. Its upper bank runs
from $a$ to $-a$, and its lower bank from $-a$ to $a$. The
endpoint circles vanish as the square root of their radius,
since $|D(s)|^2=|s-a|\,|s+a|$ (as in Exercise~20.10.7).
Reversing the limits of the upper-bank integral gives
\begin{align*}
F(t)&=\frac1{2\pi i}\int_{-a}^a e^{xt}
\big[f(x-i0)-f(x+i0)\big]dx
=\frac1{2\pi i}\int_{-a}^a\frac{2i\,e^{xt}}{\sqrt{a^2-x^2}}dx\\
&=\frac1\pi\int_{-a}^a\frac{e^{xt}}{\sqrt{a^2-x^2}}dx
=\frac1\pi\int_0^\pi e^{at\cos\theta}d\theta
=I_0(at).
\end{align*}
In the substitution $x=a\cos\theta$,
$dx=-a\sin\theta\,d\theta$ and $\sqrt{a^2-x^2}=a\sin\theta$
for $0\le\theta\le\pi$; the limits $x=-a,a$ become
$\theta=\pi,0$, and the reversed limits cancel the sign in
$dx$. The last equality is the integral representation of
Exercise~14.5.14 with $n=0$.

(c) Set $s=(a/2)(z+z^{-1})$, taking the exterior inverse
$z=(s+\sqrt{s^2-a^2})/a$, for which $z\sim2s/a$ at infinity.
Since $s^2-a^2=\frac{a^2}4(z^2+2+z^{-2}-4)
=\frac{a^2}4(z-z^{-1})^2$ and $ds=\frac a2(1-z^{-2})dz$,
\[
\sqrt{s^2-a^2}=\frac a2(z-z^{-1}),\qquad
\frac{ds}{\sqrt{s^2-a^2}}
=\frac{(1-z^{-2})dz}{z-z^{-1}}=\frac{dz}{z};
\]
the sign of the root is fixed by $\sqrt{s^2-a^2}\sim s\sim az/2$
at infinity, and $z-z^{-1}=z(1-z^{-2})$. Deform the $s$
contour around its cut as in (b). For $x=a\cos\theta$,
$0<\theta<\pi$, the bank values $D(x\pm i0)$ of (b) give
$z=(x\pm i\sqrt{a^2-x^2})/a=e^{\pm i\theta}$. So the upper
bank ($x$ from $a$ to $-a$, $\theta$ from $0$ to $\pi$) maps to
the upper unit semicircle from $1$ to $-1$,
and the lower bank to the lower semicircle from $-1$ to $1$.
Their union is one counterclockwise unit circle, giving
\[
F(t)=\frac1{2\pi i}\oint_{|z|=1}
\exp\left[\frac{at}{2}(z+z^{-1})\right]\frac{dz}{z},
\]
which is Exercise~14.5.3(a) with $n=0$ and $x=at$. Directly,
expand the two exponentials on this fixed circle:
\[
\frac1z\sum_{n,m\ge0}
\frac{(at/2)^{n+m}}{n!m!}z^{n-m}.
\]
On $|z|=1$ the double series is dominated by
$\sum_{n,m}(at/2)^{n+m}/(n!m!)=e^{at}$, so it may be integrated
term by term, and $\oint z^{n-m-1}dz$ is $2\pi i$ for $n=m$ and
$0$ otherwise. Only $n=m$ contributes to the coefficient of
$z^{-1}$, so the residue is
$\sum_{n\ge0}(at/2)^{2n}/(n!)^2=I_0(at)$.
\erratum{For $at\ne0$, the singularity at $z=0$ is essential,
not the simple pole called for in the Instructor's Manual; the
residue calculation still applies. Also the image of the
undeformed vertical Bromwich line is a curve, not exactly a
vertical line in $z$. Mapping the already deformed cut contour
avoids that unwarranted replacement.}
The answer is even in $a$, so real negative $a$ gives the
same result with $|a|$. At $a=0$ the transform is $1/s$ and
its inverse is $1$.

\Needspace{5\baselineskip}
\par\bigskip
\noindent\textbf{20.10.9}\quad Show that
\[
L^{-1}\left\{\frac{\ln s}{s}\right\}=-\ln t-\gamma,
\]
where $\gamma=0.5772\ldots$ is the Euler--Mascheroni constant.
\solution
For real $q>0$ and $s>0$, the gamma substitution $u=st$ gives
\[
L\left\{\frac{t^{q-1}}{\Gamma(q)}\right\}(s)
=\frac1{\Gamma(q)}\int_0^\infty e^{-st}t^{q-1}dt
=\frac{s^{-q}}{\Gamma(q)}\int_0^\infty e^{-u}u^{q-1}du=s^{-q}.
\]
Differentiate with respect to $q$ near $q=1$. For
$|q-1|\le\frac12$ the $q$-derivative of the integrand,
$e^{-st}t^{q-1}\ln t$, is bounded by
$e^{-st}(t^{-1/2}+t^{1/2})|\ln t|$, which is integrable on
$(0,\infty)$ (exponential decay at infinity, an integrable
singularity at zero), so differentiation inside the transform
is justified. With
$\partial_qt^{q-1}=t^{q-1}\ln t$,
$\partial_q[1/\Gamma(q)]=-\Gamma'(q)/\Gamma(q)^2$ and
$\partial_qs^{-q}=-(\ln s)s^{-q}$,
\[
L\left\{\frac{t^{q-1}}{\Gamma(q)}
[\ln t-\psi(q)]\right\}(s)=-(\ln s)s^{-q},
\qquad\psi(q)=\frac{\Gamma'(q)}{\Gamma(q)}.
\]
Since $\Gamma(1)=1$ and $\psi(1)=-\gamma$ (Section~13.2),
setting $q=1$ gives $L\{\ln t+\gamma\}=-(\ln s)/s$, and
multiplying by $-1$ yields
\[
\boxed{L^{-1}\left\{\frac{\ln s}{s}\right\}(t)
=-\ln t-\gamma},\qquad t>0.
\]
For complex $s$ in $\Re s>0$, take the branch of $\ln s$
that is real on $s>0$; both sides are analytic in $\Re s>0$,
so the identity extends there. The logarithmic singularity of the
answer at $t=0$ is integrable, so it causes no failure of the
ordinary Laplace transform.

\Needspace{5\baselineskip}
\par\bigskip
\noindent\textbf{20.10.10}\quad Evaluate the Bromwich integral
for
\[
f(s)=\frac{s}{(s^{2}+a^{2})^{2}}.
\]
\solution
For real $a\ne0$, close the Bromwich line to the left
when $t>0$. The integrand
\[
\frac{se^{st}}{(s-ia)^2(s+ia)^2}
\]
has double poles at $s=\pm ia$ and a rational factor
$O(s^{-3})$: on $|s|=R>|a|$ it is at most $R/(R^2-a^2)^2$.
The closing arc therefore vanishes by the arc estimate of
Exercise~20.10.4(b). At a double pole $s_0$, the residue of
$\varphi(s)/(s-s_0)^2$, with $\varphi$ analytic at $s_0$, is
$\varphi'(s_0)$. At $s=ia$ take $\varphi(s)=se^{st}/(s+ia)^2$;
the product and quotient rules give
\[
\varphi'(s)=e^{st}\left[\frac{1+st}{(s+ia)^2}
-\frac{2s}{(s+ia)^3}\right].
\]
With $(2ia)^2=-4a^2$ and $(2ia)^3=-8ia^3$,
\begin{align*}
R_+=\varphi'(ia)
&=e^{iat}\left[\frac{1+iat}{(2ia)^2}
-\frac{2ia}{(2ia)^3}\right]\\
&=e^{iat}\left[-\frac{1+iat}{4a^2}+\frac1{4a^2}\right]
=-\frac{it}{4a}e^{iat}.
\end{align*}
Similarly, at $s=-ia$ the factor is $se^{st}/(s-ia)^2$, with
$(s-ia)^2=-4a^2$ and $(s-ia)^3=8ia^3$ there:
\[
R_-=e^{-iat}\left[\frac{1-iat}{-4a^2}-\frac{-2ia}{8ia^3}\right]
=e^{-iat}\left[-\frac{1-iat}{4a^2}+\frac1{4a^2}\right]
=\frac{it}{4a}e^{-iat}.
\]
Summing, with $e^{-iat}-e^{iat}=-2i\sin at$,
\[
F(t)=R_++R_-=
\frac{it}{4a}(e^{-iat}-e^{iat})
=\boxed{\frac t{2a}\sin at}.
\]
As a check, the rule $L\{tG(t)\}=-g'(s)$ for $g=L\{G\}$, with
$G=\sin at$ and $g=a/(s^2+a^2)$, gives
$L\{t\sin at\}=2as/(s^2+a^2)^2$, which is $2a$ times the
given transform.
\erratum{At $a=0$ the two double poles merge and the reduced
transform is $s^{-3}$; its inverse is $t^2/2$, the continuous
limit of the formula above.}

\Needspace{5\baselineskip}
\par\bigskip
\noindent\textbf{20.10.11}\quad Heaviside expansion theorem. If
the transform $f(s)$ may be written as a ratio
\[
f(s)=\frac{g(s)}{h(s)},
\]
where $g(s)$ and $h(s)$ are analytic functions, with $h(s)$
having simple, isolated zeros at $s=s_i$,
\erratum{and with a vanishing left closing-contour contribution
and a convergent residue sum}, show that
\[
F(t)=L^{-1}\left\{\frac{g(s)}{h(s)}\right\}
=\sum_{i}\frac{g(s_{i})}{h'(s_{i})}\,e^{s_{i}t}.
\]
\par\emph{Hint.}~See Exercise~11.6.3.
\par\erratum{A sufficient elementary setting is a strictly proper rational transform with finitely many simple poles, for $t>0$. The unqualified analytic-function assertion in the textbook is false, as the counterexample below shows.}
\solution
At a simple zero $s_i$ of $h$,
\[
h(s)=h'(s_i)(s-s_i)+O((s-s_i)^2),\qquad h'(s_i)\ne0.
\]
The residue of the inverse-transform integrand is therefore,
since $(s-s_i)/h(s)=(s-s_i)/[h(s)-h(s_i)]\to1/h'(s_i)$
(this is Exercise~11.6.3 with $f_1=e^{st}g$, $f_2=h$),
\[
\mathop{\rm Res}_{s=s_i}\frac{e^{st}g(s)}{h(s)}
=\lim_{s\to s_i}\frac{(s-s_i)e^{st}g(s)}{h(s)}
=\frac{g(s_i)}{h'(s_i)}e^{s_it}.
\]
Place the Bromwich line to the right of all the $s_i$, and
close it to the left by contours on which the contribution
vanishes in the limit. If the residue sums inside these
contours converge, the residue theorem applied to each closed
contour gives in the limit
\[
\boxed{F(t)=\sum_i\frac{g(s_i)}{h'(s_i)}e^{s_it}},\qquad t>0.
\]
A particularly simple sufficient case is a strictly proper
rational function $g/h$ with finitely many simple poles: then
$|g/h|\le M/|s|$ on large circles, so the arc estimate of
Exercise~20.10.4(b) applies, and one large circle encloses all
the poles. A zero of $h$ at which $g$ also vanishes is a
removable singularity of $g/h$; its residue is zero, and so is
its term $g(s_i)/h'(s_i)$.
\erratum{Analytic $g,h$ and simple zeros of $h$ alone do not
imply the stated expansion: one must also control the closing
contour, and an infinite residue sum needs its own convergence
prescription.}
For a concrete counterexample satisfying the printed analytic
conditions, choose $g(s)=e^{-s}$, $h(s)=s+1$. Its only pole
is $-1$, where $g(-1)/h'(-1)=e$, so the proposed unrestricted
sum would be $e\cdot e^{-t}=e^{1-t}$. But direct
transformation, with $v=t-1$, gives
\[
\int_1^\infty e^{-st}e^{-(t-1)}dt
=e^{-s}\int_0^\infty e^{-(s+1)v}dv=\frac{e^{-s}}{s+1},
\]
whose causal inverse is $e^{-(t-1)}u(t-1)$: it equals $e^{1-t}$
for $t>1$ but is zero for $0<t<1$.
For these times the integrand $e^{st}g/h=e^{s(t-1)}/(s+1)$ has
$|e^{s(t-1)}|=e^{(t-1)\Re s}$, which grows as $\Re s\to-\infty$
on a left closing arc, explaining exactly why the omitted
hypothesis fails. A polynomial part of a rational transform likewise
requires impulses at zero rather than this ordinary
exponential sum.

\Needspace{5\baselineskip}
\par\bigskip
\noindent\textbf{20.10.12}\quad Using the Bromwich integral,
invert $f(s)=s^{-2}e^{-ks}$, \erratum{$k\ge0$}. Express $F(t)=L^{-1}\{f(s)\}$ in
terms of the (shifted) unit step function $u(t-k)$.
\par\nobreak\hfill\emph{ANS.}\enspace$F(t)=(t-k)u(t-k)$.
\solution
\erratum{The restriction $k\ge0$ is needed: for $k<0$ the ramp
$(t-k)u(t-k)$ does not have the transform $s^{-2}e^{-ks}$, as
the check at the end shows.}
Since $e^{st}e^{-ks}=e^{s(t-k)}$, the inversion integral is
\[
F(t)=\frac1{2\pi i}\int_{\beta-i\infty}^{\beta+i\infty}
\frac{e^{s(t-k)}}{s^2}ds,\qquad\beta>0.
\]
Close the segment from $\beta-iT$ to $\beta+iT$ by an arc of
the circle $|s|=R=\sqrt{\beta^2+T^2}$. If $t<k$, close to the
right ($\Re s\ge\beta$), where
$|e^{s(t-k)}|=e^{(t-k)\Re s}\le e^{(t-k)\beta}$. No poles are
enclosed, so $F(t)=0$. If $t>k$, close to the left
($\Re s\le\beta$), where again $|e^{s(t-k)}|\le e^{(t-k)\beta}$,
and enclose the double pole at $s=0$. Its residue is the
derivative of $e^{s(t-k)}$ at $s=0$ (as for the double poles
of Exercise~20.10.10):
\[
\left.\frac d{ds}e^{s(t-k)}\right|_{s=0}=t-k.
\]
Both arcs have length at most $2\pi R$, and $|s^{-2}|=R^{-2}$
on them, so they contribute at most
$2\pi e^{(t-k)\beta}/R\to0$. At $t=k$ the value is zero,
either by continuity or because $1/s^2$ has zero residue at
$s=0$ (the same arc bound applies).
Thus
\[
\boxed{F(t)=(t-k)u(t-k)}.
\]
A forward check with $v=t-k$ gives
\[
\int_k^\infty e^{-st}(t-k)dt
=e^{-ks}\int_0^\infty e^{-sv}v\,dv=\frac{e^{-ks}}{s^2}.
\]
If $k<0$, the ramp equals $t-k$ for all $t\ge0$, and its
one-sided transform is instead
$\int_0^\infty e^{-st}(t-k)dt=s^{-2}-k/s\ne s^{-2}e^{-ks}$; this
confirms the need for the nonnegative delay in the one-sided
formula.

\Needspace{5\baselineskip}
\par\bigskip
\noindent\textbf{20.10.13}\quad You have a Laplace transform:
\[
f(s)=\frac{1}{(s+a)(s+b)},\qquad a\ne b.
\]
Invert this transform by each of three methods:
\begin{enumerate}
\item[(a)] Partial fractions and use of tables,
\item[(b)] Convolution theorem,
\item[(c)] Bromwich integral.
\end{enumerate}
\par\nobreak\hfill\emph{ANS.}\enspace
$F(t)=\dfrac{e^{-bt}-e^{-at}}{a-b}$,\quad$a\ne b$.
\solution
(a) Since $(s+a)-(s+b)=a-b$, write
\[
\frac1{(s+a)(s+b)}
=\frac1{a-b}\left(\frac1{s+b}-\frac1{s+a}\right).
\]
The exponential transform pair $L\{e^{-ct}\}=1/(s+c)$ gives
\[
F(t)=\frac{e^{-bt}-e^{-at}}{a-b}.
\]
(b) The inverse factors are $e^{-at}$ and $e^{-bt}$.
Their convolution, with
$e^{-au}e^{-b(t-u)}=e^{-bt}e^{-(a-b)u}$ and
$e^{-bt}e^{-(a-b)t}=e^{-at}$, is
\begin{align*}
F(t)&=\int_0^t e^{-au}e^{-b(t-u)}du
=e^{-bt}\int_0^t e^{-(a-b)u}du\\
&=e^{-bt}\frac{1-e^{-(a-b)t}}{a-b}
=\frac{e^{-bt}-e^{-at}}{a-b}.
\end{align*}
(c) Choose $\beta>\max(-\Re a,-\Re b)$ and close the
Bromwich line to the left for $t>0$. The rational factor
is $O(s^{-2})$, so the arc vanishes by the arc estimate of
Exercise~20.10.4(b). The residues at the simple poles, found
by cancelling each pole factor and setting $s$ to the pole,
are
\[
\mathop{\rm Res}_{s=-a}\frac{e^{st}}{(s+a)(s+b)}
=\frac{e^{-at}}{b-a},\qquad
\mathop{\rm Res}_{s=-b}\frac{e^{st}}{(s+a)(s+b)}
=\frac{e^{-bt}}{a-b}.
\]
Their sum gives the same result:
\[
\boxed{F(t)=\frac{e^{-bt}-e^{-at}}{a-b}}.
\]
Its initial value is zero and its initial slope is
$F'(0)=(-b+a)/(a-b)=1$.
The excluded coincident case has the regular limit
\[
\frac{e^{-bt}-e^{-at}}{a-b}
=e^{-at}\,\frac{e^{(a-b)t}-1}{a-b}\to te^{-at}
\qquad(b\to a),
\]
since $(e^{\epsilon t}-1)/\epsilon\to t$ as $\epsilon=a-b\to0$;
it is also obtained from the double pole of $e^{st}/(s+a)^2$, whose
residue is $\frac d{ds}e^{st}\big|_{s=-a}=te^{-at}$.
